\section{Local covers, global quotients and cusp carriers}\label{sec:III:carriers}

\subsection{The Kummer cover extends across the toric cusp}
\label{wb:cusp-cover-section}

This section determines exactly what is true in the assertion that the
degree-$N^2$ cover of the regular fibres extends by a toric power map.  The
answer is positive for the completed \emph{local} cusp family, but the source
action has to be changed, and two different degrees must not be conflated.

Retain the cusp data
\[
 B_0=\begin{pmatrix}0&1\\-1&0\end{pmatrix},\qquad
 Z(s)=sB_0+C(t),\qquad t=e^{2\pi i s},
\]
where $C$ is holomorphic on a disc $\Delta_t$ about zero.  Put
\[
 c_{\bar\lambda}(t)
 :=\exp\!\bigl(2\pi iC(t)\bar\lambda\bigr)
 \in(\mathbb C^*)^2,
 \qquad \bar\lambda\in\bar\Lambda\simeq\mathbb Z^2.
\]
Let $\Sigma$ be the fan in
$N'=\mathbb Z^2\oplus\mathbb Z$ consisting of the cones over the cells of
the translation-invariant $A_2$ triangulation at height one, and write
$Y=Y_\Sigma$.  On the dense torus, with coordinates $(x_1,x_2,t)$, the
target action is
\begin{equation}
 \Psi_{\bar\lambda}(x,t)
 =\bigl(c_{\bar\lambda}(t)t^{B_0\bar\lambda}x,t\bigr).
 \label{wb:cusp-cover-target-action}
\end{equation}
Its quotient is the completed cusp family
$f_0\colon N_0=Y/\bar\Lambda\to\Delta_t$.

Fix an integer $N\geq1$.  Let $v$ be a coordinate on a second disc and
make the base change
\begin{equation}
 \rho_N\colon\Delta_v\longrightarrow\Delta_t,\qquad
 \rho_N(v)=v^N,\qquad t=v^N.
 \label{wb:cusp-cover-base-change}
\end{equation}
On a copy $Y^{(N)}$ of the same toric variety, with dense-torus coordinates
$(\xi_1,\xi_2,v)$, define
\begin{equation}
 C_N(v):=\frac1N C(v^N),\qquad
 c^{(N)}_{\bar\lambda}(v)
 :=\exp\!\bigl(2\pi iC_N(v)\bar\lambda\bigr)
 \label{wb:cusp-cover-source-factor}
\end{equation}
and
\begin{equation}
 \Psi^{(N)}_{\bar\lambda}(\xi,v)
 :=\bigl(c^{(N)}_{\bar\lambda}(v)
          v^{B_0\bar\lambda}\xi,v\bigr).
 \label{wb:cusp-cover-source-action}
\end{equation}

\begin{theorem}[Exact Kummer extension through the cusp]
\label{wb:cusp-cover-extension-theorem}
\StatementRecord{wb:cusp-cover-extension-theorem}{proved}{Historical events: GPT-6 Astra (OpenAI), ChatGPT Pro mode, effort Pro, S6\_09\_PAPERS, 17 November 2026: targeted proof check. Also checked in this writing. Pre-continuation writing events: GPT-6 Astra (OpenAI), ChatGPT Pro mode, effort Pro, S6\_10\_PAPERS, 18 November 2026: targeted proof check. Continuation metadata events: GPT-6 Astra (OpenAI), ChatGPT Pro mode, effort Pro, S6\_10\_PAPERS continuation, 19 November 2026: used and rechecked in this run; status-vocabulary reconciliation; metadata reconciliation only; no new mathematical check. \emph{Pre-continuation recorded proof scope.} GPT-6 Astra (OpenAI, ChatGPT Pro mode), S6\_09\_PAPERS, effort Pro, 2026-11-17; also checked in this writing by GPT-6 Astra (OpenAI), ChatGPT Pro mode, effort Pro, 18 November 2026}{the hypotheses, notation and definitions stated here; \cref{wb:cusp-cover-source-action}; \cref{wb:cusp-cover-target-action}; \cref{wb:cusp-cover-source-factor}; \cref{wb:cusp-cover-maximal-chart-map}; \cref{wb:cusp-cover-basechange-local-ring}; \cref{wb:cusp-cover-triple-fibre-length}; \cref{wb:cusp-cover-double-fibre-length}; \cref{wb:cusp-cover-additive-fibre-map}; \cref{wb:cusp-cover-equivariance-computation}}

\mbox{}\ 
After shrinking the $v$-disc, the following statements hold.

\begin{enumerate}
\item The maps in \eqref{wb:cusp-cover-source-action} extend to a free,
properly discontinuous action on $Y^{(N)}$.  The quotient
\[
 f_N\colon N_0^{\langle N\rangle}
 :=Y^{(N)}/\bar\Lambda\longrightarrow\Delta_v
\]
is a proper holomorphic map from a smooth threefold.  Its central fibre
$W_N$ is reduced, irreducible, and nonnormal, and its normalization is a
toric del Pezzo surface of degree six.

\item Multiplication by $N$ on the cocharacter lattice defines a finite
flat toric morphism
\begin{equation}
 P_N\colon Y^{(N)}\longrightarrow Y,\qquad
 P_N(\xi_1,\xi_2,v)=(\xi_1^N,\xi_2^N,v^N),
 \label{wb:cusp-cover-toric-power-map}
\end{equation}
of degree $N^3$.  It is equivariant for
\eqref{wb:cusp-cover-source-action} and
\eqref{wb:cusp-cover-target-action}, hence descends to a finite flat map
\begin{equation}
 \kappa_N\colon N_0^{\langle N\rangle}\longrightarrow N_0,
 \qquad f_0\circ\kappa_N=\rho_N\circ f_N,
 \label{wb:cusp-cover-completed-map}
\end{equation}
of degree $N^3$.

\item The map $\kappa_N$ is etale over $v\ne0$, and
\begin{equation}
 \kappa_N^*W=NW_N,
 \qquad
 \operatorname{Ram}(\kappa_N)=(N-1)W_N.
 \label{wb:cusp-cover-direct-ramification}
\end{equation}
For $N>1$, the branch support of the map to the original cusp family
is the whole central fibre. For $N=1$, the displayed formulas give the identity and the branch support is empty.

\item Let
\begin{equation}
 N_0[N]:=N_0\times_{\Delta_t}\Delta_v,
 \qquad t=v^N.
 \label{wb:cusp-cover-base-changed-target}
\end{equation}
Then $N_0[N]$ is normal and $\kappa_N$ factors uniquely as
\begin{equation}
 N_0^{\langle N\rangle}
 \xrightarrow{\ Q_N\ }N_0[N]
 \longrightarrow N_0.
 \label{wb:cusp-cover-relative-factorization}
\end{equation}
The map $Q_N$ has generic degree $N^2$.  For $N>1$ it is quasi-etale but
not flat: it is etale over the smooth locus of $N_0[N]$, and its non-etale
and nonflat loci are exactly the inverse images of the double and triple
strata of the base-changed central fibre.  For $N>1$, its degree $N^2>1$ proves that $Q_N$ is not a
normalization map. For $N=1$ it is the identity of the normal target.

\item On the reduced central fibres, the induced map
$W_N\to W$ has generic degree $N^2$.  It lifts to a commutative diagram
\begin{equation}
\begin{CD}
 E_N\simeq\mathrm{dP}_6 @>{[N]}>> E\simeq\mathrm{dP}_6\\
 @V{\nu_N}VV @VV{\nu}V\\
 W_N @>>> W,
\end{CD}
 \label{wb:cusp-cover-normalization-diagram}
\end{equation}
where $[N]$ is toric multiplication by $N$, of degree $N^2$, ramified
with index $N$ along each of the six toric boundary curves.  The diagram
respects the paired conductor curves and the two triple-point classes.

\item Over $v\ne0$, $Q_N$ is exactly the cover defined by
\[
 \mathbb Z\langle\widehat\gamma,\widehat u,
 N\widehat w,N\widehat\delta\rangle
 \subset
 \mathbb Z\langle\widehat\gamma,\widehat u,
 \widehat w,\widehat\delta\rangle.
\]
Its deck group is $(\mathbb Z/N\mathbb Z)^2$.  The pulled-back cusp
monodromy is $M_0^N$, and in the displayed cover basis it is $M_0$:
\begin{equation}
 S_N^{-1}M_0^NS_N=M_0,\qquad
 S_N=\operatorname{diag}(1,1,N,N).
 \label{wb:cusp-cover-monodromy-conjugacy}
\end{equation}
\end{enumerate}
\end{theorem}



\begin{proof}
\begin{enumerate}[label=\arabic*.,leftmargin=*]
\item We prove the assertions in the order in which the constructions occur.

\item \medskip\noindent
\emph{The source action.}
For $\bar\lambda\in\bar\Lambda$, the integral linear map
\[
 \varphi_{\bar\lambda}(y,y_3)
 =(y+y_3B_0\bar\lambda,y_3)
\]
translates the height-one slice by the integral vector
$B_0\bar\lambda$.  It preserves the $A_2$ triangulation and hence induces
a toric automorphism $\Phi_{\bar\lambda}$ of $Y^{(N)}$.  On the dense
torus this automorphism is
\[
 \Phi_{\bar\lambda}(\xi,v)
 =(v^{B_0\bar\lambda}\xi,v).
\]
Multiplication by the nowhere-zero holomorphic function
$c^{(N)}_{\bar\lambda}(v)$ extends across every toric chart.  Since
$C_N(v)\bar\lambda$ is additive in $\bar\lambda$, the resulting maps
satisfy
$\Psi^{(N)}_{\bar\lambda}\Psi^{(N)}_{\bar\mu}
=\Psi^{(N)}_{\bar\lambda+\bar\mu}$.

\item For completeness, freeness and proper discontinuity can be checked
directly.  Put
\[
 R_N(v):=-2\pi\operatorname{Im}C_N(v).
\]
It is bounded at $v=0$.  At a dense-torus point fixed by
$\bar\lambda$, logarithmic absolute values give
\[
 \log|v|B_0\bar\lambda+R_N(v)\bar\lambda=0.
\]
Choose the disc so small that
$|\log|v||>2\|R_N\|_\infty$.  Because $B_0$ is orthogonal, the last
equality is impossible for $\bar\lambda\ne0$.  On the central fibre,
$\Psi^{(N)}_{\bar\lambda}$ translates every bounded cell by the nonzero
vector $B_0\bar\lambda$, so it fixes no toric orbit.

\item For a dense-torus point set
\[
 y(\xi,v)=\frac{(\log|\xi_1|,\log|\xi_2|)}{\log|v|},
 \qquad
 B_{N,v}=B_0+\frac{R_N(v)}{\log|v|}.
\]
Then
\[
 y(\Psi^{(N)}_{\bar\lambda}(\xi,v))
 =y(\xi,v)+B_{N,v}\bar\lambda,
 \qquad
 |B_{N,v}\bar\lambda|\ge\tfrac12|\bar\lambda|.
\]
% checked carrier correction: compact-chart-control
The logarithmic positions of an entire affine toric chart need not be
bounded. The required assertion is about its compact subsets, and follows
as follows. In a maximal unimodular chart write its height-one rays as
$(b_0,1),(b_1,1),(b_2,1)$ and its coordinates as $w_0,w_1,w_2$.
On its dense torus,
\[
 v=w_0w_1w_2,\qquad
 y=\frac{\log|\xi|}{\log|v|}
   =\sum_{i=0}^2\lambda_i b_i,\qquad
 \lambda_i=\frac{\log|w_i|}{\log|v|},\qquad
 \sum_i\lambda_i=1.
\]
For a compact subset of this chart all $|w_i|$ have a common upper
bound $M\ge1$. When $|v|<1$, this gives
$\lambda_i\ge-\log M/|\log|v||$. Thus its logarithmic positions remain
in a bounded neighbourhood of the fixed triangle as $v\to0$.
On a closed annulus away from zero, the product identity and upper
bounds also bound each $|w_i|$ below. Hence the same boundedness holds
there. Every point of a compact set has a neighbourhood whose closure
is compact in one maximal chart.  Choose a finite such cover;
the intersections with these compact closures give finitely many
compact chart pieces covering the original set.  Applying the
preceding estimate to these pieces bounds all of its logarithmic
positions on the dense-torus locus.

\item For the modified action,
$y\mapsto y+B_{N,v}\bar\lambda$ and
$|B_{N,v}\bar\lambda|\ge|\bar\lambda|/2$.
Consequently only finitely many translates of one compact set can
meet another on the dense torus. At $v=0$, choose finitely many
maximal charts covering the two compact sets. A translate meeting
both must send a face of one of their triangles to a face of another.
There are finitely many pairs of faces and at most one translation
for each pair; thus the same finiteness holds at the central fibre.
This proves proper discontinuity without asserting boundedness of
whole affine charts.

\item To prove that $f_N$ is proper, reduce
$B_{N,v}^{-1}y$ modulo $\mathbb Z^2$ for $v\ne0$. The resulting $y$
lies in one bounded set, uniformly on the chosen small disc. Choose
an $A_2$ triangle containing $y$. Its barycentric coordinates are
nonnegative, and the corresponding toric coordinates satisfy
$|w_i|\le1$. Only finitely many triangles meet that bounded set.
Fix a closed subdisc strictly inside the chosen action disc.
Every orbit over that subdisc has a representative in finitely
many closed unit polydiscs intersected with its inverse image.
These intersections are closed subsets of compact polydiscs
in the ambient smooth Hausdorff toric space, and all lie in the
domain where the modified action is defined. At $v=0$, vertex translations move a component containing the
point to a fixed component; that component is the compact toric
surface with hexagonal star. Adding this compact component gives
a compact set mapping onto the quotient above a closed subdisc.
The inverse image of any compact subset of the base is a closed
subset of this compact image, and is therefore compact.

\item The central-fibre and
normalization assertions follow exactly from the height-one fan: locally
$v$ is the product of the coordinates belonging to the rays through a
cell, every component occurs with multiplicity one, the lattice action is
transitive on vertices, and the star of a vertex is the complete smooth
hexagonal fan.

\item \medskip\noindent
\emph{The toric map and equivariance.}
Multiplication by $N$ on $N'$ sends every cone $\sigma\in\Sigma$ onto
itself as a real cone.  It therefore defines the toric morphism $P_N$.
For a maximal cone, let $z_0,z_1,z_2$ be the target characters dual to its
three primitive height-one ray generators, and let $w_0,w_1,w_2$ be the
corresponding source characters.  The map on convergent local rings is
\begin{equation}
 \mathbb C\{z_0,z_1,z_2\}\longrightarrow
\mathbb C\{w_0,w_1,w_2\},\qquad z_i\longmapsto w_i^N.
 \label{wb:cusp-cover-maximal-chart-map}
\end{equation}
The target-chart preimage is the corresponding source chart, and the
right-hand ring is free of rank $N^3$ over the left-hand ring, with basis
\[
 \{w_0^{a_0}w_1^{a_1}w_2^{a_2}:0\le a_i<N\}.
\]
Thus $P_N$ is finite flat of degree $N^3$ on every maximal chart, hence
globally.

\item The choice \eqref{wb:cusp-cover-source-factor} satisfies equivariance by the following exact identity.
Indeed, writing powers coordinatewise,
\begin{align}
 P_N\Psi^{(N)}_{\bar\lambda}(\xi,v)
 &=\bigl((c^{(N)}_{\bar\lambda}(v))^N
          v^{NB_0\bar\lambda}\xi^N,v^N\bigr)\\
 &=\bigl(c_{\bar\lambda}(v^N)
          (v^N)^{B_0\bar\lambda}\xi^N,v^N\bigr)\\
 &=\Psi_{\bar\lambda}P_N(\xi,v).
 \label{wb:cusp-cover-equivariance-computation}
\end{align}
Consequently $P_N$ descends to \eqref{wb:cusp-cover-completed-map}.  To see
finiteness and flatness after the infinite-lattice quotient without any
implicit descent theorem, take an evenly covered open set in $N_0$ and
one of its lifts $V\subset Y$.  Equivariance gives a unique representative
of every source orbit above that open set whose image lies in $V$; two
such representatives cannot differ by a nontrivial lattice translate.
The descended map is therefore locally identified with
$P_N^{-1}(V)\to V$.  It is finite flat of degree $N^3$ by
\eqref{wb:cusp-cover-maximal-chart-map}.

\item In the same chart the base characters are
\[
 t=z_0z_1z_2,\qquad v=w_0w_1w_2,\qquad P_N^*t=v^N.
\]
The Jacobian determinant of \eqref{wb:cusp-cover-maximal-chart-map} is
\[
 N^3(w_0w_1w_2)^{N-1}=N^3v^{N-1}.
\]
This proves \eqref{wb:cusp-cover-direct-ramification}; the local quotient
maps are biholomorphisms, so the calculation is unchanged downstairs.

\item \medskip\noindent
\emph{The relative degree-$N^2$ factor.}
The equality
\((f_0\circ\kappa_N)([\xi,v])=v^N=(\rho_N\circ f_N)([\xi,v])\)
gives the factorization
\eqref{wb:cusp-cover-relative-factorization}.  At a triple point the exact
local-ring map is
\begin{equation}
 \begin{gathered}
 A_N:=
 \frac{\mathbb C\{z_0,z_1,z_2,v\}}
      {(z_0z_1z_2-v^N)},
 \qquad
 B:=\mathbb C\{w_0,w_1,w_2\},\\
 A_N\longrightarrow B,\qquad
 z_i\longmapsto w_i^N,\qquad
 v\longmapsto w_0w_1w_2.
 \end{gathered}
 \label{wb:cusp-cover-basechange-local-ring}
\end{equation}
Let
\[
 G_N=\{(\varepsilon_0,\varepsilon_1,\varepsilon_2)\in\mu_N^3:
       \varepsilon_0\varepsilon_1\varepsilon_2=1\}
 \simeq\mu_N^2
\]
act diagonally on $B$.  A monomial
$w_0^{a_0}w_1^{a_1}w_2^{a_2}$ is $G_N$-invariant if and only if
\[
 a_0\equiv a_1\equiv a_2\pmod N.
\]
Writing their common residue as $r$, with $0\le r<N$, expresses that
monomial uniquely in the form
\[
 (w_0w_1w_2)^r
 (w_0^N)^{b_0}(w_1^N)^{b_1}(w_2^N)^{b_2}.
\]
The same coefficientwise argument applies to convergent power series.
Hence
\begin{equation}
 A_N=B^{G_N}.
 \label{wb:cusp-cover-invariant-ring}
\end{equation}
It follows that $A_N$ is normal, that the map has generic degree
$|G_N|=N^2$, and that it is a finite quotient rather than a normalization.
The corresponding face-chart calculations prove normality along the
double and smooth strata as well, so $N_0[N]$ is normal everywhere.

\item The stabilizer in $G_N$ is trivial when none or exactly one of the $w_i$
vanishes, is isomorphic to $\mu_N$ when exactly two vanish, and is all of
$G_N$ when all three vanish.  Thus $Q_N$ is etale in codimension one and
has no ramification divisor, but it is not etale above the double and
triple strata.  It is also nonflat at every such point when $N>1$.  At a
triple point the closed fibre of \eqref{wb:cusp-cover-basechange-local-ring}
has length
\begin{align}
 \dim_{\mathbb C}
 \frac{\mathbb C\{w_0,w_1,w_2\}}
      {(w_0^N,w_1^N,w_2^N,w_0w_1w_2)}
 &=N^3-(N-1)^3\\
 &=3N^2-3N+1,
 \label{wb:cusp-cover-triple-fibre-length}
\end{align}
which exceeds the generic rank $N^2$.  At a point of a double stratum,
localizing at the nonzero third coordinate and choosing its holomorphic
$N$th root reduces the ramified factor to
\[
 \frac{\mathbb C\{z_0,z_1,v\}}{(z_0z_1-v^N)}
\longrightarrow\mathbb C\{w_0,w_1\},\qquad
 z_i\longmapsto w_i^N,\quad v\longmapsto w_0w_1.
\]
Its local generic rank is $N$, whereas its closed fibre has length
\begin{equation}
 \dim_{\mathbb C}
 \frac{\mathbb C\{w_0,w_1\}}
      {(w_0^N,w_1^N,w_0w_1)}=2N-1>N.
 \label{wb:cusp-cover-double-fibre-length}
\end{equation}
The remaining transverse nonzero torus coordinate is etale locally.
Equations \eqref{wb:cusp-cover-triple-fibre-length} and
\eqref{wb:cusp-cover-double-fibre-length}, together with etaleness elsewhere,
identify the nonflat locus exactly.

\item \medskip\noindent
\emph{The regular fibres and monodromy.}
Choose a logarithm $r$ with $v=e^{2\pi ir}$.  The target period matrix at
$t=v^N$ is
\[
 Z(Nr)=NrB_0+C(v^N),
\]
whereas the source residual action is determined by
\[
 Z_N(r)=rB_0+\frac1NC(v^N)=\frac1N Z(Nr).
\]
The source fibre is therefore
\[
 \frac{\mathbb C^2}{\mathbb Z^2+Z_N(r)\mathbb Z^2},
\]
and multiplication by $N$ on $\mathbb C^2$ maps it to the target fibre.
Equivalently, in the unscaled additive coordinate it is the identity map
\begin{equation}
 \frac{\mathbb C^2}
      {N\mathbb Z^2+Z(Nr)\mathbb Z^2}
 \longrightarrow
 \frac{\mathbb C^2}
      {\mathbb Z^2+Z(Nr)\mathbb Z^2}.
 \label{wb:cusp-cover-additive-fibre-map}
\end{equation}
The imaginary part of $Z(Nr)$ is invertible on a regular fibre, so its
period lattice meets $\mathbb Z^2$ trivially in the required real sense;
the index in \eqref{wb:cusp-cover-additive-fibre-map} is $N^2$.  Its deck
group is $(\mathbb Z/N\mathbb Z)^2$.  In the ordered lattice basis
$(\widehat\gamma,\widehat u,\widehat w,\widehat\delta)$, this is exactly
the sublattice obtained by $S_N$.  Since $M_0=I+R$ with $R^2=0$ and
$R(a,b,c,d)=(0,0,b,-a)$, one has
\[
 M_0^N=I+NR,\qquad S_N^{-1}M_0^NS_N=M_0,
\]
proving \eqref{wb:cusp-cover-monodromy-conjugacy}.

\item \medskip\noindent
\emph{The central normalization and conductor.}
Multiplication by $N$ fixes every ray of $\Sigma$.  It therefore maps a
chosen source component of the toric central fibre to the corresponding
target component.  The star of that ray is the hexagonal fan, so the
restriction is the toric power map
$[N]\colon\mathrm{dP}_6\to\mathrm{dP}_6$.  On its dense two-torus it has
degree $N^2$; in each toric boundary chart it is an $N$th-power map in the
normal coordinate, hence has ramification index $N$ along all six boundary
curves.  Equation \eqref{wb:cusp-cover-equivariance-computation} says that
the source and target lattice identifications commute with this map.
Consequently opposite boundary curves, their conductor identifications,
and the two translation-orbits of triangle vertices descend compatibly.
This proves \eqref{wb:cusp-cover-normalization-diagram} and completes the
proof.
\end{enumerate}
\end{proof}


\paragraph{Exact comparison with the unmodified source action}
\label{wb:cusp-cover-naive-action-failure}
\StatementRecord{wb:cusp-cover-naive-action-failure}{proved}{Historical events: GPT-6 Astra (OpenAI), ChatGPT Pro mode, effort Pro, S6\_09\_PAPERS, 17 November 2026: targeted proof check. Also checked in this writing. Pre-continuation writing events: GPT-6 Astra (OpenAI), ChatGPT Pro mode, effort Pro, S6\_10\_PAPERS, 18 November 2026: targeted proof check. Continuation metadata events: GPT-6 Astra (OpenAI), ChatGPT Pro mode, effort Pro, S6\_10\_PAPERS continuation, 19 November 2026: used and rechecked in this run; status-vocabulary reconciliation; metadata reconciliation only; no new mathematical check. \emph{Pre-continuation recorded proof scope.} GPT-6 Astra (OpenAI, ChatGPT Pro mode), S6\_09\_PAPERS, effort Pro, 2026-11-17; also checked in this writing by GPT-6 Astra (OpenAI), ChatGPT Pro mode, effort Pro, 18 November 2026}{the hypotheses, notation and definitions stated here; \cref{wb:cusp-cover-source-factor}}

If one reuses the target action with $v$ in place of $t$ at the source,
then the two sides of equivariance have multiplicative factors
\[
 c_{\bar\lambda}(v)^N v^{NB_0\bar\lambda}
 \quad\hbox{and}\quad
 c_{\bar\lambda}(v^N)v^{NB_0\bar\lambda},
\]
respectively. Equivariance of this unmodified action is exactly the equality $c_{\bar\lambda}(v)^N=c_{\bar\lambda}(v^N)$ for every $\bar\lambda$. The displayed source factor \eqref{wb:cusp-cover-source-factor} satisfies the required equality by taking its specified holomorphic logarithm. For $N=1$ the modified and unmodified actions agree. The explicit
logarithm $2\pi iC(v^N)\bar\lambda/N$ makes this root holomorphic; no
unjustified choice of an $N$th root is involved.






\paragraph{Exact scope}
\label{wb:cusp-cover-global-scope}
\StatementRecord{wb:cusp-cover-global-scope}{proved}{Historical events: GPT-6 Astra (OpenAI), ChatGPT Pro mode, effort Pro, S6\_09\_PAPERS, 17 November 2026: targeted proof check. Also checked in this writing. Pre-continuation writing events: GPT-6 Astra (OpenAI), ChatGPT Pro mode, effort Pro, S6\_10\_PAPERS, 18 November 2026: targeted proof check. Continuation metadata events: GPT-6 Astra (OpenAI), ChatGPT Pro mode, effort Pro, S6\_10\_PAPERS continuation, 19 November 2026: used and rechecked in this run; status-vocabulary reconciliation; metadata reconciliation only; no new mathematical check. \emph{Pre-continuation recorded proof scope.} GPT-6 Astra (OpenAI, ChatGPT Pro mode), S6\_09\_PAPERS, effort Pro, 2026-11-17; also checked in this writing by GPT-6 Astra (OpenAI), ChatGPT Pro mode, effort Pro, 18 November 2026}{the hypotheses, notation and definitions stated here; the argument and qualifications immediately following}

The theorem constructs a finite morphism of completed families over the
local cusp discs.  It is a degree-$N^3$ map to the original family and a
degree-$N^2$ map only after the target is base changed.  For $N>1$, its base map $v\mapsto v^N$ differs from the identity on the disc, and the degree $N^2>1$ of the relative factor excludes a normalization map. At $N=1$ both factors are identities.

Nothing here supplies a cover over the entire base $B$.  In fact the
displayed sublattice is not preserved by the original full monodromy when
$N>1$: for example, the printed matrix $A_1$ sends
$\widehat u$ to a vector whose $\widehat w$- and
$\widehat\delta$-coefficients are $-1$ and $1$.  A further global base
cover and compatible extensions through both finite fillings would have
to be constructed separately.  The targets of the displayed Kummer morphisms are exactly $N_0$ and $N_0[N]$. The additional maps to the sphere carriers are constructed explicitly in the subsequent carrier results.





\subsection{Radial period coordinates and the compact quotient}
\label{mapsIII:carrier-radial-prerequisites}

The carrier constructions use the following real quotient of the period
family of \cref{wb:prop:period-input-audit}.  Fix the lattice basis
\[
 \Lambda=\mathbb Z\langle\widehat\gamma,\widehat u,
                         \widehat w,\widehat\delta\rangle.
\]
The cusp coordinate and its holomorphic residual functions give
\[
 \Pi=[Z\mid I_2],\qquad
 Z=s\begin{pmatrix}0&1\\-1&0\end{pmatrix}
       +\begin{pmatrix}6\mu(t)&h(t)\\b(t)-h(t)&\mu(t)\end{pmatrix},
 \qquad t=e^{2\pi i s}.
\]
Here $\mu,b,h$ are holomorphic on a disc containing zero.  On the positive
radial ray put $t=e^{-2\pi T}$ and $s=iT$.  Choose $T_0$ large enough that
this ray lies in that disc for $T\geq T_0$, and set
\begin{equation}
 \begin{aligned}
 \mu_T&=\alpha_T+im_T,& h_T&=\eta_T+iH_T,& b_T&=\xi_T+iB_T,\\
 c_T&=\xi_T-\eta_T,& L_T&=T+H_T,& q_T&=L_T-B_T,\\
 D_T&=L_Tq_T+6m_T^2.
 \end{aligned}
 \label{wb:cusp-spec-radial-notation}
\end{equation}
All these quantities are real.  Write
$m=\operatorname{Im}\mu(0)$, $H=\operatorname{Im}h(0)$ and
$B=\operatorname{Im}b(0)$.  Taylor expansion at zero yields
\begin{equation}
 m_T=m+O(e^{-2\pi T}),\qquad
 H_T=H+O(e^{-2\pi T}),\qquad
 B_T=B+O(e^{-2\pi T}),
 \label{wb:cusp-spec-holomorphic-errors}
\end{equation}
and the analogous estimates for the real parts.  In particular
$D_T=T^2+O(T)>0$ after increasing $T_0$.

Identify $\mathbb C^2$ with $\mathbb R^4$ in the order
$(x_1,y_1,x_2,y_2)$ and equip it with the Euclidean metric.  The induced
flat metric on a quotient is the metric used below.  This metric is an
explicit choice in these coordinates.

\begin{proposition}[Real periods and their dual coordinates]
\label{mapsIII:carrier-real-periods}
\StatementRecord{mapsIII:carrier-real-periods}{proved}{Historical events: GPT-6 Astra (OpenAI), ChatGPT Pro mode, effort Pro, S6\_09\_PAPERS, 17 November 2026: targeted proof check. Also checked in this writing. Pre-continuation writing events: GPT-6 Astra (OpenAI), ChatGPT Pro mode, effort Pro, S6\_10\_PAPERS, 18 November 2026: targeted proof check. Continuation metadata events: GPT-6 Astra (OpenAI), ChatGPT Pro mode, effort Pro, S6\_10\_PAPERS continuation, 19 November 2026: used and rechecked in this run; status-vocabulary reconciliation; metadata reconciliation only; no new mathematical check. \emph{Pre-continuation recorded proof scope.} GPT-6 Astra (OpenAI, ChatGPT Pro mode), effort Pro, 17 November 2026; also checked in this writing by GPT-6 Astra (OpenAI), ChatGPT Pro mode, effort Pro, 18 November 2026}{the hypotheses, notation and definitions stated here; the explicit constructions and calculations in the proof}

The radial fibre and its real period matrix are
\[
 F_T=\mathbb R^4/P_T\mathbb Z^4,\qquad
 P_T=\begin{pmatrix}
 6\alpha_T&\eta_T&1&0\\
 6m_T&L_T&0&0\\
 c_T&\alpha_T&0&1\\
 -q_T&m_T&0&0
 \end{pmatrix}.
\]
Its determinant is $-D_T$.  Its complete inverse transpose is
\begin{equation}
 P_T^{-\mathsf T}=
 \begin{pmatrix}
 0&0&1&0\\[2pt]
 \dfrac{m_T}{D_T}&\dfrac{q_T}{D_T}&
 \dfrac{-6m_T\alpha_T-q_T\eta_T}{D_T}&
 \dfrac{-m_Tc_T-q_T\alpha_T}{D_T}\\[7pt]
 0&0&0&1\\[2pt]
 -\dfrac{L_T}{D_T}&\dfrac{6m_T}{D_T}&
 \dfrac{6(L_T\alpha_T-m_T\eta_T)}{D_T}&
 \dfrac{L_Tc_T-6m_T\alpha_T}{D_T}
 \end{pmatrix}.
 \label{wb:cusp-spec-full-inverse-transpose}
\end{equation}
Thus $(P_T\mathbb Z^4)^*=P_T^{-\mathsf T}\mathbb Z^4$.
\end{proposition}



\begin{proof}
\begin{enumerate}[label=\arabic*.,leftmargin=*]
\item At $s=iT$ the complex period matrix is
\[
 Z_T=\begin{pmatrix}6\mu_T&h_T+iT\\
                          b_T-h_T-iT&\mu_T\end{pmatrix}.
\]
Taking real and imaginary parts of its columns and of the identity
columns gives the displayed $P_T$.  Expansion along its last two
columns gives $\det P_T=-L_Tq_T-6m_T^2=-D_T$.

\item To solve $P_T^{\mathsf T}x=(r,s,n_1,n_2)^{\mathsf T}$, its last two
equations first give $x_1=n_1$ and $x_3=n_2$.  Put
\[
 R=r-6\alpha_Tn_1-c_Tn_2,\qquad
 S=s-\eta_Tn_1-\alpha_Tn_2.
\]
The remaining two equations and their solution are
\[
 \begin{pmatrix}6m_T&-q_T\\L_T&m_T\end{pmatrix}
 \binom{x_2}{x_4}=\binom R S,
 \qquad
 \binom{x_2}{x_4}=
 \frac1{D_T}\begin{pmatrix}m_T&q_T\\-L_T&6m_T\end{pmatrix}
 \binom R S.
\]
Substituting $R,S$ proves every entry in
\eqref{wb:cusp-spec-full-inverse-transpose}.  Finally a real covector
$\ell$ is integral on $P_T\mathbb Z^4$ exactly when
$P_T^{\mathsf T}\ell\in\mathbb Z^4$, proving the dual-lattice formula.
\end{enumerate}
\end{proof}

Let $e_1,\ldots,e_4$ denote the coordinate unit vectors in this real
order and let $v_{\gamma,T},v_{u,T}$ be the first two columns of $P_T$.
Then
\[
 P_T\mathbb Z^4=
 \mathbb Zv_{\gamma,T}\oplus\mathbb Zv_{u,T}
                 \oplus\mathbb Ze_1\oplus\mathbb Ze_3.
\]

\begin{proposition}[Exact compact quotient of the native cusp fibre]
\label{wb:cusp-spec-native-quotient}
\StatementRecord{wb:cusp-spec-native-quotient}{proved}{Historical events: GPT-6 Astra (OpenAI), ChatGPT Pro mode, effort Pro, S6\_09\_PAPERS, 17 November 2026: targeted proof check. Also checked in this writing. Pre-continuation writing events: GPT-6 Astra (OpenAI), ChatGPT Pro mode, effort Pro, S6\_10\_PAPERS, 18 November 2026: targeted proof check. Continuation metadata events: GPT-6 Astra (OpenAI), ChatGPT Pro mode, effort Pro, S6\_10\_PAPERS continuation, 19 November 2026: used and rechecked in this run; status-vocabulary reconciliation; metadata reconciliation only; no new mathematical check. \emph{Pre-continuation recorded proof scope.} GPT-6 Astra (OpenAI, ChatGPT Pro mode), effort Pro, 17 November 2026; also checked in this writing by GPT-6 Astra (OpenAI), ChatGPT Pro mode, effort Pro, 18 November 2026}{the hypotheses, notation and definitions stated here; the explicit constructions and calculations in the proof}

Define the orthogonal projection and its projected period matrix by
\[
 \rho(x_1,x_2,x_3,x_4)=(x_2,x_4),\qquad
 Y_T=\begin{pmatrix}6m_T&L_T\\-q_T&m_T\end{pmatrix}.
\]
Then $\det Y_T=D_T>0$, so $Q_T=\mathbb R^2/Y_T\mathbb Z^2$ is a
compact real torus, and projection induces an exact sequence of compact
real Lie groups
\[
 0\longrightarrow
 \mathbb T^2_{\mathrm{re}}=
 \frac{\mathbb Re_1\oplus\mathbb Re_3}
      {\mathbb Ze_1\oplus\mathbb Ze_3}
 \longrightarrow F_T\xrightarrow{\ \pi_T\ }Q_T
 \longrightarrow0.
\]
This is a Riemannian quotient for the displayed flat metrics.  There is a
constant $C$, independent of $T\geq T_0$, such that
\[
 \|Y_T(a,b)^{\mathsf T}\|\geq(T-C)\sqrt{a^2+b^2},\qquad
 \operatorname{sys}(Q_T)\geq T-C.
\]
For every $R>0$, the radius-$R$ balls of $Q_T$ and $Q_T^2$ are exactly
Euclidean once $T-C>4R$.  Consequently their pointed limits are
$\mathbb R^2$ and $\mathbb R^4$, respectively, both locally isometrically
and smoothly, and in pointed Gromov--Hausdorff topology.
\end{proposition}



\begin{proof}
\begin{enumerate}[label=\arabic*.,leftmargin=*]
\item The determinant calculation gives
$\det Y_T=6m_T^2+L_Tq_T=D_T$.  Projection kills the last two period
columns and sends the first two to those of $Y_T$.  Therefore
$\rho(P_T\mathbb Z^4)=Y_T\mathbb Z^2$, which proves that $\pi_T$ is
well defined and surjective.

\item If $\pi_T([x])=0$, choose $a,b\in\mathbb Z$ such that
$\rho(x)=Y_T(a,b)^{\mathsf T}$.  The vector
$x-av_{\gamma,T}-bv_{u,T}$ belongs to
$\mathbb Re_1\oplus\mathbb Re_3$ and represents the same point of
$F_T$.  Conversely every point in this real plane maps to zero.
If a lattice vector lies in that plane, its first two lattice
coefficients satisfy $Y_T(a,b)^{\mathsf T}=0$.  Invertibility forces
$a=b=0$, so the intersection lattice is exactly
$\mathbb Ze_1\oplus\mathbb Ze_3$.  This proves the asserted kernel and
its injective inclusion.  Since $\rho$ is the orthogonal projection
along that plane, its Euclidean quotient metric is the displayed one.

\item Write
\[
 Y_T=T\begin{pmatrix}0&1\\-1&0\end{pmatrix}
       +\begin{pmatrix}6m_T&H_T\\B_T-H_T&m_T\end{pmatrix}.
\]
The second matrix has uniformly bounded operator norm by holomorphy.
Choose $C$ above this bound.  The first matrix is $T$ times an orthogonal
matrix, and the reverse triangle inequality gives the stated estimate.
A nonzero integer vector has norm at least one, which proves the
systolic bound; it also proves the same bound for the product lattice.

\item For $x,y$ in the Euclidean radius-$R$ ball and a nonzero lattice vector
$\lambda$, the estimate
\[
 \|x-y+\lambda\|\geq\|\lambda\|-\|x-y\|
 \geq(T-C)-2R>2R\geq\|x-y\|
\]
holds when $T-C>4R$.  Thus no lattice translation shortens a distance
between points of this ball.  Every point of the quotient ball has a
lift in the Euclidean ball, by choosing a shortest representative of
its lattice coset.  The covering map therefore gives an isometry onto
the entire quotient ball.  These are smooth local isometries, proving
all the stated pointed convergences.
\end{enumerate}
\end{proof}

\begin{theorem}[The maximal regular-locus
\texorpdfstring{\(\mathbb R^4\)}{R4} quotient and its compactification
obstruction]
\label{wb:cusp-spec-globalization-obstruction}
\StatementRecord{wb:cusp-spec-globalization-obstruction}{proved}{Historical events: GPT-6 Astra (OpenAI), ChatGPT Pro mode, effort Pro, S6\_09\_PAPERS, 17 November 2026: targeted proof check. Also checked in this writing. Pre-continuation writing events: GPT-6 Astra (OpenAI), ChatGPT Pro mode, effort Pro, S6\_10\_PAPERS, 18 November 2026: targeted proof check. Continuation metadata events: GPT-6 Astra (OpenAI), ChatGPT Pro mode, effort Pro, S6\_10\_PAPERS continuation, 19 November 2026: used and rechecked in this run; status-vocabulary reconciliation; metadata reconciliation only; no new mathematical check. \emph{Pre-continuation recorded proof scope.} GPT-6 Astra (OpenAI, ChatGPT Pro mode), S6\_09\_PAPERS, effort Pro, 2026-11-17; also checked in this writing by GPT-6 Astra (OpenAI), ChatGPT Pro mode, effort Pro, 18 November 2026}{the hypotheses, notation and definitions stated here; \cref{wb:cusp-spec-native-quotient}; \cref{wb:cusp-cone-S12-route}; \cref{wb:cusp-cone-S24-local-route}; \cref{wb:cusp-cone-S24-route}; \cite[\S15, Proposition~15.1 and the following remark,
p.~434]{BorelSerre1953}; \cref{wb:cusp-spec-invariant-character-lattice}; \cref{wb:cusp-spec-full-inverse-transpose}; \cref{wb:cusp-spec-radial-notation}; \cref{wb:cusp-spec-holomorphic-errors}; \cref{wb:cusp-spec-invariant-flag-calculation}; \cref{wb:cusp-spec-finite-overlap-identity}; \cref{wb:cusp-spec-cusp-nonextension-chart}}

\mbox{}\ 
Let
\[
 \Lambda=\mathbb Z\langle\widehat\gamma,\widehat u,
                   \widehat w,\widehat\delta\rangle,\qquad
 K_\Lambda=\mathbb Z\langle\widehat w,\widehat\delta\rangle,
\]
and define
\[
 \varepsilon_\gamma
 (a\widehat\gamma+b\widehat u+c\widehat w+d\widehat\delta)=a,\qquad
 L_\Lambda=\ker\varepsilon_\gamma
 =\mathbb Z\langle\widehat u,\widehat w,\widehat\delta\rangle.
\tag{G.1}
\]
Then the following statements hold.

\smallskip
\noindent
\emph{(i) Regular group quotient.}
Over \(B^\circ\), the period family \(J=X^\circ=f^{-1}(B^\circ)\) is a
holomorphic family of complex Lie groups, not merely an affine torus
bundle.  The formula
\begin{equation}
 \chi_\gamma\bigl([z,\Pi(z)x]\bigr)
 =\bigl([z],\varepsilon_\gamma(x)\bmod\mathbb Z\bigr)
 \label{wb:cusp-spec-global-circle-map}
\end{equation}
defines a global smooth real-analytic fibrewise group epimorphism
\begin{equation}
 \chi_\gamma:J\longrightarrow B^\circ\times(\mathbb R/\mathbb Z).
 \label{wb:cusp-spec-global-circle-quotient}
\end{equation}
Write
\begin{equation}
 \theta_\gamma:=\operatorname{pr}_{\mathbb R/\mathbb Z}\circ\chi_\gamma:
 J\longrightarrow\mathbb R/\mathbb Z
 \label{wb:cusp-spec-global-circle-coordinate}
\end{equation}
for its circle coordinate.  Thus
\(\chi_\gamma=(f,\theta_\gamma)\); the two symbols have different
codomains and will not be interchanged below.
Its kernel is the connected real three-torus family
\begin{equation}
 \mathcal H:=
 \frac{\mathfrak h_z^\circ\times
       ((L_\Lambda\otimes\mathbb R)/L_\Lambda)}{\Delta},
 \label{wb:cusp-spec-global-kernel-family}
\end{equation}
% checked carrier correction: deck-group
where \(\Delta\) is the actual triangle deck group of
\(\mathfrak h_z^\circ\to B^\circ\).  It acts diagonally by its deck action
and by the restricted integral monodromy on \(L_\Lambda\).
The deck representation maps \(\pi_1(B^\circ)\) onto \(\Delta\);
these two groups are not identified.  Equivalently, one can replace
\(\mathfrak h_z^\circ\) by the universal cover of \(B^\circ\) and
take the quotient by \(\pi_1(B^\circ)\).
The invariant-lattice calculation uses the image of the monodromy
representation, while finite base covers correspond to finite-index
subgroups of \(\pi_1(B^\circ)\).  Thus every fibre of \(\mathcal H\) is
\((L_\Lambda\otimes\mathbb R)/L_\Lambda\); the family need not be a
product.
The quotient is real and is not a holomorphic quotient of complex tori.

\smallskip
\noindent
\emph{(ii) Exact radial quotient metric.}
On the positive radial cusp, the restriction of
\cref{wb:cusp-spec-global-circle-map} to
\(F_T=\mathbb R^4/P_T\mathbb Z^4\) is
\begin{equation}
 \chi_{\gamma,T}([y])
 =\left\langle P_T^{-\mathsf T}e_\gamma,y\right\rangle
   \bmod\mathbb Z,\qquad e_\gamma=(1,0,0,0)^{\mathsf T}.
 \label{wb:cusp-spec-radial-circle-character}
\end{equation}
For the Riemannian quotient of the Euclidean metric, the target circle has
circumference
\begin{equation}
 \ell_T
 =\frac1{\lVert P_T^{-\mathsf T}e_\gamma\rVert}
 =\frac{D_T}{\sqrt{L_T^2+m_T^2}}
 =T+H-B+O(T^{-1}).
 \label{wb:cusp-spec-global-circle-circumference}
\end{equation}
In particular, \(\ell_T\to\infty\).

\smallskip
\noindent
\emph{(iii) Four copies and minimality.}
The source
\begin{equation}
 J^{[4]}=J\times_{B^\circ}J\times_{B^\circ}J\times_{B^\circ}J
 \label{wb:cusp-spec-fourfold-fibre-product}
\end{equation}
is the actual fourfold fibre product over all of \(B^\circ\).  The product
map
\begin{equation}
 \chi_\gamma^{[4]}=(\chi_\gamma,\chi_\gamma,\chi_\gamma,\chi_\gamma):
 J^{[4]}\longrightarrow B^\circ\times(\mathbb R/\mathbb Z)^4
 \label{wb:cusp-spec-fourfold-circle-map}
\end{equation}
has kernel
\(\mathcal H^{[4]}=\mathcal H\times_{B^\circ}\cdots
\times_{B^\circ}\mathcal H\), whose fibres are
\(\mathbb T^{12}\).  Along the synchronized positive
radial cusp, its quotient fibre, with product quotient metric, is
\begin{equation}
 G_T=(\mathbb R/\ell_T\mathbb Z)^4,\qquad
 \operatorname{sys}(G_T)=\ell_T.
 \label{wb:cusp-spec-four-circle-fibre}
\end{equation}
Consequently
\begin{equation}
 (G_T,0)\longrightarrow(\mathbb R^4,0)
 \label{wb:cusp-spec-four-circle-limit}
\end{equation}
in pointed local-isometric, pointed smooth, and pointed
Gromov--Hausdorff topology.  If \(\ell_T>4R\), the radius-\(R\) ball
about zero is exactly Euclidean.

Equivalently, the external Cartesian product gives a quotient
\[
 J^4\longrightarrow(B^\circ)^4\times(\mathbb R/\mathbb Z)^4,
\]
and \(J^{[4]}\) is its pullback along the diagonal
\(B^\circ\to(B^\circ)^4\).  At four independent cusp parameters its target
lattice is
\(\operatorname{diag}(\ell_{T_1},\ldots,\ell_{T_4})\mathbb Z^4\), so the
same pointed limit holds whenever
\(\min_i\ell_{T_i}\to\infty\).

Among connected torus quotients obtained from \(r\) copies by a primitive
full-monodromy-invariant kernel and whose quotient systole tends to infinity
on the radial cusp, the quotient rank is at most \(r\).  Thus four copies
are minimal for a four-dimensional all-expanding quotient, and
\cref{wb:cusp-spec-fourfold-circle-map} attains the bound.

The cusp-local two-copy map
\begin{equation}
 \pi_T\times\pi_T:F_T\times F_T\longrightarrow Q_T\times Q_T
 \label{wb:cusp-spec-doubled-local-quotient}
\end{equation}
from \cref{wb:cusp-spec-native-quotient} also has pointed limit
\(\mathbb R^4\).  It is not a quotient over the full \(B^\circ\), however,
and no finite base cover makes it one.

\smallskip
\noindent
\emph{(iv) Finite fillings.}
At \(p_j\), use the notation of
\cref{wb:prop:candidate-finite-fillings} and put
\[
 (m_1,v_1,\kappa_1)
 =\bigl(3,\widehat\gamma+2\widehat u-4\widehat w,1\bigr),
\]
\[
 (m_2,v_2,\kappa_2)
 =\bigl(4,-\widehat\gamma-3\widehat u+3\widehat w,-1\bigr),\qquad
 \kappa_j=\varepsilon_\gamma(v_j).
\tag{G.2}
\]
In the un-conjugated flat coordinates, the finite filling generator is
\begin{equation}
 \Phi_j(s,x)=
 \left(e^{-2\pi i/m_j}s,A_jx+\frac{v_j}{m_j}\right).
 \label{wb:cusp-spec-finite-affine-action}
\end{equation}
The raw character is not invariant:
\begin{equation}
 \varepsilon_\gamma\left(A_jx+\frac{v_j}{m_j}\right)
 =\varepsilon_\gamma(x)+\frac{\kappa_j}{m_j}\pmod{\mathbb Z}.
 \label{wb:cusp-spec-finite-character-shift}
\end{equation}
It instead descends equivariantly to the Seifert circle family
\begin{equation}
 C_j=
 \frac{\Delta_j\times(\mathbb R/\mathbb Z)}
 {(s,\theta)\sim
  (e^{-2\pi i/m_j}s,\theta+\kappa_j/m_j)}.
 \label{wb:cusp-spec-finite-Seifert-target}
\end{equation}
The total space \(C_j\) is a smooth solid torus, the map
\(C_j\to D_j\) has a multiple central circle of order \(m_j\), and the
induced map \(N_j\to C_j\) is a smooth proper \(\mathbb T^3\)-bundle.
The target translation in \cref{wb:cusp-spec-finite-Seifert-target} is
exactly compatible with the logarithmic gluing map \(h_j\).

For four synchronized copies, the smooth simultaneous-root extension is
\begin{equation}
 \widetilde N_j^{[4]}
 =\frac{(J'_j)^{\times_{\Delta_j}4}}
 {\langle(\widetilde g_j^{\log})^{\times4}\rangle},
 \label{wb:cusp-spec-finite-fourfold-normalization}
\end{equation}
with target
\begin{equation}
 C_j^{[4]}=
 \frac{\Delta_j\times(\mathbb R/\mathbb Z)^4}
 {(s,\theta_1,\ldots,\theta_4)\sim
  (e^{-2\pi i/m_j}s,
   \theta_1+\kappa_j/m_j,\ldots,\theta_4+\kappa_j/m_j)}.
 \label{wb:cusp-spec-finite-fourfold-target}
\end{equation}
The map \(\widetilde N_j^{[4]}\to C_j^{[4]}\) is a smooth proper
\(\mathbb T^{12}\)-bundle.  There is a canonical finite map
\begin{equation}
 \nu_j:\widetilde N_j^{[4]}
 \longrightarrow N_j\times_{D_j}N_j\times_{D_j}N_j\times_{D_j}N_j
 \label{wb:cusp-spec-finite-normalization-map}
\end{equation}
which is an isomorphism over \(D_j^*\), is the normalization of the
reduced ordinary filled fibre product, and has degree \(m_j^3\) on the
reduced central fibre.  These two degrees are \(27\) and \(64\).  Gluing
these two local normalizations to \(J^{[4]}\) by
\((h_j)^{\times4}\), and gluing their targets to the product
\(B^\circ\times(\mathbb R/\mathbb Z)^4\) by
\((h_j^{S^1})^{\times4}\), gives the corresponding normalization and
Seifert target over \(B\setminus\{p_0\}\).

\smallskip
\noindent
\emph{(v) Cusp nonextension.}
Neither the literal map \(\chi_\gamma\) nor
\(\chi_\gamma^{[4]}\) extends continuously through the toric cusp fibre
\(W\).  Thus the four-copy quotient is global over \(B^\circ\), and its
finite-filling normalization is global over \(B\setminus\{p_0\}\), but
the displayed map \(\chi_\gamma^{[4]}\) has no continuous extension to
the compact fourfold fibre product or to \(X^4\) with the stated
regular-fibre restriction, and it does not factor through
\(X^{\wedge4}\cong S^{24}\).  The cohomology class of
\(\chi_\gamma\) does extend over
\(N_0\); this weaker homotopy statement does not extend the displayed
quotient map or its metric.

\smallskip
\noindent
\emph{(vi) Products, smashes, and sphere constructions.}
The homotopy-sphere recognition in \cref{maps:src:B.5.2}, together with $\Theta_6=0$ \cite[Theorem~5.1 and table, p.~512]{manKM63}, gives a based diffeomorphism $h:X\to S^6$. Its Cartesian powers give
\begin{equation}
 X^2\cong_{\mathrm{diff}}S^6\times S^6\not\simeq S^{12},\qquad
 H_6(X^2;\mathbb Z)\cong\mathbb Z^2,
 \label{wb:cusp-spec-product-not-sphere}
\end{equation}
and
\begin{equation}
 \begin{aligned}
 X^4&\cong_{\mathrm{diff}}(S^6)^4\not\simeq S^{24},\\
 H_6(X^4;\mathbb Z)&\cong\mathbb Z^4,&
 H_{12}(X^4;\mathbb Z)&\cong\mathbb Z^6,&
 H_{18}(X^4;\mathbb Z)&\cong\mathbb Z^4.
 \end{aligned}
 \label{wb:cusp-spec-four-product-not-sphere}
\end{equation}
The synchronized regular fourfold \(J^{[4]}\) has real dimension \(18\),
not \(24\).  Choosing a base point gives based homeomorphisms
\begin{equation}
X\wedge X\cong S^{12},\qquad X^{\wedge4}\cong S^{24},
\label{wb:cusp-spec-smash-spheres}
\end{equation}
but these smash spheres do not carry either \emph{literal fixed-target}
quotient above.  The later cone constructions
\cref{wb:cusp-cone-S12-route} and \cref{wb:cusp-cone-S24-local-route} and \cref{wb:cusp-cone-S24-route}
change the codomain by retaining the radial parameter and collapsing the
obstructing boundary cross-section; they therefore do not contradict this
nonextension theorem.  Moreover,
neither \(S^{12}\) nor \(S^{24}\), with any differentiable structure,
admits an almost-complex structure
\cite[\S15, Proposition~15.1 and the following remark,
p.~434]{BorelSerre1953}; hence neither can be promoted to a compact
complex manifold.  This is an obstruction to a proposed complex
\(S^{12}\)- or \(S^{24}\)-globalization, not to the real quotient or its
pointed metric limit.  Precisely,
\(\pi_T\times\pi_T\) is a map on \(F_T^2\), and it fails to factor
through \(F_T\wedge F_T\); \(\chi_{\gamma,T}^{[4]}\) is a map on
\(F_T^4\), and it fails to factor through \(F_T^{\wedge4}\).
Neither map is a map on \(X^2\) or \(X^4\), and no extension with the
stated fibre restriction exists.

There is one exact but different sphere construction.  If
\(C_T=\mathbb R/\ell_T\mathbb Z\) is based at zero, postcomposition with
the target fat-wedge collapse gives
\begin{equation}
 F_T^4\xrightarrow{\chi_{\gamma,T}^{[4]}}C_T^4
 \longrightarrow C_T^{\wedge4}\cong S^4.
 \label{wb:cusp-spec-target-smash-map}
\end{equation}
This map does factor through \(F_T^{\wedge4}\).  With the quotient length
metric and base point
\([\ell_T/2,\ell_T/2,\ell_T/2,\ell_T/2]\), its pointed limit is
\(\mathbb R^4\), with every radius-\(R\) ball exactly Euclidean once
\(\ell_T>4R\).  The metric is singular at the collapsed fat wedge, and
the map still does not extend to \(X^4\) or \(X^{\wedge4}\).
\end{theorem}



\begin{proof}
\begin{enumerate}[label=\arabic*.,leftmargin=*]
\item \medskip\noindent\emph{Invariant lattices and the regular group law.}\par
The matrices in \cref{wb:prop:monodromy-certificate} give
\[
 (M_0-I)(a,b,c,d)^{\mathsf T}=(0,0,b,-a)^{\mathsf T},
\]
so
\begin{equation}
 K_\Lambda=\ker(M_0-I)=\operatorname{im}(M_0-I).
 \label{wb:cusp-spec-cusp-kernel}
\end{equation}
They also give
\begin{equation}
 \begin{aligned}
 A_1\widehat u&=-\widehat w+\widehat\delta,&
 A_1\widehat w&=\widehat u-\widehat w,&
 A_1\widehat\delta&=\widehat\delta,\\
 A_2\widehat u&=\widehat w,&
 A_2\widehat w&=-\widehat u+\widehat\delta,&
 A_2\widehat\delta&=\widehat\delta.
 \end{aligned}
 \label{wb:cusp-spec-invariant-flag-calculation}
\end{equation}
Thus \(L_\Lambda\) is invariant.  It is the smallest invariant lattice
containing \(K_\Lambda\), because
\(\widehat u=A_1\widehat w+\widehat w\).  Both lattices are primitive.

\item For a row covector \((a,b,c,d)\), invariance under \(A_1\) gives
\(b=2c\) and \(d=3c\), whereas invariance under \(A_2\) gives
\(b=c\) and \(d=2c\).  Hence \(b=c=d=0\), and
\begin{equation}
 \operatorname{Hom}(\Lambda,\mathbb Z)^\Delta
 =\mathbb Z\varepsilon_\gamma.
 \label{wb:cusp-spec-invariant-character-lattice}
\end{equation}

\item The family over the cover is
\((\mathfrak h_z^\circ\times\mathbb C^2)/\Lambda\), and the action used
to form \(J\) is
\[
 [z,\zeta]\longmapsto[gz,R_g(z)\zeta].
\]
It is complex linear in every fibre and fixes the zero section.
Fibrewise addition, inversion, and zero are therefore equivariant and
descend to \(J\).  In the flat real-analytic marking
\[
 \Xi(z,x+\Lambda)=[z,\Pi(z)x],
\]
equation (2.3d) makes the action
\[
 (z,x+\Lambda)\longmapsto(gz,A_gx+\Lambda).
\]
Equation \cref{wb:cusp-spec-invariant-character-lattice} proves that
\cref{wb:cusp-spec-global-circle-map} is well defined.  Primitivity of
\(\varepsilon_\gamma\) proves surjectivity and connectedness of the
kernel.  The induced action on \(\Lambda/L_\Lambda\simeq\mathbb Z\) is
trivial, so the target is the product circle bundle.  Since the kernel has
odd real dimension, this is not a holomorphic quotient.

\item \medskip\noindent\emph{The quotient metric and the four-copy limit.}\par
The first column of
\cref{wb:cusp-spec-full-inverse-transpose} is
\begin{equation}
 P_T^{-\mathsf T}e_\gamma
 =\left(0,\frac{m_T}{D_T},0,-\frac{L_T}{D_T}\right)^{\mathsf T}.
 \label{wb:cusp-spec-circle-dual-vector}
\end{equation}
The quotient metric on \(\mathbb R/\mathbb Z\) is therefore
\(\ell_T^2d\theta^2\).  Equivalently,
\[
 P_T(L_\Lambda\otimes\mathbb R)
 =\operatorname{span}_{\mathbb R}\{e_1,e_3,v_{u,T}\}
\]
has unit normal
\[
 \nu_T=\frac{(0,m_T,0,-L_T)}{\sqrt{L_T^2+m_T^2}},
\]
and
\[
 \operatorname{dist}
 \bigl(v_{\gamma,T},P_T(L_\Lambda\otimes\mathbb R)\bigr)
 =\left|\langle v_{\gamma,T},\nu_T\rangle\right|
 =\frac{L_Tq_T+6m_T^2}{\sqrt{L_T^2+m_T^2}}
 =\ell_T.
\]
The expansions in
\cref{wb:cusp-spec-radial-notation} and \cref{wb:cusp-spec-holomorphic-errors} give
\[
 D_T=T^2+(2H-B)T+O(1),\qquad
 \sqrt{L_T^2+m_T^2}=T+H+O(T^{-1}),
\]
which proves \cref{wb:cusp-spec-global-circle-circumference}.

\item After rescaling the circle coordinate by \(\ell_T\), the product target is
\(\mathbb R^4/\ell_T\mathbb Z^4\).  If
\(x,y\in B_R(0)\) and \(0\ne\lambda\in\ell_T\mathbb Z^4\), then
\[
 \lVert x-y+\lambda\rVert
 \ge\lVert\lambda\rVert-\lVert x-y\rVert
 \ge\ell_T-2R>2R\ge\lVert x-y\rVert
\]
when \(\ell_T>4R\).  Hence the covering projection is an exact isometry
on \(B_R(0)\), proving all three convergences in
\cref{wb:cusp-spec-four-circle-limit}.  The quotient map is global over
\(B^\circ\); only the metric limit is evaluated along the radial cusp.
The separately chosen Euclidean fibre metric is not claimed to be a
canonical global metric away from that cusp.

\item \medskip\noindent\emph{Minimality and the finite-base-cover obstruction.}\par
Let \(N\subset\Lambda^{\oplus r}\) be a primitive invariant kernel whose
connected quotient has radial systole tending to infinity.  Every element
of \(K_\Lambda^{\oplus r}\) has uniformly bounded physical length, since
\[
 P_T\widehat w=e_1,\qquad P_T\widehat\delta=e_3.
\]
If such an element survived in the quotient lattice, its quotient norm
would be no greater than its source norm and would give a nonzero uniformly
bounded period.  Therefore \(K_\Lambda^{\oplus r}\subset N\).
Invariance and
\cref{wb:cusp-spec-invariant-flag-calculation} then force
\(L_\Lambda^{\oplus r}\subset N\), and
\[
 \operatorname{rank}(\Lambda^{\oplus r}/N)\le4r-3r=r.
\]
This proves the lower bound \(r\ge4\), while
\(N=L_\Lambda^{\oplus4}\) proves sharpness.

\item For the finite-cover assertion, put
\begin{equation}
 H_*=A_2A_1
 =\begin{pmatrix}
 1&0&0&0\\
 6&1&1&0\\
 0&0&1&0\\
 -5&0&-1&1
 \end{pmatrix}=I+N_*.
 \label{wb:cusp-spec-no-finite-cover-matrix}
\end{equation}
Direct multiplication gives \(N_*^2=0\), and hence
\begin{equation}
 H_*^n\widehat w
 =\widehat w+n(\widehat u-\widehat\delta),\qquad
 H_*^n\widehat\delta=\widehat\delta.
 \label{wb:cusp-spec-no-finite-cover-orbit}
\end{equation}
Thus \(H_*^nK_\Lambda\ne K_\Lambda\) for every \(n\ne0\).
Every finite-index subgroup of the base fundamental group contains a
positive power of the loop with monodromy \(H_*\).  Such a subgroup cannot
preserve \(K_\Lambda\), proving that no finite base cover globalizes the
native two-torus quotient.

\item \medskip\noindent\emph{The affine finite fillings and the actual overlap
maps.}\par
Equation \cref{wb:cusp-spec-invariant-character-lattice} gives
\(\varepsilon_\gamma A_j=\varepsilon_\gamma\), so
\cref{wb:cusp-spec-finite-character-shift} follows.  Define
\[
 \overline\Phi_j(s,\theta)
 =\left(e^{-2\pi i/m_j}s,\theta+\frac{\kappa_j}{m_j}\right).
\]
Then \((s,x)\mapsto(s,\varepsilon_\gamma(x))\) is
\((\Phi_j,\overline\Phi_j)\)-equivariant; denote the descended map by
\(\overline\chi_j:N_j\to C_j\).  A nonidentity power which fixes
the base must lie over \(s=0\), where it translates the circle by
\(k\kappa_j/m_j\).  Since \(\gcd(\kappa_j,m_j)=1\), this is nonintegral
for \(0<k<m_j\).  The target action is free; the source action is free by
\cref{wb:prop:candidate-finite-fillings}.

\item For an explicit solid-torus chart, put
\(\lambda_j=e^{2\pi i/m_j}\), \(u=e^{2\pi i\theta}\), and choose
\(r_j\in\mathbb Z\) with \(r_j\kappa_j\equiv1\pmod {m_j}\).  Then
\[
 w=su^{r_j},\qquad U=u^{m_j}
\]
are invariant and identify \(C_j\) real-analytically with a solid torus.
In these coordinates
\[
 t_j=s^{m_j}=w^{m_j}U^{-r_j},
\]
which proves the central multiplicity.  Equivariance and the primitive
kernel give the proper \(\mathbb T^3\)-bundle \(N_j\to C_j\).

\item It remains to check the actual logarithmic gluing.  Let
\(\Pi_{\mathbb R}(z):\Lambda\otimes\mathbb R\to\mathbb C^2\) be the
underlying real isomorphism.  On a logarithm branch \(\ell\), define
\(b_j^\ell(z)\) uniquely by
\begin{equation}
 \Pi_{\mathbb R}(z)b_j^\ell(z)
 =\frac{\ell(z)}{2\pi i}\Pi(z)v_j,\qquad
 \beta_j(z)=\varepsilon_\gamma(b_j^\ell(z))\bmod\mathbb Z.
 \label{wb:cusp-spec-finite-target-conjugacy}
\end{equation}
Changing \(\ell\) by \(2\pi in\) adds \(nv_j\) to \(b_j^\ell\), hence
adds \(n\kappa_j\in\mathbb Z\) to its
\(\varepsilon_\gamma\)-value.  Thus \(\beta_j\) is branch independent.
In flat coordinates the source conjugacy (2.9e) is translation by
\(b_j^\ell\).  Applying \(\varepsilon_\gamma\) to that conjugacy gives
\begin{equation}
 \beta_j(g_jz)=\beta_j(z)-\frac{\kappa_j}{m_j}\pmod{\mathbb Z}.
 \label{wb:cusp-spec-finite-target-cocycle}
\end{equation}
Therefore
\[
 h_j^{S^1}(z,\theta)=(z,\theta+\beta_j(z))
\]
conjugates \(\overline\Phi_j\) to the linear target action
\((z,\theta)\mapsto(g_jz,\theta)\), and the overlap identity is
\begin{equation}
 \chi_\gamma\circ h_j=h_j^{S^1}\circ\overline\chi_j.
 \label{wb:cusp-spec-finite-overlap-identity}
\end{equation}

\item The diagonal action on
\((J'_j)^{\times_{\Delta_j}4}\) is free, and the same coprimality argument
proves that \(C_j^{[4]}\) is smooth.  Every fibre of the resulting map is
a \(\mathbb T^{12}\)-torsor; over \(D_j^*\) its structure torus is the
fourfold fibre product of \(\mathcal H\).  No group zero is asserted at
the affine filling.
In local etale coordinates along the central fibre, the ordinary filled
fibre product has equations
\begin{equation}
 s_i^{m_j}=s_1^{m_j},\qquad i=2,3,4.
 \label{wb:cusp-spec-finite-fibre-product-branches}
\end{equation}
It has \(m_j^3\) smooth branches along the reduced central fibre.  The
simultaneous-root model separates them, is smooth and hence normal, and is
an isomorphism away from the central fibre.  The map \(\nu_j\) is proper:
both spaces are proper over \(D_j\), and the inverse image of a compact
subset of the target is closed in the inverse image of its compact image
in \(D_j\).  It is quasi-finite, since its fibres are singletons away from
the central fibre and have at most \(m_j^3\) points over the central
fibre.  Hence \(\nu_j\) is finite.  The universal property of
normalization now proves that it is the normalization of the reduced
ordinary fibre product.  On the central fibre the relative root choices form
\((\mathbb Z/m_j)^4/\operatorname{diag}(\mathbb Z/m_j)\), of order
\(m_j^3\), which proves the asserted degrees.

\item On \(D_j^*\), the source identification is \((h_j)^{\times4}\), and
\cref{wb:cusp-spec-finite-overlap-identity} identifies the target overlap
with \((h_j^{S^1})^{\times4}\).  These are precisely the fourfolds of
the one-copy overlap maps and therefore satisfy the cocycle identities.
They glue the two local normalizations and Seifert targets to
\(J^{[4]}\) and to
\(B^\circ\times(\mathbb R/\mathbb Z)^4\), respectively.  Equivalently,
the source gluing follows from uniqueness of normalization.

\item The fixed product-target trivialization does not extend over \(X\), even if
base winding is allowed.  Indeed, in the presentation of
\cref{mapsI:top:prop:geometric-pi1}, let a homomorphism to \(\mathbb Z\) restrict to
\(n\varepsilon_\gamma\) on \(\Lambda\) and take values \(a,b\) on
\(\widehat\rho_1,\widehat\rho_2\).  The three attachment relations give
\[
 3a=n,\qquad4b=-n,\qquad a+b=0.
\]
The first two give \(a=n/3\), \(b=-n/4\), while the third gives
\(n/12=0\).  Hence \(n=0\).  This is the exact global obstruction behind
the fact that the two Seifert targets do not combine into a fixed-circle
map on \(X\).

\item \medskip\noindent\emph{The explicit toric cusp obstruction.}\par
On the positive radial ray, write
\[
 Z_T=iTB_0+C(t_c),\qquad
 B_0=\begin{pmatrix}0&1\\-1&0\end{pmatrix},\qquad
 t_c=e^{-2\pi T}.
\]
For \(a=(a_1,a_2)^{\mathsf T}\in\mathbb R^2\), take the point of the
regular fibre represented by \(Z_Ta\), with zero short-period coordinates,
and exponentiate it to the dense torus of the cusp model:
\[
 x_T(a)=\exp(2\pi iZ_Ta)\in(\mathbb C^*)^2.
\]
Its logarithmic position is
\begin{equation}
 y_T(a)=
 \frac{(\log|x_{T,1}(a)|,\log|x_{T,2}(a)|)}{\log|t_c|}
 =B_0a+\frac{\operatorname{Im}C(t_c)a}{T}.
 \label{wb:cusp-spec-cusp-nonextension-position}
\end{equation}
The circle character has the constant value
\begin{equation}
 \chi_{\gamma,T}(x_T(a))=a_1\bmod\mathbb Z.
 \label{wb:cusp-spec-cusp-nonextension-character}
\end{equation}

\item Use the affine toric chart of
\(\sigma=\operatorname{conv}(0,e_1,e_2)\).  Its coordinates are
\begin{equation}
 z_0=\frac{t_c}{x_1x_2},\qquad z_1=x_1,\qquad z_2=x_2.
 \label{wb:cusp-spec-cusp-nonextension-chart}
\end{equation}
If \(B_0a=(y_1,y_2)\) is in the interior of \(\sigma\), then
\(y_1>0\), \(y_2>0\), and \(1-y_1-y_2>0\).  For all sufficiently large
\(T\), the three absolute values in
\cref{wb:cusp-spec-cusp-nonextension-chart} are
\[
 |t_c|^{1-y_{T,1}-y_{T,2}},\qquad
 |t_c|^{y_{T,1}},\qquad |t_c|^{y_{T,2}},
\]
and all tend to zero.  Every such path converges to the same toric fixed
point and hence to the same triple point of \(W\).

\item Take
\[
 a=\left(-\frac14,\frac14\right),\qquad
 a'=\left(-\frac12,\frac14\right).
\]
Then
\[
 B_0a=\left(\frac14,\frac14\right),\qquad
 B_0a'=\left(\frac14,\frac12\right)
\]
are both in the interior of \(\sigma\).  The paths converge to the same
point of \(W\), whereas
\[
 \chi_{\gamma,T}(x_T(a))=\frac34,\qquad
 \chi_{\gamma,T}(x_T(a'))=\frac12.
\]
No continuous extension of the literal map \(\chi_\gamma\) exists.
Choosing identical marked radial paths in the other three factors and
using these two paths in one factor proves the same assertion for
\(\chi_\gamma^{[4]}\).

\item The radial family of translations by \(\widehat u/2\) in \(\mathcal H\)
fails to extend as well.  Take two paths
whose limiting logarithmic positions are
\[
 y=\left(\frac18,\frac18\right),\qquad
 y'=\left(\frac38,\frac38\right).
\]
Both approach the fixed point of
\(\operatorname{conv}(0,e_1,e_2)\).  Translation by
\(\widehat u/2\in\mathcal H_T\) adds the exact displacement
\[
 \frac12B_0e_2+
 \frac1{2T}\operatorname{Im}C(t_c)e_2
 \longrightarrow\left(\frac12,0\right)
\]
to logarithmic position.  The translated limits are
\[
 \left(\frac58,\frac18\right)
 \quad\hbox{and}\quad
 \left(\frac78,\frac38\right).
\]
The first lies in the interior of
\(\operatorname{conv}(0,e_1,e_2)\), while the second lies in the interior
of \(\operatorname{conv}(e_1,e_2,e_1+e_2)\).  These two triangle types are
the two translation orbits in the \(A_2\) triangulation and give the two
distinct triple points of \(W\), as proved in
\cref{wb:thm:candidate-cusp-geometry}.  Thus one regular-fibre translation
sends two paths with the same limit to paths with different limits; it has
no continuous extension to \(N_0\).

\item This pointwise obstruction must not be confused with the cohomology class.
By \cref{mapsI:top:prop:geometric-pi1},
\[
 \pi_1(N_0)=\Lambda/K_\Lambda,
\]
and \(\varepsilon_\gamma\) kills \(K_\Lambda\).  It therefore defines an
integral \(H^1\)-class on \(N_0\), hence a map to the Eilenberg--Mac Lane
space \(S^1\) whose restriction is homotopic to \(\chi_\gamma\).  The
explicit paths prove that no such representative agrees with
\(\chi_\gamma\) on the whole punctured cusp.  The phase torus
\((K_\Lambda\otimes\mathbb R)/K_\Lambda\) extends as the compact real torus
action of the toric model, generally with boundary stabilizers; the full
radial \(\widehat u/2\)-translation and the full quotient do not.

\item \medskip\noindent\emph{Products and smash quotients.}\par
The integral Kunneth theorem gives
\[
 H_k(X^n;\mathbb Z)\cong
 \begin{cases}
  \mathbb Z^{\binom nj},&k=6j,\quad0\le j\le n,\\
  0,&\text{otherwise},
 \end{cases}
 \qquad \chi(X^n)=2^n,
\]
which proves
\cref{wb:cusp-spec-product-not-sphere} and \cref{wb:cusp-spec-four-product-not-sphere}.
The synchronized regular \(n\)-fold fibre product has dimension
\[
 6n-2(n-1)=4n+2,
\]
which gives \(18\) for four copies.

\item A product map factors through a smash quotient if and only if it is
constant on the corresponding fat wedge.  On
\(F_T\times\{0\}\), the local two-copy quotient is
\[
 x\longmapsto(\pi_T(x),0),
\]
which is nonconstant.  Thus the literal conclusion is that
\(\pi_T\times\pi_T\) does not factor through \(F_T\wedge F_T\); it would
be a type error to call it a map on \(X\wedge X\).  The restriction of
\(\chi_{\gamma,T}^{[4]}\) to \(F_T\times\{0\}^3\) is similarly
\[
 x\longmapsto(\chi_{\gamma,T}(x),0,0,0),
\]
so it does not factor through \(F_T^{\wedge4}\).

\item There is no global extension which repairs the domain mismatch:
\(H^1(X^n;\mathbb Z)=0\), whereas both local quotient maps pull nonzero
target \(H^1\)-classes back to nonzero classes on the cusp fibre.  Hence no
map on \(X^2\) or \(X^4\) has the required restriction, and precomposition
with the smash collapses in \cref{wb:cusp-spec-smash-spheres} cannot produce
the local quotient maps.

\item Identify the metric quotient circle with
\(C_T=\mathbb R/\ell_T\mathbb Z\) by
\(\theta\mapsto\ell_T\theta\).  For
\cref{wb:cusp-spec-target-smash-map}, target-smash collapse is constant
whenever one source coordinate is the identity, so the modified map does
factor through \(F_T^{\wedge4}\).  Its chosen base point has distance
\(\ell_T/2\) from the collapsed target fat wedge.  If two points lie
within distance \(R\) of it, a route through the collapsed point has
length at least \(\ell_T-2R\), while their direct Euclidean distance is at
most \(2R\).  When \(\ell_T>4R\), neither the collapse nor circle
wrap-around shortens distances in that ball.

\item Let \(u_i\in H^1(F_T^4;\mathbb Z)\) be the pullback of the generator from
the \(i\)-th quotient circle.  The orientation class of
\(S^4=C_T^{\wedge4}\) pulls back under
\cref{wb:cusp-spec-target-smash-map} to
\[
 \pm u_1\smile u_2\smile u_3\smile u_4\ne0.
\]
But \(H^4(X^4;\mathbb Z)=0\).  Thus even this modified \(S^4\)-target map
does not extend to \(X^4\), and a fortiori does not arise from a map on
\(X^{\wedge4}\cong S^{24}\).

\item For completeness, Borel--Serre prove that \(S^{2n}\) admits no
almost-complex structure for \(n\ge4\).  Their argument applies to every
differentiable structure on the sphere: on \(S^{2n}\), all lower Chern
classes vanish, while their reduced-power calculation forces the top
Chern class to vanish modulo a suitable odd prime \(p<n\); on the other
hand that class is the Euler class, namely \(2h\) for a generator
\(h\in H^{2n}(S^{2n};\mathbb Z)\), a contradiction.  Taking \(n=6\) and
\(n=12\) gives the asserted obstructions for \(S^{12}\) and \(S^{24}\).
This topological obstruction does not itself prohibit the real
four-circle target \(G_T\), the target-smash \(S^4\), or their metric
limits.

\item Finally,
\(X*X\cong S^{13}\) and \(X^{*4}\cong S^{27}\), so joins do not produce
the proposed \(S^{12}\) or \(S^{24}\) dimensions.  The based suspension map $\Sigma^k h:\Sigma^k X\to\Sigma^k S^6\cong S^{6+k}$ gives $S^{12}$ at $k=6$ and $S^{24}$ at $k=18$. The displayed fixed-target fibre quotients are obstructed from extending by the preceding cohomology calculations; the subsequent cone constructions give the specified comparison maps with modified targets.
\end{enumerate}
\end{proof}

\subsection{Cone-compactified sphere carriers for the native cusp quotient}
\label{wb:cusp-cone-section}

This subsection answers a topological question which is strictly different
from the finite-cover construction.  No period is enlarged by passing to a
cover.  Instead, the native expanding quotient is retained on the cusp, the
parts on which that quotient is unavailable are collapsed, and the resulting
based map is passed through a smash product.  The outcome is an exact global
map with source \(S^{12}\), and there is an analogous \(S^{24}\) map from the
full-monodromy-invariant circle quotient.  The targets are conically
stratified spaces.  The Euclidean limits below use the intrinsic metrics
inherited from the flat quotient tori on the regular slices; they are not
ambient metric limits in one fixed metric
on the compact targets.

Fix a closed coordinate disc
\begin{equation}
 D_0=\{t\in\mathbb C:|t|\leq\epsilon\}\subset B,\qquad
 p_0=\{t=0\},\qquad N_0=f^{-1}(D_0),\qquad W=f^{-1}(0),
 \label{wb:cusp-cone-disc}
\end{equation}
small enough that it contains no singular value other than \(p_0\).  Put
\begin{equation}
 K_\Lambda=\mathbb Z\langle\widehat w,\widehat\delta\rangle,\qquad
 \overline\Lambda=\Lambda/K_\Lambda,\qquad
 \overline T^2=
 \frac{\overline\Lambda\otimes\mathbb R}{\overline\Lambda}.
 \label{wb:cusp-cone-quotient-lattice}
\end{equation}
The ordered basis of \(\overline\Lambda\) is the image of
\((\widehat\gamma,\widehat u)\).

\begin{lemma}[The native quotient over the whole punctured cusp]
\label{wb:cusp-cone-native-descent}
\StatementRecord{wb:cusp-cone-native-descent}{proved}{Historical events: GPT-6 Astra (OpenAI), ChatGPT Pro mode, effort Pro, S6\_09\_PAPERS, 17 November 2026: targeted proof check. Also checked in this writing. Pre-continuation writing events: GPT-6 Astra (OpenAI), ChatGPT Pro mode, effort Pro, S6\_10\_PAPERS, 18 November 2026: targeted proof check. Continuation metadata events: GPT-6 Astra (OpenAI), ChatGPT Pro mode, effort Pro, S6\_10\_PAPERS continuation, 19 November 2026: used and rechecked in this run; status-vocabulary reconciliation; metadata reconciliation only; no new mathematical check. \emph{Pre-continuation recorded proof scope.} GPT-6 Astra (OpenAI, ChatGPT Pro mode), S6\_09\_PAPERS, effort Pro, 2026-11-17; also checked in this writing by GPT-6 Astra (OpenAI), ChatGPT Pro mode, effort Pro, 18 November 2026}{the hypotheses, notation and definitions stated here; \cref{wb:cusp-spec-native-quotient}}

There is a well-defined real-analytic fibrewise group epimorphism
\begin{equation}
 \overline\pi:N_0\setminus W
 \longrightarrow D_0^*\times\overline T^2
 \label{wb:cusp-cone-native-family-map}
\end{equation}
whose kernel is the connected two-torus family defined by \(K_\Lambda\).
The target local system is trivial.  On the positive radial fibre \(F_T\),
\(t=e^{-2\pi T}\), the isomorphism
\begin{equation}
 \iota_T:\overline T^2\longrightarrow Q_T,\qquad
 [(a,b)]\longmapsto[Y_T(a,b)^{\mathsf T}]
 \label{wb:cusp-cone-radial-isometry}
\end{equation}
satisfies
\begin{equation}
 \pi_T=\iota_T\circ\overline\pi|_{F_T}.
 \label{wb:cusp-cone-native-restriction}
\end{equation}
If \(\overline T^2\) is given the pullback by \(\iota_T\) of the
Euclidean quotient metric on \(Q_T=\mathbb R^2/Y_T\mathbb Z^2\), then
\(\iota_T\) is an isometry and
\begin{equation}
 \operatorname{sys}(\overline T^2,g_T)
 =\operatorname{sys}(Q_T)\geq T-C.
 \label{wb:cusp-cone-native-systole}
\end{equation}
\end{lemma}



\begin{proof}
\begin{enumerate}[label=\arabic*.,leftmargin=*]
\item On the universal cover of \(D_0^*\), use the real-analytic flat marking
\[
 (z,x+\Lambda)\longmapsto[z,\Pi(z)x].
\]
The deck generator acts by
\[
 (z,x+\Lambda)\longmapsto(g_0z,M_0x+\Lambda).
\]
The displayed matrix \(M_0\) gives
\begin{equation}
 (M_0-I)(a,b,c,d)^{\mathsf T}=(0,0,b,-a)^{\mathsf T},\qquad
 \operatorname{im}(M_0-I)=K_\Lambda.
 \label{wb:cusp-cone-monodromy-image}
\end{equation}
Consequently \(M_0\) preserves \(K_\Lambda\) and induces the identity on
\(\overline\Lambda\).  The formula
\[
 [z,\Pi(z)x]\longmapsto
 \bigl([z],[x\bmod K_\Lambda]\bigr)
\]
is invariant under the deck action and proves
\cref{wb:cusp-cone-native-family-map}.  The lattice \(K_\Lambda\) is primitive,
so the kernel is connected.  On the radial fibre, projection to the
\((e_2,e_4)\)-coordinates sends
\(a\widehat\gamma+b\widehat u\) to \(Y_T(a,b)^{\mathsf T}\) and kills
\(K_\Lambda\).  This proves
\cref{wb:cusp-cone-radial-isometry} and \cref{wb:cusp-cone-native-restriction}.  The metric
and systolic assertion is the quotient calculation in
\cref{wb:cusp-spec-native-quotient}.
\end{enumerate}
\end{proof}

Define
\begin{equation}
 \mathcal A=
 \frac{[0,\epsilon]\times S^1\times\overline T^2}
 {\{0\}\times S^1\times\overline T^2\sim v_0,\quad
  \{\epsilon\}\times S^1\times\overline T^2\sim v_1}.
 \label{wb:cusp-cone-A-definition}
\end{equation}
Thus
\begin{equation}
 \mathcal A\cong
 \operatorname{Susp}(S^1\times\overline T^2)
 \cong\operatorname{Susp}(T^3),
 \label{wb:cusp-cone-A-suspension}
\end{equation}
where \(\operatorname{Susp}\) denotes the unreduced suspension.  The extra
circle is the angular coordinate of the punctured base.  If one deliberately
forgets that angle, the evident map
\(\mathcal A\to\operatorname{Susp}(\overline T^2)\) gives a smaller target
with the same restriction to every fixed radial fibre.  We retain
\(\mathcal A\) because it records the entire punctured-cusp quotient.

\begin{theorem}[Exact localized \(S^{12}\) route]
\label{wb:cusp-cone-S12-route}
\StatementRecord{wb:cusp-cone-S12-route}{proved}{Historical events: GPT-6 Astra (OpenAI), ChatGPT Pro mode, effort Pro, S6\_09\_PAPERS, 17 November 2026: targeted proof check. Also checked in this writing. Pre-continuation writing events: GPT-6 Astra (OpenAI), ChatGPT Pro mode, effort Pro, S6\_10\_PAPERS, 18 November 2026: targeted proof check. Continuation metadata events: GPT-6 Astra (OpenAI), ChatGPT Pro mode, effort Pro, S6\_10\_PAPERS continuation, 19 November 2026: used and rechecked in this run; status-vocabulary reconciliation; metadata reconciliation only; no new mathematical check. \emph{Pre-continuation recorded proof scope.} GPT-6 Astra (OpenAI, ChatGPT Pro mode), S6\_09\_PAPERS, effort Pro, 2026-11-17; also checked in this writing by GPT-6 Astra (OpenAI), ChatGPT Pro mode, effort Pro, 18 November 2026}{the hypotheses, notation and definitions stated here; \cref{wb:cusp-cone-A-definition}; \cref{wb:cusp-cone-A-suspension}; \cref{wb:cusp-cone-native-descent}; \cref{wb:cusp-cone-native-family-map}; \cite[Chapter~0, Smash Product]{Hatcher2002}; \cref{wb:cusp-cone-native-restriction}}

Choose \(x_0\in W\), and base \(\mathcal A\) at \(v_0\).  There is a
continuous surjective based quotient map
\begin{equation}
 q_2:(X,x_0)\longrightarrow(\mathcal A,v_0)
 \label{wb:cusp-cone-q2}
\end{equation}
with the following exact description.  For \(x\in N_0\setminus W\), write
\[
 t(f(x))=re^{i\phi},\qquad 0<r<\epsilon,
\]
and let \(\overline\pi_f(x)\in\overline T^2\) be the fibre component of
\(\overline\pi(x)\).  Then
\begin{equation}
 q_2(x)=[r,e^{i\phi},\overline\pi_f(x)],\qquad
 q_2(W)=v_0,\qquad
 q_2\bigl(X\setminus\operatorname{int}N_0\bigr)=v_1.
 \label{wb:cusp-cone-q2-formula}
\end{equation}
On \(f^{-1}(|t|=\epsilon)\), the value is \(v_1\).

The smash square is a global based map
\begin{equation}
 q_2^{\wedge2}:X\wedge X\longrightarrow
 \mathcal A\wedge\mathcal A.
 \label{wb:cusp-cone-S12-map}
\end{equation}
Using \(X\cong_{\mathrm{diff}}S^6\), a based choice of the diffeomorphism
identifies its source with
\begin{equation}
 X\wedge X\cong S^6\wedge S^6\cong S^{12}.
 \label{wb:cusp-cone-S12-source}
\end{equation}
For \(0<e^{-2\pi T}<\epsilon\), the canonical product-to-smash quotient
maps restrict to embeddings on \(F_T^2\) and on the corresponding interior
target-product cross-section.  Under those embeddings and \(\iota_T^2\),
the induced smash map restricts exactly to
\begin{equation}
 \pi_T\times\pi_T:F_T\times F_T\longrightarrow Q_T\times Q_T.
 \label{wb:cusp-cone-S12-fibre-map}
\end{equation}
With the Euclidean \(\ell^2\)-product flat quotient metric,
\begin{equation}
 \operatorname{sys}(Q_T\times Q_T)\geq T-C,\qquad
 (Q_T\times Q_T,(0,0))\longrightarrow(\mathbb R^4,0)
 \label{wb:cusp-cone-S12-R4-limit}
\end{equation}
in pointed local-isometric, pointed smooth, and pointed
Gromov--Hausdorff topology.  Once \(T-C>4R\), the radius-\(R\) ball is
exactly Euclidean.

The integral homology of the targets is
\begin{equation}
 \widetilde H_k(\mathcal A;\mathbb Z)\cong
 \begin{cases}
  \mathbb Z^3,&k=2,3,\\
  \mathbb Z,&k=4,\\
  0,&\text{otherwise},
 \end{cases}
 \label{wb:cusp-cone-A-homology}
\end{equation}
and
\begin{equation}
 \widetilde H_k(\mathcal A\wedge\mathcal A;\mathbb Z)\cong
 \begin{cases}
  \mathbb Z^9,&k=4,\\
  \mathbb Z^{18},&k=5,\\
  \mathbb Z^{15},&k=6,\\
  \mathbb Z^6,&k=7,\\
  \mathbb Z,&k=8,\\
  0,&\text{otherwise}.
 \end{cases}
 \label{wb:cusp-cone-A-smash-homology}
\end{equation}
In particular, the target is not \(S^4\), \(S^{12}\), or a compact
four-manifold.
\end{theorem}



\begin{proof}
\begin{enumerate}[label=\arabic*.,leftmargin=*]
\item The equivalence relation in \cref{wb:cusp-cone-A-definition} is closed: it is
the union of the diagonal with the squares of the two closed end
cross-sections.  A compact Hausdorff space modulo a closed equivalence
relation is compact Hausdorff.  Rescaling \(r\) to the unit interval gives
exactly the unreduced suspension in \cref{wb:cusp-cone-A-suspension}.

\item Well-definedness of the interior formula follows from
\cref{wb:cusp-cone-native-descent}.  A full turn in \(\phi\) changes the
marking by \(M_0\), but \(M_0\) is the identity modulo \(K_\Lambda\).

\item We check both collapsed strata.  If \(U\) is an open neighborhood of
\(v_0\), its inverse image in the cylinder is open and contains the compact
set \(\{0\}\times S^1\times\overline T^2\).  Hence, for some \(\delta>0\),
it contains
\([0,\delta)\times S^1\times\overline T^2\).  Therefore
\(q_2^{-1}(U)\) contains \(f^{-1}(|t|<\delta)\), a neighborhood of \(W\).
The same compactness argument at the other end shows that every neighborhood
of \(v_1\) contains the image of
\((\epsilon-\delta,\epsilon]\times S^1\times\overline T^2\).  Its preimage
together with \(X\setminus\operatorname{int}N_0\) is a neighborhood of the
outer boundary.  Continuity on the interior is
\cref{wb:cusp-cone-native-family-map}; the closed pasting lemma gives
continuity on \(X\).

\item Every interior point has a preimage because \(\overline\pi\) is fibrewise
surjective.  The nonempty sets \(W\) and
\(X\setminus\operatorname{int}N_0\) map onto \(v_0\) and \(v_1\),
respectively.  Thus \(q_2\) is surjective.  A continuous surjection from
compact \(X\) to Hausdorff \(\mathcal A\) is a quotient map.

\item The map is based, so functoriality gives \cref{wb:cusp-cone-S12-map}.  The
standard CW quotient gives \(S^m\wedge S^n\cong S^{m+n}\)
\cite[Chapter~0, Smash Product]{Hatcher2002}, proving
\cref{wb:cusp-cone-S12-source}.

\item The base point lies in \(W\), so \(F_T^2\) is disjoint from the source fat
wedge.  Its image lies in an interior cross-section of
\(\mathcal A^2\), disjoint from the target fat wedge.  The source and target
fat wedges are closed in the compact Hausdorff products \(X^2\) and
\(\mathcal A^2\), so both smash quotients are Hausdorff.  Their restrictions
to the relevant compact subspaces are continuous and one-to-one, and are
therefore embeddings.
Equations \cref{wb:cusp-cone-native-restriction} and \cref{wb:cusp-cone-q2-formula} give
\cref{wb:cusp-cone-S12-fibre-map}.

\item The period lattice of the product target is
\(\operatorname{diag}(Y_T,Y_T)\mathbb Z^4\).  Its systole is the systole of
\(Y_T\mathbb Z^2\), proving the first assertion in
\cref{wb:cusp-cone-S12-R4-limit}.  If \(x,y\in B_R(0)\subset\mathbb R^4\)
and \(0\ne\lambda\) is a product period, then
\[
 \|x-y+\lambda\|\geq\|\lambda\|-\|x-y\|
 \geq(T-C)-2R>2R\geq\|x-y\|
\]
when \(T-C>4R\).  The quotient projection is an exact isometry on the
Euclidean radius-\(R\) ball.  Conversely, any point of quotient distance
less than \(R\) has a shortest lattice representative of norm less than
\(R\), so this embedded ball is the entire quotient radius-\(R\) ball.
This proves all three pointed convergences.

\item Reduced suspension shifts reduced homology by one.  Since
\[
 H_1(T^3)=\mathbb Z^3,\qquad
 H_2(T^3)=\mathbb Z^3,\qquad H_3(T^3)=\mathbb Z,
\]
\cref{wb:cusp-cone-A-homology} follows.  All groups are free, so the Kunneth
theorem for smash products has no Tor terms.  The ranks in degrees
\(4,5,6,7,8\) are
\[
 3\cdot3,\quad 2(3\cdot3),\quad
 2(3\cdot1)+3\cdot3,\quad 2(3\cdot1),\quad1,
\]
which proves \cref{wb:cusp-cone-A-smash-homology}.
\end{enumerate}
\end{proof}

\paragraph{The localized one-circle construction.}
The full-monodromy-invariant character also gives a smaller cusp-supported
target.  Define
\begin{equation}
 \mathcal B=
 \frac{[0,\epsilon]\times S^1\times S^1}
 {\{0\}\times S^1\times S^1\sim b_0,\quad
  \{\epsilon\}\times S^1\times S^1\sim b_1}
 \cong\operatorname{Susp}(T^2),
 \label{wb:cusp-cone-B-definition}
\end{equation}
where the two circle coordinates in an interior cross-section are the base
angle and the value of \(\theta_\gamma\).

\begin{proposition}[Localized \(S^{24}\) route with target
\(\operatorname{Susp}(T^2)^{\wedge4}\)]
\label{wb:cusp-cone-S24-local-route}
\StatementRecord{wb:cusp-cone-S24-local-route}{proved}{Historical events: GPT-6 Astra (OpenAI), ChatGPT Pro mode, effort Pro, S6\_09\_PAPERS, 17 November 2026: targeted proof check. Also checked in this writing. Pre-continuation writing events: GPT-6 Astra (OpenAI), ChatGPT Pro mode, effort Pro, S6\_10\_PAPERS, 18 November 2026: targeted proof check. Continuation metadata events: GPT-6 Astra (OpenAI), ChatGPT Pro mode, effort Pro, S6\_10\_PAPERS continuation, 19 November 2026: used and rechecked in this run; status-vocabulary reconciliation; metadata reconciliation only; no new mathematical check. \emph{Pre-continuation recorded proof scope.} GPT-6 Astra (OpenAI, ChatGPT Pro mode), S6\_09\_PAPERS, effort Pro, 2026-11-17; also checked in this writing by GPT-6 Astra (OpenAI), ChatGPT Pro mode, effort Pro, 18 November 2026}{the hypotheses, notation and definitions stated here; \cref{wb:cusp-cone-S24-R4-limit}; \cref{wb:cusp-cone-S12-route}}

There is a continuous surjective based quotient map
\begin{equation}
 q_{1,\mathrm{cusp}}:(X,x_0)\longrightarrow(\mathcal B,b_0)
 \label{wb:cusp-cone-q1-local}
\end{equation}
given by
\begin{equation}
 q_{1,\mathrm{cusp}}(x)
 =[r,e^{i\phi},\theta_\gamma(x)]
 \quad(0<|t|=r<\epsilon),\qquad
 q_{1,\mathrm{cusp}}(W)=b_0,\qquad
 q_{1,\mathrm{cusp}}
 \bigl(X\setminus\operatorname{int}N_0\bigr)=b_1.
 \label{wb:cusp-cone-q1-local-formula}
\end{equation}
Its fourth smash power has source \(X^{\wedge4}\cong S^{24}\), and on
four synchronized radial fibres it is exactly
\begin{equation}
 \chi_{\gamma,T}^{\times4}:F_T^4\longrightarrow C_T^4.
 \label{wb:cusp-cone-S24-local-fibre-map}
\end{equation}
Consequently its intrinsic target slices satisfy the Euclidean limit
\cref{wb:cusp-cone-S24-R4-limit}.  The target homology is
\begin{equation}
 \widetilde H_k(\mathcal B;\mathbb Z)\cong
 \begin{cases}
  \mathbb Z^2,&k=2,\\
  \mathbb Z,&k=3,\\
  0,&\text{otherwise},
 \end{cases}
 \label{wb:cusp-cone-B-homology}
\end{equation}
and
\begin{equation}
 \widetilde H_k(\mathcal B^{\wedge4};\mathbb Z)\cong
 \begin{cases}
  \mathbb Z^{16},&k=8,\\
  \mathbb Z^{32},&k=9,\\
  \mathbb Z^{24},&k=10,\\
  \mathbb Z^8,&k=11,\\
  \mathbb Z,&k=12,\\
  0,&\text{otherwise}.
 \end{cases}
 \label{wb:cusp-cone-B-smash-homology}
\end{equation}
\end{proposition}



\begin{proof}
\begin{enumerate}[label=\arabic*.,leftmargin=*]
\item The monodromy invariance
\(\varepsilon_\gamma M_0=\varepsilon_\gamma\) makes the interior formula
single-valued after a turn in the base angle.  The two endpoint-continuity
arguments are identical to those in the proof of
\cref{wb:cusp-cone-S12-route}, with compact cross-section \(T^2\) in place of
\(T^3\).  Fibrewise surjectivity of \(\theta_\gamma\), together with the
nonempty inner fibre and exterior, proves surjectivity.  Compact-to-Hausdorff
makes the map a quotient map.

\item Since \(x_0\in W\), the fourth smash power exists.  The radial regular
fibres and their circle targets avoid the two closed fat wedges.  The
restrictions of the Hausdorff smash quotients to these compact subspaces are
continuous and one-to-one, hence embeddings, and the fibre map is exactly
\cref{wb:cusp-cone-S24-local-fibre-map}.  The intrinsic metric statement is
therefore the four-circle calculation
\cref{wb:cusp-cone-S24-R4-limit}.  Finally, suspension shifts the free homology
of \(T^2\), giving \cref{wb:cusp-cone-B-homology}.  In the fourth smash power,
choosing \(j\) of the four factors in degree \(3\) and the other
\(4-j\) factors in degree \(2\) gives degree \(8+j\) and rank
\(\binom4j2^{4-j}\).  This is \cref{wb:cusp-cone-B-smash-homology}.
\end{enumerate}
\end{proof}

\paragraph{The full-base circle target before coning.}
Choose pairwise disjoint closed coordinate discs \(D_0,D_1,D_2\) about the
three singular values and put
\begin{equation}
 P=B\setminus\bigcup_{j=0}^2\operatorname{int}D_j.
 \label{wb:cusp-cone-pair-of-pants}
\end{equation}
For \(j=1,2\), let \(C_j\) be the solid-torus Seifert target in
\cref{wb:cusp-spec-finite-Seifert-target}, and let \(h_j^{S^1}\) be the
target overlap map in \cref{wb:cusp-spec-finite-overlap-identity}.  Glue
\begin{equation}
 \mathcal M=(P\times S^1)
 \underset{h_1^{S^1}}{\cup}C_1
 \underset{h_2^{S^1}}{\cup}C_2.
 \label{wb:cusp-cone-M-definition}
\end{equation}
The gluing is along boundary tori over \(\partial D_1,\partial D_2\).
The remaining boundary is
\begin{equation}
 \partial\mathcal M=\partial D_0\times S^1\cong T^2.
 \label{wb:cusp-cone-M-boundary}
\end{equation}
Define
\begin{equation}
 \mathcal C=\mathcal M\cup_{\partial\mathcal M}
 \operatorname{Cone}(\partial\mathcal M),
 \label{wb:cusp-cone-C-definition}
\end{equation}
and write \(c_0\) for the cone vertex.
This changes the codomain at the cusp: the whole boundary two-torus,
including the target circle, is collapsed to \(c_0\).  It therefore does
not contradict the proved nonextension of the literal map
\(\chi_\gamma\) to the fixed product target \(B\times S^1\).

\begin{theorem}[Exact cone-compactified \(S^{24}\) route]
\label{wb:cusp-cone-S24-route}
\StatementRecord{wb:cusp-cone-S24-route}{proved}{Historical events: GPT-6 Astra (OpenAI), ChatGPT Pro mode, effort Pro, S6\_09\_PAPERS, 17 November 2026: targeted proof check. Also checked in this writing. Pre-continuation writing events: GPT-6 Astra (OpenAI), ChatGPT Pro mode, effort Pro, S6\_10\_PAPERS, 18 November 2026: targeted proof check. Continuation metadata events: GPT-6 Astra (OpenAI), ChatGPT Pro mode, effort Pro, S6\_10\_PAPERS continuation, 19 November 2026: used and rechecked in this run; status-vocabulary reconciliation; metadata reconciliation only; no new mathematical check. \emph{Pre-continuation recorded proof scope.} GPT-6 Astra (OpenAI, ChatGPT Pro mode), S6\_09\_PAPERS, effort Pro, 2026-11-17; also checked in this writing by GPT-6 Astra (OpenAI), ChatGPT Pro mode, effort Pro, 18 November 2026}{the hypotheses, notation and definitions stated here; \cref{wb:cusp-spec-finite-Seifert-target}; \cref{wb:cusp-spec-global-circle-circumference}; \cref{wb:cusp-cone-S12-route}}

There is a continuous surjective based quotient map
\begin{equation}
 q_1:(X,x_0)\longrightarrow(\mathcal C,c_0),\qquad x_0\in W,
 \label{wb:cusp-cone-q1}
\end{equation}
defined as follows:
\begin{enumerate}
\item over \(P\), it is
\((f,\theta_\gamma):f^{-1}(P)\to P\times S^1\), equivalently the
restriction of \(\chi_\gamma\) to \(f^{-1}(P)\);
\item over \(D_j\), \(j=1,2\), it is the finite-filling map
\(N_j\to C_j\) from \cref{wb:cusp-spec-finite-Seifert-target};
\item on \(N_0\setminus W\), writing \(t=re^{i\phi}\), it is
\begin{equation}
 x\longmapsto[r,(e^{i\phi},\theta_\gamma(x))]
 \in\operatorname{Cone}(\partial D_0\times S^1),
 \label{wb:cusp-cone-q1-cusp-formula}
\end{equation}
and every point of \(W\) maps to \(c_0\).
\end{enumerate}
The fourth smash power is a global based map
\begin{equation}
 q_1^{\wedge4}:X^{\wedge4}\longrightarrow\mathcal C^{\wedge4}.
 \label{wb:cusp-cone-S24-map}
\end{equation}
Using \(X\cong_{\mathrm{diff}}S^6\), its source is
\begin{equation}
 X^{\wedge4}\cong(S^6)^{\wedge4}\cong S^{24}.
 \label{wb:cusp-cone-S24-source}
\end{equation}
On four synchronized radial fibres, the canonical product-to-smash quotient
maps restrict to embeddings on the source fibre product and the interior
target-product cross-section.  Under those embeddings,
\cref{wb:cusp-cone-S24-map} restricts exactly to
\begin{equation}
 \chi_{\gamma,T}^{\times4}:F_T^4\longrightarrow C_T^4,\qquad
 C_T=\mathbb R/\ell_T\mathbb Z.
 \label{wb:cusp-cone-S24-fibre-map}
\end{equation}
With its Euclidean \(\ell^2\)-product flat quotient metric,
\begin{equation}
 \operatorname{sys}(C_T^4)=\ell_T
 =T+H-B+O(T^{-1}),\qquad
 (C_T^4,0)\longrightarrow(\mathbb R^4,0)
 \label{wb:cusp-cone-S24-R4-limit}
\end{equation}
in pointed local-isometric, pointed smooth, and pointed
Gromov--Hausdorff topology.

The target homology is
\begin{equation}
 \widetilde H_k(\mathcal C;\mathbb Z)\cong
 \begin{cases}
  \mathbb Z,&k=2,3,\\
  0,&\text{otherwise},
 \end{cases}
 \label{wb:cusp-cone-C-homology}
\end{equation}
and
\begin{equation}
 \widetilde H_k(\mathcal C^{\wedge4};\mathbb Z)\cong
 \begin{cases}
  \mathbb Z,&k=8,12,\\
  \mathbb Z^4,&k=9,11,\\
  \mathbb Z^6,&k=10,\\
  0,&\text{otherwise}.
 \end{cases}
 \label{wb:cusp-cone-C-smash-homology}
\end{equation}
\end{theorem}



\begin{proof}
\begin{enumerate}[label=\arabic*.,leftmargin=*]
\item The compact spaces \(P\times S^1,C_1,C_2\) are glued along closed boundary
tori by homeomorphisms, producing a compact orientable three-manifold
\(\mathcal M\) with boundary \(\partial D_0\times S^1\).  Attaching the cone
along this closed boundary gives a compact Hausdorff finite CW space
\(\mathcal C\).

\item The overlap identity
\[
 \chi_\gamma\circ h_j=h_j^{S^1}\circ\overline\chi_j
\]
proves that the first two pieces of \(q_1\) agree.  On \(\partial D_0\), the
cone base is glued to the remaining boundary of \(P\times S^1\), so the
first and third pieces agree.  If \(U\) is a neighborhood of \(c_0\), its
inverse image in the cone contains \([0,\delta)\times T^2\) for some
\(\delta>0\), by compactness of the cone base.  Therefore \(q_1^{-1}(U)\)
contains \(f^{-1}(|t|<\delta)\), a neighborhood of \(W\).  This proves
continuity at the new stratum; continuity elsewhere follows from the input
maps and the closed pasting lemma.

\item Every regular circle quotient and both finite Seifert targets are
surjective, and \(W\) maps onto the cone vertex.  Thus \(q_1\) is a
continuous surjection from compact \(X\) to Hausdorff \(\mathcal C\), hence
a quotient map.  It is based, so its fourth smash power exists.  Iterating
\(S^m\wedge S^n\cong S^{m+n}\) gives \cref{wb:cusp-cone-S24-source}.

\item The base point lies in \(W\), while each radial regular fibre is disjoint
from \(W\).  Hence \(F_T^4\) and its product-circle image avoid the closed
source and target fat wedges.  The canonical quotient maps from the source
product and the indicated interior target-product cross-section to their
smash products restrict to embeddings.  Under those embeddings, the induced
smash map restricts to \(\chi_{\gamma,T}^{\times4}\), which itself retains
its displayed nontrivial torus kernel.
The quotient-metric calculation in
\cref{wb:cusp-spec-global-circle-circumference} gives \(\ell_T\).
The product lattice is \(\ell_T\mathbb Z^4\), so the radius-\(R\) argument,
including the shortest-representative proof that the embedded Euclidean ball
is the whole quotient ball, in the proof of \cref{wb:cusp-cone-S12-route} proves
\cref{wb:cusp-cone-S24-R4-limit}.

\item For homology, let \(\rho_0,\rho_1,\rho_2\) be the boundary loops of the pair
of pants, with \(\rho_0+\rho_1+\rho_2=0\), and let \(\eta\) be the regular
target-circle class.  With the boundary-loop and circle orientations used
in the finite-filling presentation, the meridian relations are
\begin{equation}
 3\rho_1-\eta=0,\qquad4\rho_2+\eta=0.
 \label{wb:cusp-cone-Seifert-meridians}
\end{equation}
Reversing the circle orientation reverses both displayed \(\eta\)-signs
and changes none of the following groups.  Thus
\begin{equation}
 H_1(\mathcal M;\mathbb Z)
 =\frac{\mathbb Z\langle\eta,\rho_1,\rho_2\rangle}
 {\langle3\rho_1-\eta,\;4\rho_2+\eta\rangle}
 \cong\mathbb Z\langle g\rangle,
 \label{wb:cusp-cone-M-H1}
\end{equation}
where
\begin{equation}
 \rho_0=-g,\qquad\rho_1=4g,\qquad
 \rho_2=-3g,\qquad\eta=12g.
 \label{wb:cusp-cone-M-boundary-map}
\end{equation}
In particular,
\(H_1(\partial\mathcal M)=\mathbb Z\langle\rho_0,\eta\rangle
\to H_1(\mathcal M)\) is surjective.  The orientable three-manifold
\(\mathcal M\) has one torus boundary, \(H_3(\mathcal M)=0\), and
\(\chi(\mathcal M)=\frac12\chi(\partial\mathcal M)=0\); hence
\(H_2(\mathcal M;\mathbb Q)=0\).  To obtain the integral statement without
suppressing torsion, Poincare--Lefschetz duality and the universal
coefficient theorem give
\[
 H_2(\mathcal M;\mathbb Z)
 \cong H^1(\mathcal M,\partial\mathcal M;\mathbb Z)
 \cong \operatorname{Hom}
 \bigl(H_1(\mathcal M,\partial\mathcal M;\mathbb Z),\mathbb Z\bigr),
\]
because \(H_0(\mathcal M,\partial\mathcal M;\mathbb Z)=0\).  Thus
\(H_2(\mathcal M;\mathbb Z)\) is free; its rank is zero, so it vanishes.

\item The cone attachment is the homotopy cofiber of
\(\partial\mathcal M\hookrightarrow\mathcal M\).  Since the boundary
inclusion is a cofibration,
\[
 \widetilde H_k(\mathcal C;\mathbb Z)
 \cong H_k(\mathcal M,\partial\mathcal M;\mathbb Z).
\]
Poincare--Lefschetz duality and
\(H^0(\mathcal M)=H^1(\mathcal M)=\mathbb Z\),
\(H^2(\mathcal M)=H^3(\mathcal M)=0\), give
\cref{wb:cusp-cone-C-homology}.  The homology is free, so the smash Kunneth
theorem says that four factors of degree \(2\) or \(3\) produce ranks
\(\binom4j\) in degree \(8+j\), proving
\cref{wb:cusp-cone-C-smash-homology}.
\end{enumerate}
\end{proof}

\begin{definition}[The three radial target-slice embeddings]
\label{wb:cusp-cone-target-slice-embeddings}
\StatementRecord{wb:cusp-cone-target-slice-embeddings}{proved}{Historical events: GPT-6 Astra (OpenAI), ChatGPT Pro mode, effort Pro, S6\_09\_PAPERS, 17 November 2026: targeted proof check. Also checked in this writing. Pre-continuation writing events: GPT-6 Astra (OpenAI), ChatGPT Pro mode, effort Pro, S6\_10\_PAPERS, 18 November 2026: targeted proof check. Continuation metadata events: GPT-6 Astra (OpenAI), ChatGPT Pro mode, effort Pro, S6\_10\_PAPERS continuation, 19 November 2026: used and rechecked in this run; status-vocabulary reconciliation; metadata reconciliation only; no new mathematical check. \emph{Pre-continuation recorded proof scope.} GPT-6 Astra (OpenAI, ChatGPT Pro mode), S6\_09\_PAPERS, effort Pro, 2026-11-17; also checked in this writing by GPT-6 Astra (OpenAI), ChatGPT Pro mode, effort Pro, 18 November 2026}{the hypotheses, notation and definitions stated here; the argument and qualifications immediately following}

Put \(r_T=e^{-2\pi T}\) and let
\[
 \vartheta_T:\mathbb R/\mathbb Z\longrightarrow
 C_T=\mathbb R/\ell_T\mathbb Z,\qquad [a]\longmapsto[\ell_Ta].
\]
For \(y_1,y_2\in Q_T\), define
\[
 j_T(y_1,y_2)=
 \bigl[[r_T,1,\iota_T^{-1}(y_1)],
       [r_T,1,\iota_T^{-1}(y_2)]\bigr]_{\wedge}
 \in\mathcal A^{\wedge2}.
\]
For \(c_1,\ldots,c_4\in C_T\), define
\[
 k_T^{\mathcal B}(c_1,\ldots,c_4)=
 \bigl[[r_T,1,\vartheta_T^{-1}(c_1)],\ldots,
       [r_T,1,\vartheta_T^{-1}(c_4)]\bigr]_{\wedge}
 \in\mathcal B^{\wedge4},
\]
and, using the cusp-cone piece of \(\mathcal C\),
\[
 k_T^{\mathcal C}(c_1,\ldots,c_4)=
 \bigwedge_{j=1}^4[r_T,(1,\vartheta_T^{-1}(c_j))]
 \in\mathcal C^{\wedge4}.
\]
These are embeddings because every displayed factor lies away from the smash
basepoint, so equality in the smash quotient is equality in every factor;
compact-to-Hausdorff injectivity then gives a topological embedding.
\end{definition}



\begin{proposition}[Intrinsic expansion versus ambient collapse]
\label{wb:cusp-cone-metric-scope}
\StatementRecord{wb:cusp-cone-metric-scope}{proved}{Historical events: GPT-6 Astra (OpenAI), ChatGPT Pro mode, effort Pro, S6\_09\_PAPERS, 17 November 2026: targeted proof check. Also checked in this writing. Pre-continuation writing events: GPT-6 Astra (OpenAI), ChatGPT Pro mode, effort Pro, S6\_10\_PAPERS, 18 November 2026: targeted proof check. Continuation metadata events: GPT-6 Astra (OpenAI), ChatGPT Pro mode, effort Pro, S6\_10\_PAPERS continuation, 19 November 2026: used and rechecked in this run; status-vocabulary reconciliation; metadata reconciliation only; no new mathematical check. \emph{Pre-continuation recorded proof scope.} GPT-6 Astra (OpenAI, ChatGPT Pro mode), S6\_09\_PAPERS, effort Pro, 2026-11-17; also checked in this writing by GPT-6 Astra (OpenAI), ChatGPT Pro mode, effort Pro, 18 November 2026}{the hypotheses, notation and definitions stated here; \cref{wb:cusp-cone-S12-R4-limit}; \cref{wb:cusp-cone-S24-R4-limit}; \cref{wb:cusp-cone-native-systole}}

Let \(D_{12}\), \(D_{24}^{\mathcal B}\), and
\(D_{24}^{\mathcal C}\) be arbitrary fixed metrics compatible with the
topologies of \(\mathcal A^{\wedge2}\), \(\mathcal B^{\wedge4}\), and
\(\mathcal C^{\wedge4}\), respectively.  In their corresponding ambient
metrics, the radial target
slices satisfy
\begin{equation}
 \operatorname{diam}_{\mathrm{amb}}
 \bigl(j_T(Q_T^2)\bigr)\longrightarrow0,\qquad
 \operatorname{diam}_{\mathrm{amb}}
 \bigl(k_T^{\mathcal B}(C_T^4)\bigr)\longrightarrow0,\qquad
 \operatorname{diam}_{\mathrm{amb}}
 \bigl(k_T^{\mathcal C}(C_T^4)\bigr)\longrightarrow0.
 \label{wb:cusp-cone-ambient-collapse}
\end{equation}
By contrast, in the intrinsic quotient metrics,
\begin{equation}
 \operatorname{diam}(Q_T^2)\geq\tfrac12(T-C)\longrightarrow\infty,\qquad
 \operatorname{diam}(C_T^4)\geq\tfrac12\ell_T\longrightarrow\infty.
 \label{wb:cusp-cone-intrinsic-expansion}
\end{equation}
Thus \cref{wb:cusp-cone-S12-R4-limit} and \cref{wb:cusp-cone-S24-R4-limit} assert
pointed local-isometric limits for the specified flat quotient metrics on the varying slices. The identity from each such slice to its embedded image pulls back the fixed ambient metric to a different metric: its diameter tends to zero by the displayed calculation.
\end{proposition}



\begin{proof}
\begin{enumerate}[label=\arabic*.,leftmargin=*]
\item Every neighborhood of the relevant cone vertex contains every sufficiently
small radial cross-section, uniformly in all compact angular and torus
coordinates.  In a compatible metric, every such cross-section eventually
lies in an arbitrarily small ball about the vertex.  Products followed by
the smash quotient preserve this uniform convergence, proving
\cref{wb:cusp-cone-ambient-collapse}.

\item For a flat torus with systole \(s\), choose a shortest primitive period
\(\lambda\).  The class of \(\lambda/2\) has distance \(s/2\) from zero:
for every period \(\mu\), the nonzero period \(\lambda-2\mu\) has length at
least \(s\).  Hence the diameter is at least \(s/2\).
Equations \cref{wb:cusp-cone-native-systole} and \cref{wb:cusp-cone-S24-R4-limit} prove
\cref{wb:cusp-cone-intrinsic-expansion}.
\end{enumerate}
\end{proof}

The obstruction to an almost-complex structure on \(S^{12}\),
with all three complete classical proofs, is given in
\cref{wb:cusp-cone-S12-almost-complex-obstruction}.


\begin{theorem}[Sharp obstructions to stronger interpretations]
\label{wb:cusp-cone-obstructions}
\StatementRecord{wb:cusp-cone-obstructions}{proved}{Historical events: GPT-6 Astra (OpenAI), ChatGPT Pro mode, effort Pro, S6\_09\_PAPERS, 17 November 2026: targeted proof check. Also checked in this writing. Pre-continuation writing events: GPT-6 Astra (OpenAI), ChatGPT Pro mode, effort Pro, S6\_10\_PAPERS, 18 November 2026: targeted proof check. Continuation metadata events: GPT-6 Astra (OpenAI), ChatGPT Pro mode, effort Pro, S6\_10\_PAPERS continuation, 19 November 2026: used and rechecked in this run; status-vocabulary reconciliation; metadata reconciliation only; no new mathematical check. \emph{Pre-continuation recorded proof scope.} GPT-6 Astra (OpenAI, ChatGPT Pro mode), S6\_09\_PAPERS, effort Pro, 2026-11-17; also checked in this writing by GPT-6 Astra (OpenAI), ChatGPT Pro mode, effort Pro, 18 November 2026}{the hypotheses, notation and definitions stated here; \cref{wb:cusp-cone-S12-route}; \cref{wb:cusp-cone-S12-almost-complex-obstruction}; \cite[\S15, Proposition~15.1 and the following remark,
p.~434]{BorelSerre1953}; \cref{wb:cusp-cone-kernel-contained}; \cref{wb:cusp-cone-finite-index-orbit}}

The sphere routes above have the following exact limitations.
\begin{enumerate}
\item Let \(r\geq1\), and let \(N\subset\Lambda^{\oplus r}\) be a
primitive kernel defining a connected fibrewise torus quotient over the
full regular base, or over any finite cover of that base.  Equip the quotient
with the norm induced by the unrescaled marked Euclidean metric on the
original fibres.  If its quotient systole tends to infinity on the radial
cusp, then
\begin{equation}
 \operatorname{rank}(\Lambda^{\oplus r}/N)\leq r.
 \label{wb:cusp-cone-rank-bound}
\end{equation}
Thus two copies cannot give a full-base four-torus quotient.  The
\(S^{12}\) map in \cref{wb:cusp-cone-S12-route} evades the obstruction only by
localizing to the cusp and collapsing the exterior.

\item No nontrivial finite base cover, including one extending with
branching over the original singular values, makes \(K_\Lambda\) a
sub-local system over the full regular base.  The local cyclic change
\(t=s^M\) used for the auxiliary degree-\(M^2\) cover cannot extend as a
connected cover of \(B\cong\mathbb P^1\) branched only at \(p_0\).

\item There is no map from \(X^2\), \(X\wedge X\), \(X^4\), or
\(X^{\wedge4}\) to a fixed \(T^4\) whose restriction retains all four
target character classes of the displayed product quotient.  There is no
map to a fixed \(S^4\) whose restriction is a degree-one collapse of
\(Q_T^2\), or of the four-circle target, for even one regular fibre.  The
conical targets evade these obstructions because the fibre character and
orientation classes do not extend as global target classes.

\item Cartesian products do not produce the named spheres:
\begin{equation}
 X^2\simeq S^6\times S^6\not\simeq S^{12},\qquad
 X^4\simeq(S^6)^4\not\simeq S^{24}.
 \label{wb:cusp-cone-products-not-spheres}
\end{equation}
The two- and four-fold synchronized fibre products have real dimensions
\(10\) and \(18\), respectively, and generally have singular completed
fibres.  Joins give
\begin{equation}
 X*X\cong S^{13},\qquad X^{*4}\cong S^{27}.
 \label{wb:cusp-cone-join-dimensions}
\end{equation}

\item Since \(X\cong S^6\) is simply connected, every connected unbranched
finite cover of \(X\), \(X^r\), \(X\wedge X\cong S^{12}\), or
\(X^{\wedge4}\cong S^{24}\) is trivial.  The local degree-\(M^2\)
unspooling cover is not the restriction of a nontrivial connected
unbranched cover of any of these compact sources.  The finite flat normal covers of \cref{maps:src:C.1.5} use the explicit base cover and normalization recorded there; the present assertion concerns unbranched covers of the listed compact sources.

\item By \cref{wb:cusp-cone-S12-almost-complex-obstruction}, no smooth manifold
homeomorphic to \(S^{12}\) admits an almost-complex structure.  Nor does
\(S^{24}\), with any differentiable structure, by the Borel--Serre theorem
\cite[\S15, Proposition~15.1 and the following remark,
p.~434]{BorelSerre1953}.  Thus the smash sources cannot be promoted to
compact complex spheres.
\end{enumerate}
\end{theorem}



\begin{proof}
\begin{enumerate}[label=\arabic*.,leftmargin=*]
\item Every element of \(K_\Lambda^{\oplus r}\) has uniformly bounded physical
length on the radial cusp.  If one survived in the quotient lattice, it
would be a nonzero period of bounded quotient norm.  Hence
\begin{equation}
 K_\Lambda^{\oplus r}\subset N.
 \label{wb:cusp-cone-kernel-contained}
\end{equation}
For the full monodromy group,
\(A_1\widehat w=\widehat u-\widehat w\) and invariance force every copy of
\(\widehat u\) into \(N\).  Therefore
\[
 L_\Lambda^{\oplus r}
 =\mathbb Z\langle\widehat u,\widehat w,\widehat\delta\rangle^{\oplus r}
 \subset N,
\]
which proves \cref{wb:cusp-cone-rank-bound}.

\item The same conclusion holds after a finite base cover.  Let \(H\) be its
finite-index subgroup and put \(H_*=A_2A_1=I+N_*\), \(N_*^2=0\).
Every finite-index subgroup contains \(H_*^n\) for some \(n>0\), while
\begin{equation}
 H_*^n\widehat w
 =\widehat w+n(\widehat u-\widehat\delta).
 \label{wb:cusp-cone-finite-index-orbit}
\end{equation}
Equations
\cref{wb:cusp-cone-kernel-contained} and \cref{wb:cusp-cone-finite-index-orbit} give
\(n\widehat u\in N\).  Primitivity gives \(\widehat u\in N\), proving the
same rank bound.  The orbit also gives
\(H_*^nK_\Lambda\ne K_\Lambda\) for every \(n>0\), so no finite cover of
the regular base preserves the native kernel.  A compact cover branched
over the singular values restricts to such a cover.  Finally, a connected
finite cover of \(\mathbb P^1\) branched only at \(p_0\) would restrict to
a connected unbranched finite cover of
\(\mathbb P^1\setminus\{p_0\}\cong\mathbb C\), hence has degree one.

\item For a fixed \(T^4\) target, \(H^1(X^r;\mathbb Z)=0\), as does the first
cohomology of the smash spheres.  A global map therefore pulls every target
character class back to zero.  On a regular fibre the quotient homomorphisms
inject the target character lattice into the character lattice of the
source torus, giving four independent nonzero pullbacks, a contradiction.

\item For the fixed \(S^4\) obstruction, let
\(\omega_T\in H^2(Q_T;\mathbb Z)\) be the orientation class.  Since
\(\pi_T^*:H^1(Q_T)\to H^1(F_T)\) injects two independent classes and torus
cohomology is their exterior algebra,
\begin{equation}
 (\pi_T\times\pi_T)^*(\omega_T\times\omega_T)
 \ne0\quad\text{in }H^4(F_T^2;\mathbb Z).
 \label{wb:cusp-cone-fibre-H4-nonzero}
\end{equation}
A degree-one collapse \(Q_T^2\to S^4\) pulls the orientation generator
back to \(\pm(\omega_T\times\omega_T)\).  But
\[
 H^4(X^2;\mathbb Z)=H^4(S^{12};\mathbb Z)=0,
\]
so no global extension has this restriction.  The four-circle case is the
same: its orientation class pulls back to the exterior product of four
independent circle characters, while
\(H^4(X^4)=H^4(S^{24})=0\).

\item The Kunneth theorem gives \(H_6(X^2)=\mathbb Z^2\), whereas
\(H_6(S^{12})=0\), and gives nonzero groups of ranks \(4,6,4\) in degrees
\(6,12,18\) for \(X^4\), whereas \(S^{24}\) has none.  The real dimension
of an \(r\)-fold synchronized fibre product is
\(6r-2(r-1)=4r+2\), giving \(10\) and \(18\).  The join formula
\(S^m*S^n\cong S^{m+n+1}\) proves
\cref{wb:cusp-cone-join-dimensions}.  The remaining assertions follow from
simple connectivity, \cref{wb:cusp-cone-S12-almost-complex-obstruction}, and
the cited Borel--Serre theorem for \(S^{24}\).
\end{enumerate}
\end{proof}

\paragraph{Canonicity and exact scope.}
The lattice \(K_\Lambda=\ker(M_0-I)=\operatorname{im}(M_0-I)\) and the
punctured-cusp quotient are intrinsic to the marked cusp monodromy.
The compact map \(q_2\) additionally depends on the chosen cusp coordinate,
closed cusp disc, and the decision to collapse its outer boundary and the
exterior.  The identifications
\(X^{\wedge2}\cong S^{12}\) and \(X^{\wedge4}\cong S^{24}\) depend on a
based choice of \(X\cong S^6\).  None of these choices is claimed to be
holomorphic or canonical in the algebraic category.

The map \(q_2^{\wedge2}\) is the direct two-copy construction from the
native \(T^2\times\mathbb R^2\) cusp: it is globally defined on the smash
sphere, and its radial regular-fibre quotients are \(Q_T\times Q_T\), which
converge pointed local-isometrically, smoothly, and Gromov--Hausdorff to
\(\mathbb R^4\).  It is not a torus quotient over the full
regular base.  The map \(q_{1,\mathrm{cusp}}^{\wedge4}\) is the parallel
cusp-supported construction with target
\(\operatorname{Susp}(T^2)^{\wedge4}\).  The stronger map
\(q_1^{\wedge4}\) uses the quotient over every regular and finite-filled
fibre before the cusp cone is attached; four copies are minimal for that
full-base all-expanding group quotient.
Neither construction gives a compact complex \(S^{12}\) or \(S^{24}\), an
ambient Euclidean limit in a fixed compact metric target, or a claim about
the interacting Yang--Mills spectrum.

The argument does not rule out a separately constructed branched or
toroidal compactification of the auxiliary finite covers after additional
branching away from \(p_0\).  Such a construction would have to specify its
global base cover, normalize and resolve every pulled-back singular fibre,
track the lifted metric, and determine the topology of the resulting total
space.  No such smooth compact sphere construction follows from the present
data.


\subsubsection{Optimal sphere exponents and the exact Euclidean boundary}
\label{wb:cusp-cone-optimal-strengthening}

Within the class of native connected fibrewise torus quotients equipped with
the induced, unscaled quotient of the marked Euclidean metric, the two sphere
dimensions in the preceding constructions are forced by two different
integral quotient problems.  Cusp-locally, each copy has an expanding
quotient of rank two, so two copies are necessary and sufficient within that
class.  Over the full regular base, monodromy leaves at most one expanding
quotient direction per copy, so four copies are necessary and sufficient
within the same class.  We now prove this scoped optimality, identify the
compact conical carriers, and construct a noncompact comparison manifold in
which the limiting \(\mathbb R^4\) is an actual boundary fibre.  There is only
a continuous comparison map from this manifold to the compact conical
carrier: it set-theoretically collapses the new boundary, is a homeomorphism
away from that boundary, and is nonproper and nonquotient.

Throughout this subsubsection,
\[
 B_R^Z(z):=\{w\in Z:d_Z(z,w)\leq R\}
\]
denotes the closed radius-\(R\) ball; when the ambient space and center are
clear, superscripts and centers are suppressed.

\begin{theorem}[Optimality within the native connected torus-quotient class]
\label{wb:cusp-cone-optimal-exponents}
\StatementRecord{wb:cusp-cone-optimal-exponents}{proved}{Historical events: GPT-6 Astra (OpenAI), ChatGPT Pro mode, effort Pro, S6\_09\_PAPERS, 17 November 2026: targeted proof check. Also checked in this writing. Pre-continuation writing events: GPT-6 Astra (OpenAI), ChatGPT Pro mode, effort Pro, S6\_10\_PAPERS, 18 November 2026: targeted proof check. Continuation metadata events: GPT-6 Astra (OpenAI), ChatGPT Pro mode, effort Pro, S6\_10\_PAPERS continuation, 19 November 2026: used and rechecked in this run; status-vocabulary reconciliation; metadata reconciliation only; no new mathematical check. \emph{Pre-continuation recorded proof scope.} GPT-6 Astra (OpenAI, ChatGPT Pro mode), S6\_09\_PAPERS, effort Pro, 2026-11-17; also checked in this writing by GPT-6 Astra (OpenAI), ChatGPT Pro mode, effort Pro, 18 November 2026}{the hypotheses, notation and definitions stated here; the explicit constructions and calculations in the proof}

Consider only native connected fibrewise torus quotients of copies of the
original marked fibres, and equip every such quotient with the induced,
unscaled Euclidean quotient metric.

\begin{enumerate}
\item For a native connected torus quotient of \(r\) copies over the
punctured cusp, if the quotient systole tends to infinity on the positive
radial ray, then its lattice rank is at most \(2r\).  A rank-four quotient
therefore requires \(r\geq2\).  Equality is attained for \(r=2\) by
\[
 (\Lambda/K_\Lambda)^{\oplus2},
 \qquad K_\Lambda=
 \mathbb Z\langle\widehat w,\widehat\delta\rangle,
\]
whose physical period lattice is
\[
 L_T^{(12)}=\operatorname{diag}(Y_T,Y_T)\mathbb Z^4.
\]
The corresponding smash source is
\(X^{\wedge2}\cong S^{12}\).

\item For a native connected torus quotient of \(r\) copies over the full
regular base, or over a finite cover of that base, radial divergence of the
quotient systole forces its lattice rank to be at most \(r\).  A rank-four
quotient therefore requires \(r\geq4\).  Equality is attained by the four
copies of the full-monodromy-invariant character \(\chi_\gamma\); on the
radial fibre its period lattice is
\[
 L_T^{(24)}=\ell_T\mathbb Z^4,
 \qquad
 \ell_T=T+H-B+O(T^{-1}),
\]
and the corresponding smash source is
\(X^{\wedge4}\cong S^{24}\).
\end{enumerate}

Thus, among precisely these native connected torus quotients with their
induced unscaled marked Euclidean metrics, \(S^{12}\) is the smallest smash
sphere produced by the cusp-local rank-two quotient, and \(S^{24}\) is the
smallest smash sphere produced by a quotient which survives the full
regular-base monodromy.  No optimality is asserted here among arbitrary
continuous quotient maps or after a \(T\)-dependent rescaling of the metric.
\end{theorem}



\begin{proof}
\begin{enumerate}[label=\arabic*.,leftmargin=*]
\item For \(r\) copies, every element of
\[
 K_\Lambda^{\oplus r}
 =\mathbb Z\langle\widehat w,\widehat\delta\rangle^{\oplus r}
\]
has uniformly bounded physical length on the radial cusp.  Let
\(N\subset\Lambda^{\oplus r}\) be a primitive sublattice defining a
connected torus quotient.  If some \(v\in K_\Lambda^{\oplus r}\) had
nonzero class in \(\Lambda^{\oplus r}/N\), that nonzero quotient period
would have quotient norm at most the original norm of \(v\), uniformly in
\(T\).  The quotient systole could not diverge.  Hence
\[
 K_\Lambda^{\oplus r}\subset N,
\]
and
\[
 \operatorname{rank}(\Lambda^{\oplus r}/N)
 \leq 4r-2r=2r.
\]
For \(r=2\), the primitive kernel \(K_\Lambda^{\oplus2}\) attains this
bound.  The cusp monodromy is the identity modulo \(K_\Lambda\), so the
quotient descends over the punctured cusp.  Its radial lattice is exactly
\(\operatorname{diag}(Y_T,Y_T)\mathbb Z^4\), and
\[
 \operatorname{sys}(L_T^{(12)})
 =\operatorname{sys}(Y_T\mathbb Z^2)\geq T-C\longrightarrow\infty.
\]

\item For a quotient over the full regular base, monodromy invariance gives a
stronger bound.  Since \(\widehat w\in N\) and
\[
 A_1\widehat w=\widehat u-\widehat w,
\]
one has \(\widehat u\in N\).  Together with
\(\widehat w,\widehat\delta\in N\), this places the rank-three lattice
\[
 \mathbb Z\langle\widehat u,\widehat w,
 \widehat\delta\rangle^{\oplus r}
\]
inside \(N\), and hence
\[
 \operatorname{rank}(\Lambda^{\oplus r}/N)\leq r.
\]
The same conclusion holds after a finite base cover.  Indeed, if
\(\Gamma_{\mathrm{mon}}\) is the corresponding finite-index monodromy
subgroup and \(G_*=A_2A_1=I+N_*\), then \(N_*^2=0\), and
\(\Gamma_{\mathrm{mon}}\) contains \(G_*^m\) for
some \(m>0\).  The exact orbit formula
\[
 G_*^m\widehat w
 =\widehat w+m(\widehat u-\widehat\delta)
\]
implies \(m\widehat u\in N\); primitivity of \(N\) implies
\(\widehat u\in N\).

\item The invariant character \(\chi_\gamma\) has precisely the complementary
rank-one quotient in each copy.  The full-base construction, including the
two finite Seifert fillings and the cone at the unipotent cusp, proves that
four copies attain rank four.  Finally, after choosing the based
diffeomorphism \(X\cong_{\mathrm{diff}}S^6\) already reconstructed in this
workbench,
\(X^{\wedge r}\cong(S^6)^{\wedge r}\cong S^{6r}\), which gives the two
claimed sphere dimensions.
\end{enumerate}
\end{proof}

\begin{theorem}[Exact sphere-to-sphere carriers with moving four-spheres]
\label{wb:cusp-cone-angle-forgotten-sphere-carriers}
\StatementRecord{wb:cusp-cone-angle-forgotten-sphere-carriers}{proved}{Historical events: GPT-6 Astra (OpenAI), ChatGPT Pro mode, effort Pro, S6\_09\_PAPERS, 17 November 2026: targeted proof check. Also checked in this writing. Pre-continuation writing events: GPT-6 Astra (OpenAI), ChatGPT Pro mode, effort Pro, S6\_10\_PAPERS, 18 November 2026: targeted proof check. Continuation metadata events: GPT-6 Astra (OpenAI), ChatGPT Pro mode, effort Pro, S6\_10\_PAPERS continuation, 19 November 2026: used and rechecked in this run; status-vocabulary reconciliation; metadata reconciliation only; no new mathematical check. \emph{Pre-continuation recorded proof scope.} GPT-6 Astra (OpenAI, ChatGPT Pro mode), S6\_09\_PAPERS, effort Pro, 2026-11-17; also checked in this writing by GPT-6 Astra (OpenAI), ChatGPT Pro mode, effort Pro, 18 November 2026}{the hypotheses, notation and definitions stated here; the explicit constructions and calculations in the proof}

There are continuous surjective based quotient maps
\[
 \widehat q_3:(X,x_0)\longrightarrow(S^3,*),
 \qquad
 \widehat q_\gamma:(X,x_0)\longrightarrow(S^2,*),
\]
whose smash powers are continuous surjective quotient maps
\begin{equation}
 \widehat q_3^{\wedge2}:S^{12}\twoheadrightarrow S^6,
 \qquad
 \widehat q_\gamma^{\wedge4}:S^{24}\twoheadrightarrow S^8.
 \label{wb:cusp-cone-sphere-to-sphere-maps}
\end{equation}
Use the marked presentation \(\overline T^2=\mathbb R^2/\mathbb Z^2\) and
put
\[
 K^{(1)}
 =((\mathbb R/\mathbb Z)\times\{0\})
  \cup(\{0\}\times(\mathbb R/\mathbb Z)),
 \qquad
 c:\overline T^2\longrightarrow\overline T^2/K^{(1)}\cong S^2.
\]
Here \(c\) is the quotient collapsing the coordinate one-skeleton.  For
every sufficiently large \(T\), put
\begin{align*}
 r_T&=e^{-2\pi T},\\
 K_T^{(1)}&:=\iota_T(K^{(1)})\subset Q_T,\\
 S_T^2&:=\overline T^2/K^{(1)}\cong Q_T/K_T^{(1)},\\
 c_T&:=c\circ\iota_T^{-1}:Q_T\longrightarrow S_T^2,
\end{align*}
where the identification of the two displayed quotients is induced by
\(\iota_T\).  Within this theorem and its proof, \(c_T\) denotes the
displayed collapse map; the earlier scalar \(c_T=\xi_T-\eta_T\) does not
enter.
At reduced-suspension cross-section height \(r_T\), the first target contains
a based topologically embedded four-sphere, equipped with its intrinsic
quotient length metric,
\[
 \mathbb S_{12,T}^4=S_T^2\wedge S_T^2\hookrightarrow S^6,
\]
and the second target contains, at the same cross-section height, a based
topologically embedded four-sphere, equipped with its intrinsic quotient
length metric,
\[
 \mathbb S_{24,T}^4=C_T^{\wedge4}\hookrightarrow S^8.
\]
On synchronized radial fibres, the exact restrictions factor as
\begin{align}
 F_T^2&\xrightarrow{\ \pi_T\times\pi_T\ }
 Q_T^2\xrightarrow{\ c_T\times c_T\ }(S_T^2)^2
 \xrightarrow{\ \rho_{12,T}\ }S_T^2\wedge S_T^2
 =\mathbb S_{12,T}^4,
 \label{wb:cusp-cone-S12-S4-radial-map}\\
 F_T^4&\xrightarrow{\ \chi_{\gamma,T}^{\times4}\ }
 C_T^4\xrightarrow{\ \rho_{24,T}\ }C_T^{\wedge4}
 =\mathbb S_{24,T}^4,
 \label{wb:cusp-cone-S24-S4-radial-map}
\end{align}
Here \(c_T\) collapses \(K_T^{(1)}\), and \(\rho_{12,T}\) and
\(\rho_{24,T}\) are the smash quotient maps.  Thus the
full collapsed locus in \(Q_T^2\) is exactly
\begin{equation}
 W_{12,T}
 =\bigl(K_T^{(1)}\times Q_T\bigr)
  \cup\bigl(Q_T\times K_T^{(1)}\bigr),
 \label{wb:cusp-cone-S12-exact-collapsed-locus}
\end{equation}
and the full collapsed locus in \(C_T^4\) is exactly
\begin{equation}
 W_{24,T}
 =\{(x_1,x_2,x_3,x_4)\in C_T^4:
       x_j=0\pmod{\ell_T}\ \text{for some }j\}.
 \label{wb:cusp-cone-S24-exact-collapsed-locus}
\end{equation}
The maps
\(\rho_{12,T}\circ(c_T\times c_T):Q_T^2\to
S_T^2\wedge S_T^2\) and
\(\rho_{24,T}:C_T^4\to C_T^{\wedge4}\) send exactly \(W_{12,T}\) and
\(W_{24,T}\), respectively, to the smash basepoint and are injective on
their complements.  The full composites from \(F_T^2\) and \(F_T^4\)
still identify every pair of points identified by
\(\pi_T\times\pi_T\) and \(\chi_{\gamma,T}^{\times4}\), respectively; in
particular, their first arrows retain their nontrivial torus kernels.
Their exact smash-basepoint preimages in the full radial sources are
\[
 \bigl(\pi_T^{-1}(K_T^{(1)})\times F_T\bigr)
 \cup\bigl(F_T\times\pi_T^{-1}(K_T^{(1)})\bigr)
\]
and
\[
 \{(x_1,x_2,x_3,x_4)\in F_T^4:
   \chi_{\gamma,T}(x_j)=0\ \text{in }C_T\ \text{for some }j\},
\]
respectively.  Outside these basepoint preimages, two source points have the
same final image if and only if their images under the corresponding first
arrow in \cref{wb:cusp-cone-S12-S4-radial-map} and \cref{wb:cusp-cone-S24-S4-radial-map}
are equal.

Choose the Euclidean Riemannian \(\ell^2\)-product length metrics on
\(Q_T^2\) and \(C_T^4\) before taking the smash quotients, and use the
successive quotient length metrics specified in the proof.  With the
specified midpoint basepoints,
\begin{equation}
 (\mathbb S_{12,T}^4,\zeta_{12,T})
 \longrightarrow(\mathbb R^4,0),
 \qquad
 (\mathbb S_{24,T}^4,\zeta_{24,T})
 \longrightarrow(\mathbb R^4,0)
 \label{wb:cusp-cone-moving-S4-R4-limits}
\end{equation}
in pointed local-isometric and pointed Gromov--Hausdorff topology.  Moreover,
the convergence is smooth on the eventually Euclidean basepoint
neighborhoods.  For every \(R>0\), once \(T-C>4R\), respectively once
\(\ell_T>4R\), the corresponding quotient map restricts to an isometry from
the radius-\(R\) ball in \(Q_T^2\), respectively \(C_T^4\), \emph{onto} the
entire radius-\(R\) ball in the indicated smash sphere.  The domain ball, and
hence the quotient ball, is exactly a Euclidean four-ball.
\end{theorem}



\begin{proof}
\begin{enumerate}[label=\arabic*.,leftmargin=*]
\item Write \(\Sigma_r(Z,z_0)\) for the reduced suspension
\[
 \frac{[0,\epsilon]\times Z}
 {(\{0,\epsilon\}\times Z)\cup([0,\epsilon]\times\{z_0\})},
\]
where the displayed union is collapsed to one basepoint.

\item Use the marked presentation
\(\overline T^2=\mathbb R^2/\mathbb Z^2\), let
\[
 K^{(1)}=
 ((\mathbb R/\mathbb Z)\times\{0\})
 \cup(\{0\}\times(\mathbb R/\mathbb Z)),
\]
and let
\[
 c:\overline T^2\longrightarrow
 \overline T^2/K^{(1)}\cong S^2
\]
be the degree-one quotient obtained by collapsing the one-skeleton.  The
formula
\[
 \widehat p_2:\mathcal A\longrightarrow
 \Sigma_r(\overline T^2/K^{(1)},c(K^{(1)}))\cong S^3,
 \qquad
 [r,e^{i\phi},u]\longmapsto[r,c(u)]
\]
is well defined: the two suspension ends of \(\mathcal A\) and the entire
cylinder over \(K^{(1)}\) map to the reduced-suspension basepoint.  It is a
continuous surjection.  Set
\(\widehat q_3=\widehat p_2\circ q_2\).

\item For every sufficiently large \(T\), define the transported one-skeleton and
its collapse by
\begin{equation}
 \begin{aligned}
 K_T^{(1)}&:=\iota_T(K^{(1)})\subset Q_T,\\
 S_T^2&:=\overline T^2/K^{(1)}\cong Q_T/K_T^{(1)},\\
 c_T&:=c\circ\iota_T^{-1}:Q_T\longrightarrow S_T^2.
 \end{aligned}
 \label{wb:cusp-cone-transported-skeleton-collapse}
\end{equation}
The displayed target identification is the one induced by \(\iota_T\), so
\(c_T^{-1}(*)=K_T^{(1)}\), where
\(*=c(K^{(1)})=c_T(K_T^{(1)})\).

\item Give \(Q_T\) its Euclidean quotient length metric and \(S_T^2\) the quotient
length metric under \(c_T\).  Put
\[
 y_*=[(1/2,1/2)]\in\overline T^2,
 \qquad y_T:=\iota_T(y_*)\in Q_T,
 \qquad z_T:=c_T(y_T)\in S_T^2.
\]
Equivalently, the metric on \(\overline T^2\) used below is
\(g_T:=\iota_T^*g_{Q_T}\).  For every \(v\in\mathbb R^2\), the exact
decomposition
\[
 Y_T=T\begin{pmatrix}0&1\\-1&0\end{pmatrix}+E_T,
 \qquad\|E_T\|_{\mathrm{op}}\leq C,
\]
gives
\[
 \|Y_Tv\|\geq(T-C)\|v\|.
\]
The standard flat distance from \(y_*\) to \(K^{(1)}\) is \(1/2\);
therefore
\begin{equation}
 d_T:=\operatorname{dist}_{Q_T}(y_T,K_T^{(1)})
 =\operatorname{dist}_{g_T}(y_*,K^{(1)})
 \geq\frac{T-C}{2}.
 \label{wb:cusp-cone-distance-to-skeleton}
\end{equation}
For a length-space quotient which collapses a closed subset \(K\) to one
point, the distance between two noncollapsed classes represented by \(x,y\)
is
\[
 \min\{d(x,y),d(x,K)+d(y,K)\}.
\]
The two displayed terms are upper bounds, obtained respectively from a
direct path and from paths from \(x\) and \(y\) to \(K\).  Conversely, a
rectifiable quotient path either avoids the collapsed point, in which case
its compact image lifts to \(X\setminus K\) and has length at least
\(d(x,y)\), or meets that point, in which case its two pieces have total
length at least \(d(x,K)+d(y,K)\).  Taking infima proves the formula.
If \(T-C>4R\), then \(d_T>2R\).  The displayed formula and the systolic
bound show that
\[
 c_T:B_R^{Q_T}(y_T)\longrightarrow B_R^{S_T^2}(z_T)
\]
is an isometry onto the entire target ball, and both balls are exactly
Euclidean two-balls.  Indeed, the second term in the minimum cannot realize
a distance at most \(2R\), and any point of the target ball has a unique
noncollapsed representative; the same formula puts that representative in
the source ball.

\item At reduced-suspension coordinate \(r_T=e^{-2\pi T}\), \(S_T^2\) is a based
embedded subspace of \(S^3\): its point \(*\) is the global
reduced-suspension basepoint.  Hence
\[
 \mathbb S_{12,T}^4=S_T^2\wedge S_T^2
\]
embeds in \(S^3\wedge S^3=S^6\).  Equip \(Q_T^2\) first with the Euclidean
Riemannian \(\ell^2\)-product length metric, equip \((S_T^2)^2\) with the
\(\ell^2\)-product of the two quotient length metrics, and let
\[
 \rho_{12,T}:(S_T^2)^2\longrightarrow S_T^2\wedge S_T^2
\]
carry the quotient length metric obtained by collapsing the fat wedge
\[
 \mathcal W_{12,T}
 =(S_T^2\times\{*\})\cup(\{*\}\times S_T^2).
\]
Since \(c_T^{-1}(*)=K_T^{(1)}\), the inverse image of
\(\mathcal W_{12,T}\) under \(c_T\times c_T\) is exactly \(W_{12,T}\) in
\cref{wb:cusp-cone-S12-exact-collapsed-locus}.  On
\(Q_T^2\setminus W_{12,T}\), the map \(c_T\times c_T\) is injective and its
image avoids \(\mathcal W_{12,T}\); the map \(\rho_{12,T}\) is injective on
the complement of \(\mathcal W_{12,T}\).  Hence their composite identifies
all of \(W_{12,T}\) to the smash basepoint and makes no other
identifications.

\item The point \(\zeta_{12,T}=[z_T,z_T]\) has distance \(d_T\) from this fat
wedge.  The one-factor onto isometry above gives an onto isometry from the
radius-\(R\) ball about \((y_T,y_T)\) in \(Q_T^2\) to the corresponding
ball about \((z_T,z_T)\) in \((S_T^2)^2\).  Applying the same collapsed-set
distance formula to \(\rho_{12,T}\), and using \(d_T>2R\), shows that
\begin{equation}
 \rho_{12,T}\circ(c_T\times c_T):
 B_R^{Q_T^2}((y_T,y_T))
 \xrightarrow{\ \cong_{\mathrm{isom}}\ }
 B_R^{\mathbb S_{12,T}^4}(\zeta_{12,T})
 \label{wb:cusp-cone-S12-onto-exact-ball}
\end{equation}
is onto as well as isometrically embedded.  The systolic bound
\(\operatorname{sys}(Q_T^2)\geq T-C>4R\) makes its domain exactly a
Euclidean four-ball.  On a radial fibre,
\[
 (\widehat p_2\circ q_2)|_{F_T}
 =[r_T,c(\iota_T^{-1}(\pi_T(\,\cdot\,)))]
 =[r_T,c_T(\pi_T(\,\cdot\,))],
\]
and taking two factors followed by \(\rho_{12,T}\) proves the entire
factorization in \cref{wb:cusp-cone-S12-S4-radial-map}.

\item For the four-copy construction, choose \(1\in S^1\) as character basepoint
and define
\[
 \widehat p_1:\mathcal B\longrightarrow\Sigma_r(S^1,1)\cong S^2,
 \qquad
 [r,e^{i\phi},\theta]\longmapsto[r,\theta].
\]
The formula forgets the base angle and collapses the two suspension ends
and the entire cylinder with \(\theta=1\) to the reduced-suspension
basepoint.  It is a continuous surjection.  Put
\(\widehat q_\gamma=\widehat p_1\circ q_{1,\mathrm{cusp}}\).

\item At the interior suspension coordinate \(r_T=e^{-2\pi T}\), identify the
character circle \(C_T=\mathbb R/\ell_T\mathbb Z\) with
\(\mathbb R/\mathbb Z\subset S^2\) by \([x]\mapsto[x/\ell_T]\).  It is a
based embedded circle, so its fourth smash is an embedded
\[
 \mathbb S_{24,T}^4=C_T^{\wedge4}\subset(S^2)^{\wedge4}=S^8.
\]
Equip \(C_T^4\) first with the Euclidean Riemannian \(\ell^2\)-product
length metric, and let
\[
 \rho_{24,T}:C_T^4\longrightarrow C_T^{\wedge4}
\]
carry the quotient length metric obtained by collapsing exactly the fat
wedge \(W_{24,T}\) of
\cref{wb:cusp-cone-S24-exact-collapsed-locus}.  Away from that fat wedge the
smash quotient is injective, so no further points are identified.  The
midpoint
\[
 \zeta_{24,T}
 =[(\ell_T/2,\ell_T/2,\ell_T/2,\ell_T/2)]
\]
has distance \(\ell_T/2\) from that fat wedge.  If \(\ell_T>4R\), neither a
nonzero period nor a path through the collapsed fat wedge can shorten a
distance between points in the radius-\(R\) ball of \(C_T^4\).  Conversely,
the collapsed-set distance formula shows that every point of the
radius-\(R\) quotient ball has its unique noncollapsed representative in
that product ball.  Consequently
\begin{equation}
 \rho_{24,T}:
 B_R^{C_T^4}((\ell_T/2,\ell_T/2,\ell_T/2,\ell_T/2))
 \xrightarrow{\ \cong_{\mathrm{isom}}\ }
 B_R^{\mathbb S_{24,T}^4}(\zeta_{24,T})
 \label{wb:cusp-cone-S24-onto-exact-ball}
\end{equation}
is onto as well as isometrically embedded, and its domain is exactly a
Euclidean four-ball because \(\operatorname{sys}(C_T^4)=\ell_T>4R\).  On
the radial source the formula discards only the base angle and then performs
this smash quotient, giving
\cref{wb:cusp-cone-S24-S4-radial-map}.

\item Both one-copy maps are continuous surjections from compact \(X\) to
Hausdorff spheres, hence quotient maps, and they are based.  Their smash
powers are surjective; compact-to-Hausdorff again makes them quotient maps.
Using the same chosen based diffeomorphism \(X\cong_{\mathrm{diff}}S^6\),
the homeomorphisms
\[
 X^{\wedge2}\cong S^{12},\quad(S^3)^{\wedge2}\cong S^6,
 \qquad
 X^{\wedge4}\cong S^{24},\quad(S^2)^{\wedge4}\cong S^8
\]
prove \cref{wb:cusp-cone-sphere-to-sphere-maps}.  The exact Euclidean-ball
statements prove all three convergences in
\cref{wb:cusp-cone-moving-S4-R4-limits}.
\end{enumerate}
\end{proof}

\begin{corollary}[The two-copy rank-two sphere carrier is null-homotopic]
\label{wb:cusp-cone-S12-sphere-map-null}
\StatementRecord{wb:cusp-cone-S12-sphere-map-null}{proved}{Historical events: GPT-6 Astra (OpenAI), ChatGPT Pro mode, effort Pro, S6\_09\_PAPERS, 17 November 2026: targeted proof check. Also checked in this writing. Pre-continuation writing events: GPT-6 Astra (OpenAI), ChatGPT Pro mode, effort Pro, S6\_10\_PAPERS, 18 November 2026: targeted proof check. Continuation metadata events: GPT-6 Astra (OpenAI), ChatGPT Pro mode, effort Pro, S6\_10\_PAPERS continuation, 19 November 2026: used and rechecked in this run; status-vocabulary reconciliation; metadata reconciliation only; no new mathematical check. \emph{Pre-continuation recorded proof scope.} GPT-6 Astra (OpenAI, ChatGPT Pro mode), S6\_09\_PAPERS, effort Pro, 2026-11-17; also checked in this writing by GPT-6 Astra (OpenAI), ChatGPT Pro mode, effort Pro, 18 November 2026}{the hypotheses, notation and definitions stated here; \cref{wb:cusp-cone-sphere-to-sphere-maps}; \cref{wb:cusp-s12-full-carrier-null}}

After the based identifications in
\cref{wb:cusp-cone-sphere-to-sphere-maps}, the continuous surjective quotient
\[
 \widehat q_3^{\wedge2}:S^{12}\twoheadrightarrow S^6
\]
is based null-homotopic.  In fact, \(g\wedge g:S^{12}\to S^6\) is based
null-homotopic for every based map \(g:S^6\to S^3\).
\end{corollary}



\begin{proof}
\begin{enumerate}[label=\arabic*.,leftmargin=*]
\item Apply the full external-smash calculation in
\cref{wb:cusp-s12-full-carrier-null} to
\(g=\widehat q_3\).  That theorem proves the stronger assertion that
\(x\wedge y:S^{12}\to S^6\) is null-homotopic for every pair of based
maps \(x,y:S^6\to S^3\).
\end{enumerate}
\end{proof}

\begin{corollary}[Every repeated smash of the circle-quotient carrier is null]
\label{wb:cusp-cone-S24-sphere-map-null}
\StatementRecord{wb:cusp-cone-S24-sphere-map-null}{proved}{Historical events: GPT-6 Astra (OpenAI), ChatGPT Pro mode, effort Pro, S6\_09\_PAPERS, 17 November 2026: targeted proof check. Also checked in this writing. Pre-continuation writing events: GPT-6 Astra (OpenAI), ChatGPT Pro mode, effort Pro, S6\_10\_PAPERS, 18 November 2026: targeted proof check. Continuation metadata events: GPT-6 Astra (OpenAI), ChatGPT Pro mode, effort Pro, S6\_10\_PAPERS continuation, 19 November 2026: used and rechecked in this run; status-vocabulary reconciliation; metadata reconciliation only; no new mathematical check. \emph{Pre-continuation recorded proof scope.} GPT-6 Astra (OpenAI, ChatGPT Pro mode), S6\_09\_PAPERS, effort Pro, 2026-11-17; also checked in this writing by GPT-6 Astra (OpenAI), ChatGPT Pro mode, effort Pro, 18 November 2026}{the hypotheses, notation and definitions stated here; \cite[Corollary~4.24, p.~360]{Hatcher2002}; \cite[Table~1, p.~113]{IsaksenWangXu2023}; \cref{wb:cusp-cone-S24-smash-square-factorization}}

Put \(f=\widehat q_\gamma:S^6\to S^2\) after the chosen based
identification \(X\cong S^6\).  For every \(r\geq2\),
\begin{equation}
 \bigl[f^{\wedge r}\bigr]=0
 \quad\text{in}\quad
 \pi_{6r}(S^{2r}).
 \label{wb:cusp-cone-S24-null-maps}
\end{equation}
In particular,
\begin{equation}
 \bigl[\widehat q_\gamma^{\wedge2}\bigr]
 =0\in\pi_{12}(S^4),
 \qquad
 \bigl[\widehat q_\gamma^{\wedge4}\bigr]
 =0\in\pi_{24}(S^8).
 \label{wb:cusp-cone-S24-null-special-cases}
\end{equation}
Thus the displayed surjective quotient
\(\widehat q_\gamma^{\wedge4}:S^{24}\twoheadrightarrow S^8\) is globally
null-homotopic.  This does not alter either its point-set surjectivity or the
exact intrinsic \(\mathbb R^4\) limit on the moving four-spheres
\(\mathbb S_{24,T}^4\): those are properties of the specified representative
and its radial restrictions, whereas null-homotopy classifies that
representative only up to based homotopy.
\end{corollary}



\begin{proof}
\begin{enumerate}[label=\arabic*.,leftmargin=*]
\item Associativity of the smash product gives the
exact factorization
\begin{equation}
 \begin{tikzcd}[column sep=large]
 S^6\wedge S^6
 \arrow[r,"f\wedge\operatorname{id}_{S^6}"]&
 S^2\wedge S^6
 \arrow[r,"\operatorname{id}_{S^2}\wedge f"]&
 S^2\wedge S^2 .
 \end{tikzcd}
 \label{wb:cusp-cone-S24-smash-square-factorization}
\end{equation}
Under the standard based sphere identifications, the first arrow is the
sixfold suspension of \(f\), up to the harmless twist homeomorphisms dictated
by the left-versus-right suspension convention:
\begin{equation}
 f\wedge\operatorname{id}_{S^6}
 =\Sigma^6f:S^{12}\longrightarrow S^8.
 \label{wb:cusp-cone-S24-sixfold-suspension}
\end{equation}
Freudenthal's suspension theorem is already in its isomorphism range here,
because \(12\leq 2\cdot8-2\), and remains in that range after every further
suspension.  It therefore identifies \(\pi_{12}(S^8)\) with the fourth
stable stem.  The low-stem calculation gives
\phantomsection\label{qa:III:layout:11}
\begin{equation}
 \pi_{12}(S^8)\cong\pi_4^{\mathrm S}=0.
 \label{wb:cusp-cone-fourth-stable-stem}
\end{equation}
Here we use \cite[Corollary~4.24, p.~360]{Hatcher2002} and
\cite[Table~1, p.~113]{IsaksenWangXu2023}.
Hence \(\Sigma^6f\) is based null-homotopic, so the composite in
\cref{wb:cusp-cone-S24-smash-square-factorization}, namely
\(f\wedge f\), is based null-homotopic.  For \(r\geq2\), the canonical
associativity homeomorphisms identify
\begin{equation}
 f^{\wedge r}=(f\wedge f)\wedge f^{\wedge(r-2)},
\end{equation}
with the convention \(f^{\wedge0}=\operatorname{id}_{S^0}\).  Smashing a
based null-homotopy with any based map is a based null-homotopy.  This proves
\cref{wb:cusp-cone-S24-null-maps} and \cref{wb:cusp-cone-S24-null-special-cases}.
\end{enumerate}
\end{proof}

\begin{corollary}[Literal \(\mathbb R^4\) open cells in the sphere targets]
\label{wb:cusp-cone-literal-R4-open-cells}
\StatementRecord{wb:cusp-cone-literal-R4-open-cells}{proved}{Historical events: GPT-6 Astra (OpenAI), ChatGPT Pro mode, effort Pro, S6\_09\_PAPERS, 17 November 2026: targeted proof check. Also checked in this writing. Pre-continuation writing events: GPT-6 Astra (OpenAI), ChatGPT Pro mode, effort Pro, S6\_10\_PAPERS, 18 November 2026: targeted proof check. Continuation metadata events: GPT-6 Astra (OpenAI), ChatGPT Pro mode, effort Pro, S6\_10\_PAPERS continuation, 19 November 2026: used and rechecked in this run; status-vocabulary reconciliation; metadata reconciliation only; no new mathematical check. \emph{Pre-continuation recorded proof scope.} GPT-6 Astra (OpenAI, ChatGPT Pro mode), S6\_09\_PAPERS, effort Pro, 2026-11-17; also checked in this writing by GPT-6 Astra (OpenAI), ChatGPT Pro mode, effort Pro, 18 November 2026}{the hypotheses, notation and definitions stated here; \cref{wb:cusp-cone-S12-onto-exact-ball}; \cref{wb:cusp-cone-S24-onto-exact-ball}}

Let \(*\) denote the smash basepoint and set
\[
 U_{12,T}:=\mathbb S_{12,T}^4\setminus\{*\},
 \qquad
 U_{24,T}:=\mathbb S_{24,T}^4\setminus\{*\}.
\]
The embedded moving four-spheres restrict to topological embeddings
\[
 U_{12,T}\cong\mathbb R^4\hookrightarrow S^6,
 \qquad
 U_{24,T}\cong\mathbb R^4\hookrightarrow S^8.
\]
More exactly, the quotient-stage maps restrict to homeomorphisms
\begin{align}
 (Q_T\setminus K_T^{(1)})^2
 &\xrightarrow{\ \cong\ }U_{12,T},
 \label{wb:cusp-cone-U12-homeomorphism}\\
 (C_T\setminus\{0\})^4
 &\xrightarrow{\ \cong\ }U_{24,T}.
 \label{wb:cusp-cone-U24-homeomorphism}
\end{align}
These copies of \(\mathbb R^4\) are nonclosed topological open cells: they
are open in their moving four-spheres, not open in the ambient \(S^6\) or
\(S^8\), and their closures in those ambient spheres are
\(\mathbb S_{12,T}^4\) and \(\mathbb S_{24,T}^4\), respectively.  They are
not globally isometric to Euclidean \(\mathbb R^4\) for either the assigned
compact quotient metrics or the restriction of any fixed compatible metric
on the compact sphere targets.  The exact midpoint balls in
\cref{wb:cusp-cone-S12-onto-exact-ball} and \cref{wb:cusp-cone-S24-onto-exact-ball} give the
precise metric statement: as \(T\to\infty\), every fixed Euclidean
radius-\(R\) ball occurs isometrically and onto the corresponding pointed
ball.
\end{corollary}



\begin{proof}
\begin{enumerate}[label=\arabic*.,leftmargin=*]
\item For a quotient which collapses a closed subset \(K\) to a point, restriction
to the complement gives a homeomorphism
\[
 X\setminus K\xrightarrow{\ \cong\ }(X/K)\setminus\{*\}.
\]
Indeed, the complement of the quotient point is open, its inverse image is
the saturated open set \(X\setminus K\), and restriction of a quotient map
over an open subset of its target is again a quotient map.  The restricted
map is bijective, and a bijective quotient map is a homeomorphism.
Consequently
\[
 Q_T\setminus K_T^{(1)}
 \cong S_T^2\setminus\{*\}.
\]
The marked coordinate one-skeleton \(K^{(1)}\) is the \(1\)-skeleton of the
standard one-vertex CW structure on \(\overline T^2\), with one \(0\)-cell,
two \(1\)-cells, and one \(2\)-cell.  Its complement is therefore the open
\(2\)-cell \((0,1)^2\cong\mathbb R^2\).  Transport by
\(\iota_T\) gives
\[
 Q_T\setminus K_T^{(1)}\cong\mathbb R^2.
\]
Likewise \(C_T\setminus\{0\}\) is an open interval and hence is homeomorphic
to \(\mathbb R\).

\item The complement of the smash basepoint is the product of the complements of
the factor basepoints.  Therefore
\[
 (S_T^2\wedge S_T^2)\setminus\{*\}
 \cong(S_T^2\setminus\{*\})^2\cong\mathbb R^4,
\]
and
\[
 C_T^{\wedge4}\setminus\{*\}
 \cong(C_T\setminus\{0\})^4\cong\mathbb R^4.
\]
The full moving four-spheres are topologically embedded in their compact
sphere targets, so these restrictions are topological embeddings.  In each
case the omitted smash basepoint lies in the closure of the open cell, which
proves the closure assertions. These cells are not ambient open: an ambient open set in $S^6$ or $S^8$ has local integral homology concentrated in degree $6$ or $8$, respectively, whereas the displayed homeomorphisms from $\mathbb R^4$ give local homology concentrated in degree $4$. Finally, every compatible metric on a
compact target is bounded, as is each compact quotient metric assigned to a
moving four-sphere, whereas Euclidean \(\mathbb R^4\) is unbounded.  Thus no
global isometry is possible; the exact metric content is precisely the
onto-ball calculation already proved.
\end{enumerate}
\end{proof}


\paragraph{The moving four-spheres are null-homotopic in the targets}
\label{wb:cusp-cone-moving-S4-null-homotopy}
\StatementRecord{wb:cusp-cone-moving-S4-null-homotopy}{proved}{Historical events: GPT-6 Astra (OpenAI), ChatGPT Pro mode, effort Pro, S6\_09\_PAPERS, 17 November 2026: targeted proof check. Also checked in this writing. Pre-continuation writing events: GPT-6 Astra (OpenAI), ChatGPT Pro mode, effort Pro, S6\_10\_PAPERS, 18 November 2026: targeted proof check. Continuation metadata events: GPT-6 Astra (OpenAI), ChatGPT Pro mode, effort Pro, S6\_10\_PAPERS continuation, 19 November 2026: used and rechecked in this run; status-vocabulary reconciliation; metadata reconciliation only; no new mathematical check. \emph{Pre-continuation recorded proof scope.} GPT-6 Astra (OpenAI, ChatGPT Pro mode), S6\_09\_PAPERS, effort Pro, 2026-11-17; also checked in this writing by GPT-6 Astra (OpenAI), ChatGPT Pro mode, effort Pro, 18 November 2026}{the hypotheses, notation and definitions stated here; the argument and qualifications immediately following}

For a fixed interior reduced-suspension coordinate, the inclusion
\(S_T^2\hookrightarrow S^3\) contracts through based maps by sliding that
coordinate to a suspension endpoint; throughout the homotopy the collapsed
meridian remains at the global basepoint.  Hence
\(S_T^2\wedge S_T^2\hookrightarrow S^6\) is null-homotopic.  The identical
sliding argument contracts \(C_T\hookrightarrow S^2\), and therefore also
contracts \(C_T^{\wedge4}\hookrightarrow S^8\).  Thus the moving metric
four-spheres do not contradict the fixed-target orientation obstruction:
their fundamental classes die in the ambient targets.






\paragraph{Homotopical status of the two sphere-to-sphere maps}
\label{wb:cusp-cone-sphere-map-homotopy-status}
\StatementRecord{wb:cusp-cone-sphere-map-homotopy-status}{proved}{Historical events: GPT-6 Astra (OpenAI), ChatGPT Pro mode, effort Pro, S6\_09\_PAPERS, 17 November 2026: targeted proof check. Also checked in this writing. Pre-continuation writing events: GPT-6 Astra (OpenAI), ChatGPT Pro mode, effort Pro, S6\_10\_PAPERS, 18 November 2026: targeted proof check. Continuation metadata events: GPT-6 Astra (OpenAI), ChatGPT Pro mode, effort Pro, S6\_10\_PAPERS continuation, 19 November 2026: used and rechecked in this run; status-vocabulary reconciliation; metadata reconciliation only; no new mathematical check. \emph{Pre-continuation recorded proof scope.} GPT-6 Astra (OpenAI, ChatGPT Pro mode), S6\_09\_PAPERS, effort Pro, 2026-11-17; also checked in this writing by GPT-6 Astra (OpenAI), ChatGPT Pro mode, effort Pro, 18 November 2026}{the hypotheses, notation and definitions stated here; \cref{wb:cusp-cone-S24-sphere-map-null}; \cref{wb:cusp-cone-S12-sphere-map-null}; \cref{wb:cusp-null:thm}; \cref{mapsIII:angle-forgotten-one-copy-null}}

The maps \(S^{12}\to S^6\) and \(S^{24}\to S^8\) induce zero on all
positive-degree integral homology and cohomology groups and have zero
ordinary Hurewicz image.  They have no classical degree because source and
target dimensions differ.  The full four-copy map
\(\widehat q_\gamma^{\wedge4}:S^{24}\to S^8\) is based null-homotopic by
\cref{wb:cusp-cone-S24-sphere-map-null}.  The distinct two-copy map
\(\widehat q_3^{\wedge2}:S^{12}\to S^6\) is also based null-homotopic, by
\cref{wb:cusp-cone-S12-sphere-map-null}.  In both cases the vanishing is a
theorem about the full specified sphere-valued map, not merely its Hurewicz
image or the inclusion of a moving four-sphere. % checked carrier correction: one-copy-status
The one-copy class
\([\widehat q_3]\in\pi_6(S^3)\) is zero by the complete framed-torus
proof in \cref{wb:cusp-null:thm}.  Both this class and
\([\widehat q_\gamma]\in\pi_6(S^2)\) also vanish by the extension
homotopy in \cref{mapsIII:angle-forgotten-one-copy-null}.  That proof
uses the original formulas after forgetting the base angle and fixes
the chosen point of the central fibre throughout.  It supplies based
nullhomotopies of all their positive smash powers.  The retained-angle
map has a different target and is outside that assertion.





\begin{proposition}[The compact cone-on-\(T^4\) carriers]
\label{wb:cusp-cone-explicit-carriers}
\StatementRecord{wb:cusp-cone-explicit-carriers}{proved}{Historical events: GPT-6 Astra (OpenAI), ChatGPT Pro mode, effort Pro, S6\_09\_PAPERS, 17 November 2026: targeted proof check. Also checked in this writing. Pre-continuation writing events: GPT-6 Astra (OpenAI), ChatGPT Pro mode, effort Pro, S6\_10\_PAPERS, 18 November 2026: targeted proof check. Continuation metadata events: GPT-6 Astra (OpenAI), ChatGPT Pro mode, effort Pro, S6\_10\_PAPERS continuation, 19 November 2026: used and rechecked in this run; status-vocabulary reconciliation; metadata reconciliation only; no new mathematical check. \emph{Pre-continuation recorded proof scope.} GPT-6 Astra (OpenAI, ChatGPT Pro mode), S6\_09\_PAPERS, effort Pro, 2026-11-17; also checked in this writing by GPT-6 Astra (OpenAI), ChatGPT Pro mode, effort Pro, 18 November 2026}{the hypotheses, notation and definitions stated here; the explicit constructions and calculations in the proof}

Fix \(T_0\) sufficiently large and put
\(r_0=e^{-2\pi T_0}<\epsilon\).  Define
\[
 \mathcal D_{12}
 =\frac{[0,r_0]\times(\overline T^2)^2}
 {\{0\}\times(\overline T^2)^2\sim d_{12}},
 \qquad
 \mathcal D_{24}
 =\frac{[0,r_0]\times(S^1)^4}
 {\{0\}\times(S^1)^4\sim d_{24}}.
\]
Both spaces are homeomorphic to \(\operatorname{Cone}(T^4)\).  There are
topological embeddings
\begin{align}
 \delta_{12}:\mathcal D_{12}&\hookrightarrow
 \mathcal A\wedge\mathcal A,
 &
 [r,u_1,u_2]&\longmapsto
 [\,[r,1,u_1],[r,1,u_2],],
 \label{wb:cusp-cone-D12-embedding}\\
 \delta_{24}:\mathcal D_{24}&\hookrightarrow
 \mathcal C^{\wedge4},
 &
 [r,\theta_1,\ldots,\theta_4]&\longmapsto
 \bigwedge_{j=1}^4[r,(1,\theta_j)].
 \label{wb:cusp-cone-D24-embedding}
\end{align}
Here \(1\in S^1\) fixes the positive radial base angle.  For every
\(T\geq T_0\), put \(r_T=e^{-2\pi T}\).  The restrictions of the sphere
quotient maps to the synchronized radial fibres are exactly the
cross-sections at height \(r_T\):
\begin{align}
 F_T^2&\xrightarrow{\ \pi_T\times\pi_T\ }
 Q_T^2\cong\{r_T\}\times T^4
 \subset\mathcal D_{12},
 \label{wb:cusp-cone-D12-slice}\\
 F_T^4&\xrightarrow{\ \chi_{\gamma,T}^{\times4}\ }
 C_T^4\cong\{r_T\}\times T^4
 \subset\mathcal D_{24}.
 \label{wb:cusp-cone-D24-slice}
\end{align}
Every one of these embedded four-torus slices is null-homotopic in its
carrier, and hence in the full smash target.

At the cone vertex \(d_\nu\), \(\nu\in\{12,24\}\), the local integral
homology is
\begin{equation}
 H_k(\mathcal D_\nu,
       \mathcal D_\nu\setminus\{d_\nu\};\mathbb Z)
 \cong
 \begin{cases}
  \mathbb Z^4,&k=2,\\
  \mathbb Z^6,&k=3,\\
  \mathbb Z^4,&k=4,\\
  \mathbb Z,&k=5,\\
  0,&\text{otherwise}.
 \end{cases}
 \label{wb:cusp-cone-carrier-local-homology}
\end{equation}
In particular, \(d_\nu\) is not an integral homology-manifold point, with or
without boundary, of any dimension.  Consequently it is not a topological
manifold point, with or without boundary, of any dimension.
\end{proposition}



\begin{proof}
\begin{enumerate}[label=\arabic*.,leftmargin=*]
\item The two displayed domains are cones on
\((\overline T^2)^2\cong T^4\) and \((S^1)^4\cong T^4\).  At \(r=0\),
every factor on the right of
\cref{wb:cusp-cone-D12-embedding} and \cref{wb:cusp-cone-D24-embedding} is the relevant
smash basepoint, so the formulae descend from the cylinders to the cones.
For \(0<r\leq r_0\), none of the factors is the smash basepoint.  Equality
in the smash quotient is then equality in every factor, and the coordinate
descriptions of the suspension and cusp cone give equality of \(r\) and of
all torus coordinates.  Thus each \(\delta_\nu\) is injective.  Its domain
is compact and its target is Hausdorff, so it is an embedding.

\item The already proved one-factor restriction formulae, with
\(r_T=e^{-2\pi T}\), are
\[
 q_2(x)=[r_T,1,\iota_T^{-1}(\pi_T(x))]
 \quad(x\in F_T)
\]
and, in the cusp-cone piece of \(\mathcal C\),
\[
 q_1(x)
 =[r_T,(1,\vartheta_T^{-1}(\chi_{\gamma,T}(x)))]
 \quad(x\in F_T).
\]
Applying these identities in two and four factors gives
\cref{wb:cusp-cone-D12-slice} and \cref{wb:cusp-cone-D24-slice}.  Sliding the cone coordinate
from \(r_T\) to zero contracts each slice to the cone vertex.  Finally,
excision and the long exact sequence of the pair give, for \(L=T^4\),
\[
 H_k(\operatorname{Cone}(L),
     \operatorname{Cone}(L)\setminus\{v\};\mathbb Z)
 \cong\widetilde H_{k-1}(L;\mathbb Z).
\]
The ranks \(4,6,4,1\) of \(H_1,H_2,H_3,H_4\) of \(T^4\) prove
\cref{wb:cusp-cone-carrier-local-homology}.  At an interior point of an
integral homology \(n\)-manifold the local integral homology is \(\mathbb Z\)
in degree \(n\) and zero in every other degree, whereas at a boundary point
of an integral homology manifold with boundary all these local groups
vanish.  The simultaneous nonzero groups in degrees \(2,3,4,5\), including
ranks greater than one, agree with neither pattern for any \(n\).  Every
topological manifold is locally an integral homology manifold, so the final
assertion follows as well.
\end{enumerate}
\end{proof}

\begin{theorem}[Point-set quotient maps and exact homotopy types]
\label{wb:cusp-cone-target-homotopy-types}
\StatementRecord{wb:cusp-cone-target-homotopy-types}{proved}{Historical events: GPT-6 Astra (OpenAI), ChatGPT Pro mode, effort Pro, S6\_09\_PAPERS, 17 November 2026: targeted proof check. Also checked in this writing. Pre-continuation writing events: GPT-6 Astra (OpenAI), ChatGPT Pro mode, effort Pro, S6\_10\_PAPERS, 18 November 2026: targeted proof check. Continuation metadata events: GPT-6 Astra (OpenAI), ChatGPT Pro mode, effort Pro, S6\_10\_PAPERS continuation, 19 November 2026: used and rechecked in this run; status-vocabulary reconciliation; metadata reconciliation only; no new mathematical check. \emph{Pre-continuation recorded proof scope.} GPT-6 Astra (OpenAI, ChatGPT Pro mode), S6\_09\_PAPERS, effort Pro, 2026-11-17; also checked in this writing by GPT-6 Astra (OpenAI), ChatGPT Pro mode, effort Pro, 18 November 2026}{the hypotheses, notation and definitions stated here; the explicit constructions and calculations in the proof}

The smash maps are continuous surjective quotient maps
\[
 q_2^{\wedge2}:S^{12}\twoheadrightarrow\mathcal A^{\wedge2},
 \qquad
 q_1^{\wedge4}:S^{24}\twoheadrightarrow\mathcal C^{\wedge4}.
\]
Their targets have homotopy types
\begin{align}
 \mathcal A^{\wedge2}
 &\simeq
 \bigvee^9S^4\vee\bigvee^{18}S^5\vee
 \bigvee^{15}S^6\vee\bigvee^6S^7\vee S^8,
 \label{wb:cusp-cone-A-smash-wedge}\\
 \mathcal C^{\wedge4}
 &\simeq
 S^8\vee\bigvee^4S^9\vee\bigvee^6S^{10}
 \vee\bigvee^4S^{11}\vee S^{12}.
 \label{wb:cusp-cone-C-smash-wedge}
\end{align}
Both maps induce zero on every positive-degree integral homology and
cohomology group and have zero ordinary Hurewicz image.  Neither map has a
continuous homotopy section.  These conclusions do not determine whether
the unstable homotopy class of either map is zero.
\end{theorem}



\begin{proof}
\begin{enumerate}[label=\arabic*.,leftmargin=*]
\item Let \(f:(K,k_0)\to(L,\ell_0)\) be a surjective based map between compact
Hausdorff spaces.  Every non-basepoint class of \(L^{\wedge r}\) has a
representative \((y_1,\ldots,y_r)\) with each \(y_i\ne\ell_0\).  Choosing
\(x_i\in f^{-1}(y_i)\) proves surjectivity of \(f^{\wedge r}\); the smash
basepoint is also in its image.  The source smash is compact.  The target
fat wedge is closed, hence the target smash is Hausdorff.  A continuous
surjection from a compact space to a Hausdorff space is a quotient map.
This proves the first assertion for \(q_2\) and \(q_1\).

\item The suspension splitting
\[
 \Sigma(X_1\times\cdots\times X_n)
 \simeq
 \bigvee_{\varnothing\ne I\subseteq\{1,\ldots,n\}}
 \Sigma\!\left(\bigwedge_{i\in I}X_i\right)
\]
gives
\[
 \mathcal A=\Sigma T^3
 \simeq\bigvee^3S^2\vee\bigvee^3S^3\vee S^4.
\]
Finite smash products distribute over finite wedges.  The ordered degree
sums give multiplicities \(9,18,15,6,1\) in dimensions
\(4,5,6,7,8\), proving \cref{wb:cusp-cone-A-smash-wedge}.

\item For \(\mathcal C\), write \(x=\rho_1\), \(y=\rho_2\), and
\(h=\eta\).  The pair-of-pants product and the two exact Seifert fillings
give
\[
 \pi_1(\mathcal M)=
 \langle x,y,h\mid[x,h]=[y,h]=1,
 x^3h^{-1}=1,\ y^4h=1\rangle.
\]
The remaining boundary torus is generated by \(h\) and
\(\rho_0=(xy)^{-1}\).  Coning this torus kills their normal closure, and
therefore
\[
 \pi_1(\mathcal C)
 \cong\langle x,y\mid x^3=1,\ y^4=1,\ xy=1\rangle=1,
\]
because \(y=x^{-1}\) and \(\gcd(3,4)=1\).  The previously computed
relative homology gives
\[
 \widetilde H_2(\mathcal C;\mathbb Z)=\mathbb Z,
 \qquad
 \widetilde H_3(\mathcal C;\mathbb Z)=\mathbb Z,
\]
and no other reduced homology.  % checked carrier correction: relative-hurewicz
Since \(\mathcal C\) is simply connected, Hurewicz gives a based
map \(a:S^2\to\mathcal C\) inducing an isomorphism on \(H_2\)
and on \(\pi_2\).  Replace this map by its mapping-cylinder
inclusion \(S^2\hookrightarrow M_a\), with
\(M_a\simeq\mathcal C\).  Both spaces are simply connected and
the map on \(\pi_2\) is an isomorphism, so the pair is
\(2\)-connected.  Relative Hurewicz gives
\[
 \pi_3(M_a,S^2)\cong H_3(M_a,S^2;\mathbb Z).
\]
The homology exact sequence identifies the right side with
\(H_3(M_a;\mathbb Z)=\mathbb Z\): the preceding group
\(H_3(S^2)\) is zero, and the following map on \(H_2\) is an
isomorphism.  The homotopy exact sequence makes
\(\pi_3(M_a)\to\pi_3(M_a,S^2)\) surjective, since its next
map has image in the zero kernel of
\(\pi_2(S^2)\to\pi_2(M_a)\).  Naturality of Hurewicz therefore
gives a based map \(b:S^3\to\mathcal C\) representing an integral
generator of \(H_3(\mathcal C)\).  The wedge \(a\vee b\) is an
integral homology equivalence in every degree.  The homological
Whitehead theorem, with its complete proof in
\cref{mapsI:pf:theorem:5.1}, makes this map a homotopy equivalence.  Thus
\[
 \mathcal C\simeq S^2\vee S^3.
\]
Taking the fourth smash power yields \cref{wb:cusp-cone-C-smash-wedge} with
the binomial multiplicities \(1,4,6,4,1\).

\item The first target has no reduced homology above degree eight and the second
has none above degree twelve.  The source spheres have no positive homology
below degrees twelve and twenty-four.  Hence all positive-degree homology
and cohomology homomorphisms induced by the two maps vanish.  Their
Hurewicz images lie in \(H_{12}(\mathcal A^{\wedge2})=0\) and
\(H_{24}(\mathcal C^{\wedge4})=0\), respectively.  A section would split a
nonzero lower homology group of the target into the zero lower homology of
the source sphere, which is impossible.  Higher homotopy groups of the
displayed wedges need not vanish, so none of these arguments decides the
unstable null-homotopy question.
\end{enumerate}
\end{proof}

\begin{theorem}[Exact decompactification spaces with boundary
\(\mathbb R^4\)]
\label{wb:cusp-cone-exact-decompactification-space}
\StatementRecord{wb:cusp-cone-exact-decompactification-space}{proved}{Historical events: GPT-6 Astra (OpenAI), ChatGPT Pro mode, effort Pro, S6\_09\_PAPERS, 17 November 2026: targeted proof check. Also checked in this writing. Pre-continuation writing events: GPT-6 Astra (OpenAI), ChatGPT Pro mode, effort Pro, S6\_10\_PAPERS, 18 November 2026: targeted proof check. Continuation metadata events: GPT-6 Astra (OpenAI), ChatGPT Pro mode, effort Pro, S6\_10\_PAPERS continuation, 19 November 2026: used and rechecked in this run; status-vocabulary reconciliation; metadata reconciliation only; no new mathematical check. \emph{Pre-continuation recorded proof scope.} GPT-6 Astra (OpenAI, ChatGPT Pro mode), S6\_09\_PAPERS, effort Pro, 2026-11-17; also checked in this writing by GPT-6 Astra (OpenAI), ChatGPT Pro mode, effort Pro, 18 November 2026}{the hypotheses, notation and definitions stated here; the explicit constructions and calculations in the proof}

For \(\nu\in\{12,24\}\), set
\[
 A_T^{(12)}=\operatorname{diag}(Y_T,Y_T),
 \qquad
 A_T^{(24)}=\ell_T I_4,
 \qquad s_0=T_0^{-1}.
\]
On \([0,s_0]\times\mathbb R^4\), define an equivalence relation
\(\sim_\nu\) fibrewise over \(s\) by
\begin{align*}
 (0,x)\sim_\nu(0,y)&\iff x=y,\\
 (s,x)\sim_\nu(s,y)&\iff
 x-y\in A_{1/s}^{(\nu)}\mathbb Z^4
 \quad(0<s\leq s_0).
\end{align*}
Let
\[
 \mathfrak E_\nu
 =([0,s_0]\times\mathbb R^4)/\!\sim_\nu,
 \qquad p_\nu:\mathfrak E_\nu\longrightarrow[0,s_0].
\]
Then \(\mathfrak E_\nu\) is a Hausdorff, second-countable topological
five-manifold with boundary, and
\begin{align}
 p_{12}^{-1}(s)&\cong Q_{1/s}^2 &(s>0),&
 &p_{12}^{-1}(0)&\cong\mathbb R^4,
 \label{wb:cusp-cone-E12-fibres}\\
 p_{24}^{-1}(s)&\cong C_{1/s}^4 &(s>0),&
 &p_{24}^{-1}(0)&\cong\mathbb R^4.
 \label{wb:cusp-cone-E24-fibres}
\end{align}
The other boundary component is the compact four-torus over \(s=s_0\).

Give every positive fibre its Euclidean quotient metric and the zero fibre
its Euclidean metric.  For every \(R>0\), there is \(s_R>0\) such that,
for \(0\leq s<s_R\), writing
\(q_{\nu,s}(x)=[s,x]\), the quotient map restricts to an exact isometry
\[
 q_{\nu,s}:B_R^{\mathbb R^4}(0)
 \xrightarrow{\ \cong_{\mathrm{isom}}\ }
 B_R^{p_\nu^{-1}(s)}([s,0])
\]
onto the entire quotient radius-\(R\) ball, not merely onto an embedded
subset of that ball.
Consequently the positive fibres converge to the zero fibre in pointed
local-isometric, pointed smooth, and pointed Gromov--Hausdorff topology.

There are continuous surjective comparison maps which set-theoretically
collapse the zero boundary fibres,
\[
 \begin{tikzcd}[column sep=large]
 \mathbb R^4=p_{12}^{-1}(0)
 \arrow[r,hook]&
 \mathfrak E_{12}
 \arrow[r,"\beta_{12}"]&
 \mathcal D_{12}
 \arrow[r,hook,"\delta_{12}"]&
 \mathcal A^{\wedge2},
 \end{tikzcd}
\]
and
\[
 \begin{tikzcd}[column sep=large]
 \mathbb R^4=p_{24}^{-1}(0)
 \arrow[r,hook]&
 \mathfrak E_{24}
 \arrow[r,"\beta_{24}"]&
 \mathcal D_{24}
 \arrow[r,hook,"\delta_{24}"]&
 \mathcal C^{\wedge4}.
 \end{tikzcd}
\]
For \(s>0\), \(\beta_\nu\) is the specified marked homeomorphism from
\(p_\nu^{-1}(s)\) to the radial torus slice at \(r=e^{-2\pi/s}\).  Its
restriction
\[
 \beta_\nu:\mathfrak E_\nu\setminus p_\nu^{-1}(0)
 \xrightarrow{\ \cong\ }\mathcal D_\nu\setminus\{d_\nu\}
\]
is a homeomorphism, while
\[
 \beta_\nu^{-1}(d_\nu)=p_\nu^{-1}(0)\cong\mathbb R^4.
\]
Neither \(\beta_\nu\) is proper or a quotient map.  It factors through a
continuous bijection
\[
 \overline\beta_\nu:
 \mathfrak E_\nu/p_\nu^{-1}(0)\longrightarrow\mathcal D_\nu,
\]
but this bijection is not a homeomorphism; equivalently, the topology of
\(\mathcal D_\nu\) is not the quotient topology induced by \(\beta_\nu\).
Thus the limiting \(\mathbb R^4\) is an actual boundary fibre of
\(\mathfrak E_\nu\), whereas the continuous comparison map sends that entire
fibre to the cone vertex.
\end{theorem}



\begin{proof}
\begin{enumerate}[label=\arabic*.,leftmargin=*]
\item We give a common proof.  Write
\[
 L_s=A_{1/s}^{(\nu)}\mathbb Z^4\quad(s>0),
 \qquad L_0=\{0\}.
\]
The matrices vary continuously for \(s>0\), and
\[
 \operatorname{sys}(L_s)\longrightarrow\infty
 \qquad(s\downarrow0)
\]
by \(\operatorname{sys}(Y_T\mathbb Z^2)\geq T-C\) in the
\(S^{12}\) case and \(\operatorname{sys}(\ell_T\mathbb Z^4)=\ell_T\)
in the \(S^{24}\) case.

\item The equivalence relation is closed.  Indeed, suppose
\((s_j,x_j)\sim_\nu(s_j,y_j)\) and
\((s_j,x_j,y_j)\to(s,x,y)\).  If \(s>0\), write
\(y_j-x_j=A_{1/s_j}^{(\nu)}n_j\).  On a compact neighborhood of \(s\),
the inverse matrices are uniformly bounded, so boundedness of
\(y_j-x_j\) forces \(n_j\) to lie in a finite set.  Passing to a
subsequence makes \(n_j\) constant and proves \(x\sim_\nu y\).  If
\(s=0\), any nonzero \(n_j\) would give
\[
 \|y_j-x_j\|\geq\operatorname{sys}(L_{s_j})\longrightarrow\infty,
\]
contrary to convergence.  Hence \(x=y\).

\item The quotient projection \(q_\nu:[0,s_0]\times\mathbb R^4\to
\mathfrak E_\nu\) is open.  If \(U\) is open, its saturation is the union
of \(U\) with the images, for \(n\in\mathbb Z^4\), of
\[
 U\cap((0,s_0]\times\mathbb R^4)
 \quad\text{under}\quad
 (s,x)\longmapsto(s,x+A_{1/s}^{(\nu)}n),
\]
and every one of these images is open.  A closed equivalence relation with
an open quotient map has Hausdorff quotient: two inequivalent points admit
a product neighborhood disjoint from the relation, and the saturations of
the two factor neighborhoods are disjoint open sets.  Since the source is
second countable and the quotient map is open, images of a countable base
form a countable base of the quotient.

\item Fix \(x_0\in\mathbb R^4\).  Choose \(a>0\) and \(s_1>0\) such that
\(\operatorname{sys}(L_s)>2a\) for \(0<s<s_1\).  The quotient projection
is one-to-one on
\([0,s_1)\times B_a^\circ(x_0)\), where
\(B_a^\circ(x_0)=\{x:\|x-x_0\|<a\}\), and its restriction is an open continuous
bijection onto an open subset of \(\mathfrak E_\nu\); it is therefore a
homeomorphism.  These are half-space charts at \(s=0\).  The usual evenly
covered small balls give product charts for \(s>0\), and the same argument
at \(s=s_0\) gives the second boundary component.  This proves the
manifold assertion and the fibre identifications.

\item For the metric assertion, choose \(s_R\) so that
\(\operatorname{sys}(L_s)>4R\) for \(s<s_R\).  If
\(x,y\in B_R(0)\) and \(0\ne\lambda\in L_s\), then
\[
 \|x-y+\lambda\|
 \geq\|\lambda\|-\|x-y\|
 >4R-2R\geq\|x-y\|.
\]
Thus no translated representative shortens the Euclidean distance on that
ball, so \(q_{\nu,s}\) is an exact isometric embedding there.  It is also
onto the full quotient ball.  Indeed, if \(\xi=[s,x]\) satisfies
\(d([s,0],\xi)\leq R\), discreteness of \(L_s\) makes the infimum in
\[
 d([s,0],[s,x])=\min_{\lambda\in L_s}\|x+\lambda\|
\]
attained.  For a minimizing \(\lambda_0\), the representative
\(x+\lambda_0\) lies in \(B_R^{\mathbb R^4}(0)\) and maps to \(\xi\).  At
\(s=0\), \(q_{\nu,0}\) is the identity.  This proves the asserted onto
isometry, and since it holds for every fixed \(R\), it proves all three
pointed convergences.

\item It remains to define these comparison maps, which set-theoretically collapse
the zero fibres.  For \(s>0\), put
\(T=1/s\) and \(r=e^{-2\pi T}\).  In the \(S^{12}\) case, write
\(x=(x_1,x_2)\in\mathbb R^2\oplus\mathbb R^2\) and set
\[
 \beta_{12}([s,x])
 =\bigl[r,
 \iota_T^{-1}([x_1]_{Y_T\mathbb Z^2}),
 \iota_T^{-1}([x_2]_{Y_T\mathbb Z^2})\bigr].
\]
In the \(S^{24}\) case, write
\(x=(x_1,\ldots,x_4)\) and set
\[
 \beta_{24}([s,x])
 =\bigl[r,
 [x_1/\ell_T],\ldots,[x_4/\ell_T]\bigr].
\]
These formulae are well defined on the positive fibres and are continuous
there.  Define \(\beta_\nu([0,x])=d_\nu\) for every \(x\in\mathbb R^4\).
Every neighborhood of \(d_\nu\) in the cone contains all cross-sections
with sufficiently small \(r\).  Since \(e^{-2\pi/s}\to0\) as
\(s\downarrow0\), continuity at the zero fibre is uniform in \(x\).
Surjectivity follows because each positive fibre maps onto its complete
cross-section and the zero fibre maps onto the vertex.  The restriction away
from \(p_\nu^{-1}(0)\) is a homeomorphism because
\(s\mapsto e^{-2\pi/s}\) is a homeomorphism from \((0,s_0]\) to
\((0,r_0]\), and the displayed marked-torus identifications are fibrewise
homeomorphisms varying continuously with \(s\).  The vertex fibre is exactly
\(p_\nu^{-1}(0)\), so \(\beta_\nu\) is not proper.

\item It is also not a quotient map.  Choose
\(a\in\mathbb R^4\setminus\mathbb Z^4\), for example
\(a=(1/2,1/2,1/2,1/2)\), choose \(s_j\downarrow0\), and put
\[
 z_j=[s_j,A_{1/s_j}^{(\nu)}a]\in\mathfrak E_\nu.
\]
For \(\nu=12\),
\[
 \inf_{n\in\mathbb Z^4}
 \|A_{1/s_j}^{(12)}(a+n)\|
 \geq(1/s_j-C)\operatorname{dist}(a,\mathbb Z^4)
 \longrightarrow\infty,
\]
and for \(\nu=24\) the same infimum equals
\[
 \ell_{1/s_j}\operatorname{dist}(a,\mathbb Z^4)
 \longrightarrow\infty.
\]
More explicitly, for every fixed \(x_0\in\mathbb R^4\),
\[
 \inf_{n\in\mathbb Z^4}
 \|A_{1/s_j}^{(\nu)}(a+n)-x_0\|
 \geq
 \inf_{n\in\mathbb Z^4}\|A_{1/s_j}^{(\nu)}(a+n)\|-\|x_0\|
 \longrightarrow\infty.
\]
Hence \(\{z_j:j\geq1\}\) is closed in \(\mathfrak E_\nu\): it has no
accumulation in a positive fibre and the last estimate shows that it
eventually misses every half-space chart about every point of the zero
fibre.  But
\[
 \beta_\nu(z_j)=[e^{-2\pi/s_j},a]\longrightarrow d_\nu,
\]
so \(\{\beta_\nu(z_j):j\geq1\}\) is not closed in \(\mathcal D_\nu\).
Since \(\beta_\nu\) is injective off the zero fibre and the latter maps to
\(d_\nu\), the inverse image of this nonclosed set is the closed set
\(\{z_j:j\geq1\}\).  Thus \(\beta_\nu\) is not a quotient map.  This proves
the two diagrams and their exact topological scope.  Since \(\beta_\nu\) is
constant exactly on \(p_\nu^{-1}(0)\) and otherwise injective, it induces the
continuous bijection \(\overline\beta_\nu\) displayed in the theorem.  Were
that bijection a homeomorphism, \(\beta_\nu\) would be the composite of the
quotient projection with a homeomorphism and hence a quotient map, contrary
to the counterexample.  Therefore \(\overline\beta_\nu\) is not a
homeomorphism.
\end{enumerate}
\end{proof}

\begin{corollary}[Sharp global meaning of the Euclidean limit]
\label{wb:cusp-cone-sharp-global-R4-scope}
\StatementRecord{wb:cusp-cone-sharp-global-R4-scope}{proved}{Historical events: GPT-6 Astra (OpenAI), ChatGPT Pro mode, effort Pro, S6\_09\_PAPERS, 17 November 2026: targeted proof check. Also checked in this writing. Pre-continuation writing events: GPT-6 Astra (OpenAI), ChatGPT Pro mode, effort Pro, S6\_10\_PAPERS, 18 November 2026: targeted proof check. Continuation metadata events: GPT-6 Astra (OpenAI), ChatGPT Pro mode, effort Pro, S6\_10\_PAPERS continuation, 19 November 2026: used and rechecked in this run; status-vocabulary reconciliation; metadata reconciliation only; no new mathematical check. \emph{Pre-continuation recorded proof scope.} GPT-6 Astra (OpenAI, ChatGPT Pro mode), S6\_09\_PAPERS, effort Pro, 2026-11-17; also checked in this writing by GPT-6 Astra (OpenAI), ChatGPT Pro mode, effort Pro, 18 November 2026}{the hypotheses, notation and definitions stated here; \cref{wb:cusp-cone-literal-R4-open-cells}; \cref{wb:cusp-cone-explicit-carriers}; \cref{wb:cusp-cone-exact-decompactification-space}; \cref{wb:cusp-cone-carrier-local-homology}; \cref{wb:sphere-metric-fixed-realization}}

The native routes constructed above exhibit the underlying space
\(\mathbb R^4\) in the following three precise senses:
\begin{enumerate}
\item as each of the nonclosed topologically embedded open cells
\(U_{12,T}\subset S^6\) and \(U_{24,T}\subset S^8\);
\item as the pointed intrinsic Euclidean limit of the four-torus slices in the
global quotient carriers of \(S^{12}\) and \(S^{24}\); and
\item as the Euclidean decompactification boundary fibre of
\(\mathfrak E_{12}\) and \(\mathfrak E_{24}\).
\end{enumerate}
It does not produce \(\mathbb R^4\) as the whole compact sphere, as the
whole compact conical target, or as a continuous quotient target of either
compact space.  Nor does the unrescaled restriction of any fixed compatible
ambient target metric have \(\mathbb R^4\) as its slice limit, and the open
cells in item \(1\) are not globally Euclidean-isometric in such a metric.
\end{corollary}



\begin{proof}
\begin{enumerate}[label=\arabic*.,leftmargin=*]
\item The three positive assertions are
\cref{wb:cusp-cone-literal-R4-open-cells} and \cref{wb:cusp-cone-explicit-carriers} and \cref{wb:cusp-cone-exact-decompactification-space}.
For the negative assertions, a continuous image of a compact sphere or
compact conical target is compact, so it cannot be all of \(\mathbb R^4\).
The local homology calculation
\cref{wb:cusp-cone-carrier-local-homology} shows that the carrier vertex is not
an integral homology-manifold point, and hence not a manifold point, of any
dimension, with or without boundary.  Finally, in every fixed compatible
metric on a compact cone target, all sufficiently deep radial slices lie
uniformly in an arbitrarily small neighborhood of the vertex, so their ambient
diameters tend to zero.  Their intrinsic diameters tend to infinity because
their intrinsic systoles do.  Hence without a \(T\)-dependent rescaling no
fixed ambient metric realizes the intrinsic Euclidean limit on those slices.
The fixed cusp-adapted blow-ups are treated separately in
\cref{wb:sphere-metric-fixed-realization}.  The comparison diagrams record the
unrescaled distinction exactly: they collapse the new \(\mathbb R^4\)
boundary rather than identify it with the compact target.
\end{enumerate}
\end{proof}

\paragraph{Scope.}
The constructions above use only the original, unscaled fibre metrics and
the quotient maps derived from the native \(S^6\) cusp.  They do not use the
auxiliary finite-cover unspooling.  The conical carriers depend on the
chosen cusp coordinate, positive radial ray, and based identification
\(X\cong S^6\).  The spaces \(\mathfrak E_\nu\) additionally remember the
marked Euclidean quotient metrics.  None of these data defines a compact
complex structure on \(S^{12}\) or \(S^{24}\), a four-dimensional physical
spacetime, or an interacting Yang--Mills Hamiltonian.

\subsection{The full \texorpdfstring{\(S^{12}\to S^6\)}{S12-to-S6}
carrier is null-homotopic}
\label{wb:cusp-s12-carrier-homotopy}

This subsection determines the based homotopy class of the full two-copy
sphere carrier.  It is important that the calculation concerns the full
quotient map, not merely the inclusion of one of its moving four-spheres.

Recall the exact maps from
\cref{wb:cusp-cone-S12-route} and \cref{wb:cusp-cone-angle-forgotten-sphere-carriers}.  The
cusp quotient is
\[
 q_2:(X,x_0)\longrightarrow
 \mathcal A=\operatorname{Susp}(S^1\times\overline T^2),
\]
and, for the coordinate one-skeleton
\[
 K^{(1)}=
 ((\mathbb R/\mathbb Z)\times\{0\})
 \cup(\{0\}\times(\mathbb R/\mathbb Z))\subset\overline T^2,
\]
the degree-one quotient
\(c:\overline T^2\to\overline T^2/K^{(1)}\cong S^2\) defines
\[
 \widehat p_2([r,e^{i\phi},u])=[r,c(u)]
 \in\Sigma_r S^2\cong S^3.
\]
Thus
\[
 g:=\widehat q_3=\widehat p_2\circ q_2:(X,x_0)\longrightarrow(S^3,*)
\]
is the specified continuous surjective based quotient.  After the fixed
based identification \(X\cong_{\mathrm{diff}}S^6\), its smash square is the
specified continuous surjective based quotient
\[
 g^{\wedge2}:S^{12}=S^6\wedge S^6
 \longrightarrow S^3\wedge S^3=S^6.
\]

\begin{theorem}[The full external smash pairing is zero]
\label{wb:cusp-s12-full-carrier-null}
\StatementRecord{wb:cusp-s12-full-carrier-null}{proved}{Historical events: GPT-6 Astra (OpenAI), ChatGPT Pro mode, effort Pro, S6\_09\_PAPERS, 17 November 2026: targeted proof check. Also checked in this writing. Pre-continuation writing events: GPT-6 Astra (OpenAI), ChatGPT Pro mode, effort Pro, S6\_10\_PAPERS, 18 November 2026: targeted proof check. Continuation metadata events: GPT-6 Astra (OpenAI), ChatGPT Pro mode, effort Pro, S6\_10\_PAPERS continuation, 19 November 2026: used and rechecked in this run; status-vocabulary reconciliation; metadata reconciliation only; no new mathematical check. \emph{Pre-continuation recorded proof scope.} GPT-6 Astra (OpenAI, ChatGPT Pro mode), S6\_09\_PAPERS, effort Pro, 2026-11-17; also checked in this writing by GPT-6 Astra (OpenAI), ChatGPT Pro mode, effort Pro, 18 November 2026}{the hypotheses, notation and definitions stated here; \cite[opening table, p.~339]{Hatcher2002}; \cite[Table~1, p.~113]{IsaksenWangXu2023}; \cref{wb:cusp-s12-toda-smash-formula}}

For every pair of based maps \(x,y:S^6\to S^3\),
\[
 [x\wedge y]=0\qquad\text{in}\qquad\pi_{12}(S^6).
\]
In particular,
\(\widehat q_3^{\wedge2}:S^{12}\twoheadrightarrow S^6\) is simultaneously
surjective as a point-set quotient and based null-homotopic.
\end{theorem}



\begin{proof}
\begin{enumerate}[label=\arabic*.,leftmargin=*]
\item Write \(E\) for reduced suspension, using the ordered smash coordinates on
spheres.  The external smash product has the exact factorization
\begin{equation}
 x\wedge y
 =
 (x\wedge\operatorname{id}_{S^3})
 \circ
 (\operatorname{id}_{S^6}\wedge y):
 S^6\wedge S^6\longrightarrow S^6\wedge S^3
 \longrightarrow S^3\wedge S^3.
 \label{wb:cusp-s12-smash-factorization}
\end{equation}
Under the displayed ordered sphere identifications, the two factors, in
order of application, represent \(E^6y\in\pi_{12}(S^9)\) and
\(E^3x\in\pi_9(S^6)\).  Equivalently, Toda's
Barratt--Hilton reduced-join formula gives, with
\(p=q=k=h=3\),
\begin{equation}
 [x\wedge y]=(E^3[x])\circ(E^6[y]).
 \label{wb:cusp-s12-toda-smash-formula}
\end{equation}
Here Toda's sign is \((-1)^{(p+k)h}=(-1)^{18}=+1\).  Reversing the
two displayed factors introduces the separate twist sign
\((-1)^{ph}=-1\); no such reversal is used here.

\item The three exact groups needed below are
\begin{equation}
 \pi_6(S^3)\cong\mathbb Z/12,
 \qquad
 \pi_{12}(S^9)\cong\mathbb Z/24,
 \qquad
 \pi_{12}(S^6)\cong\mathbb Z/2.
 \label{wb:cusp-s12-three-stem-groups}
\end{equation}
The first and third are the integral low unstable groups displayed in
Hatcher's table extracted from Toda; the middle group is in the stable
three-stem
\cite[opening table, p.~339]{Hatcher2002};
\cite[Table~1, p.~113]{IsaksenWangXu2023}.

\item Put \(\beta=E^6[y]\in\pi_{12}(S^9)\).  Since
\(12[y]=0\), functoriality of suspension gives \(12\beta=0\).  In the
cyclic group \(\mathbb Z/24\), the elements annihilated by \(12\) are
exactly the even elements.  Hence there is
\(\gamma\in\pi_{12}(S^9)\) such that \(\beta=2\gamma\).
Postcomposition with \(E^3[x]:S^9\to S^6\) induces a group homomorphism
\[
 (E^3[x])_*:\pi_{12}(S^9)\longrightarrow\pi_{12}(S^6).
\]
The target has exponent two, so \cref{wb:cusp-s12-toda-smash-formula} gives
\[
 [x\wedge y]
 =(E^3[x])_*(\beta)
 =2(E^3[x])_*(\gamma)
 =0.
\]
This proves the vanishing of the entire external smash pairing and hence
the theorem.
\end{enumerate}
\end{proof}

\begin{corollary}[All higher smash powers]
\label{wb:cusp-s12-all-higher-smashes-null}
\StatementRecord{wb:cusp-s12-all-higher-smashes-null}{proved}{Historical events: GPT-6 Astra (OpenAI), ChatGPT Pro mode, effort Pro, S6\_09\_PAPERS, 17 November 2026: targeted proof check. Also checked in this writing. Pre-continuation writing events: GPT-6 Astra (OpenAI), ChatGPT Pro mode, effort Pro, S6\_10\_PAPERS, 18 November 2026: targeted proof check. Continuation metadata events: GPT-6 Astra (OpenAI), ChatGPT Pro mode, effort Pro, S6\_10\_PAPERS continuation, 19 November 2026: used and rechecked in this run; status-vocabulary reconciliation; metadata reconciliation only; no new mathematical check. \emph{Pre-continuation recorded proof scope.} GPT-6 Astra (OpenAI, ChatGPT Pro mode), S6\_09\_PAPERS, effort Pro, 2026-11-17; also checked in this writing by GPT-6 Astra (OpenAI), ChatGPT Pro mode, effort Pro, 18 November 2026}{the hypotheses, notation and definitions stated here; the explicit constructions and calculations in the proof}

For every based \(g:S^6\to S^3\) and every \(r\geq2\),
\[
 g^{\wedge r}:S^{6r}\longrightarrow S^{3r}
\]
is based null-homotopic.
\end{corollary}



\begin{proof}
\begin{enumerate}[label=\arabic*.,leftmargin=*]
\item With the fixed associativity homeomorphisms for smash products,
\[
 g^{\wedge r}=g^{\wedge2}\wedge g^{\wedge(r-2)}.
\]
Smashing a based null-homotopy with a based map gives a based
null-homotopy.
\end{enumerate}
\end{proof}


\paragraph{The one-copy invariant and its computed value}
\label{wb:cusp-s12-one-copy-invariant}
\StatementRecord{wb:cusp-s12-one-copy-invariant}{proved}{Historical events: GPT-6 Astra (OpenAI), ChatGPT Pro mode, effort Pro, S6\_09\_PAPERS, 17 November 2026: targeted proof check. Also checked in this writing. Pre-continuation writing events: GPT-6 Astra (OpenAI), ChatGPT Pro mode, effort Pro, S6\_10\_PAPERS, 18 November 2026: targeted proof check. Continuation metadata events: GPT-6 Astra (OpenAI), ChatGPT Pro mode, effort Pro, S6\_10\_PAPERS continuation, 19 November 2026: used and rechecked in this run; status-vocabulary reconciliation; metadata reconciliation only; no new mathematical check. \emph{Pre-continuation recorded proof scope.} GPT-6 Astra (OpenAI, ChatGPT Pro mode), S6\_09\_PAPERS, effort Pro, 2026-11-17; also checked in this writing by GPT-6 Astra (OpenAI), ChatGPT Pro mode, effort Pro, 18 November 2026}{the hypotheses, notation and definitions stated here; \cite[Theorem~5.41(b), p.~573, and pp.~577--578]{HatcherSpectralSequences}; \cref{wb:cusp-null:thm}}

The proof does not need to determine the one-copy class
\([g]\in\pi_6(S^3)\cong\mathbb Z/12\).  Toda's suspension calculation
identifies \(\pi_6(S^3)\) with the even subgroup of the stable three-stem;
Hatcher gives an integral Serre-spectral-sequence derivation of the same
injectivity and identifies twice the stable Hopf--\(\nu\) generator with a
generator of the suspended \(\mathbb Z/12\)-summand
\cite[Theorem~5.41(b), p.~573, and pp.~577--578]{HatcherSpectralSequences}.
Thus, after choosing a stable generator \(\nu_9\), there is a unique
\(a\in\mathbb Z/12\) such that
\begin{equation}
 [E^6g]=2a\nu_9\in\pi_{12}(S^9)\cong\mathbb Z/24.
 \label{wb:cusp-s12-sixfold-suspension-class}
\end{equation}
Because \(a\mapsto2a\) embeds \(\mathbb Z/12\) as the even subgroup of
\(\mathbb Z/24\), \(E^6g\) is null-homotopic if and only if \(g\) is
null-homotopic.  Thus the one-copy class and its sixfold suspension are the
same integer modulo \(12\). It is now zero by Theorem~\ref{wb:cusp-null:thm}.

The exact geometric invariant which computes \(a\) is also identifiable.
Choose \(0<r_0<\epsilon\) and
\(u_0\in\overline T^2\setminus K^{(1)}\), and let
\(y_0=[r_0,c(u_0)]\) lie in the open three-cell of
\(\Sigma_rS^2\).  From the literal formula for \(q_2\),
\begin{equation}
 g^{-1}(y_0)
 =\{x\in f^{-1}(|t|=r_0):\overline\pi_f(x)=u_0\}.
 \label{wb:cusp-s12-regular-inverse-image}
\end{equation}
The fibre of \(\overline\pi\) is the two-torus defined by
\(K_\Lambda=\mathbb Z\langle\widehat w,\widehat\delta\rangle\), and the
cusp monodromy fixes both \(\widehat w\) and \(\widehat\delta\).  Therefore
\eqref{wb:cusp-s12-regular-inverse-image} is the mapping torus of an affine translation of the kernel torus,
which is diffeomorphic to \(S^1\times T^2\cong T^3\): if the translation is
\(k\mapsto k+v\), choose a real lift \(\widetilde v\); the map
\([\theta,k]\mapsto(e^{2\pi i\theta},k+\theta\widetilde v)\) respects
\((1,k)\sim(0,k+v)\) and has the inverse obtained by subtracting
\(\theta\widetilde v\). After choosing a smooth chart in the open
target cell, the differential of \(g\) supplies its normal three-framing.
Pontryagin--Thom identifies the stabilized class of \(g\) with the framed
bordism class of this specifically framed \(T^3\).  Since stabilization on
\(\pi_6(S^3)\) is injective, computing that framing class computes \(a\)
exactly.  The unframed diffeomorphism type \(T^3\), point-set surjectivity,
and the zero ordinary homology map do not compute \(a\).






\paragraph{Scope}
\label{wb:cusp-s12-nullity-scope}
\StatementRecord{wb:cusp-s12-nullity-scope}{proved}{Historical events: GPT-6 Astra (OpenAI), ChatGPT Pro mode, effort Pro, S6\_09\_PAPERS, 17 November 2026: targeted proof check. Also checked in this writing. Pre-continuation writing events: GPT-6 Astra (OpenAI), ChatGPT Pro mode, effort Pro, S6\_10\_PAPERS, 18 November 2026: targeted proof check. Continuation metadata events: GPT-6 Astra (OpenAI), ChatGPT Pro mode, effort Pro, S6\_10\_PAPERS continuation, 19 November 2026: used and rechecked in this run; status-vocabulary reconciliation; metadata reconciliation only; no new mathematical check. \emph{Pre-continuation recorded proof scope.} GPT-6 Astra (OpenAI, ChatGPT Pro mode), S6\_09\_PAPERS, effort Pro, 2026-11-17; also checked in this writing by GPT-6 Astra (OpenAI), ChatGPT Pro mode, effort Pro, 18 November 2026}{the hypotheses, notation and definitions stated here; the argument and qualifications immediately following}

The theorem changes no point-set or metric statement.  A surjective quotient
map may be null-homotopic, and the embedded moving four-spheres and their
intrinsic pointed \(\mathbb R^4\)-limits remain exactly as constructed.  The
theorem concerns the sphere-valued postcomposition
\(\widehat q_3^{\wedge2}\); it does not imply null-homotopy of the earlier
conical-target map
\(q_2^{\wedge2}:S^{12}\to\mathcal A^{\wedge2}\).





\subsection{The original cusp coefficient: an explicit framed nullbordism}


Let \(X\) be the smooth compact source of the original construction,
\(f:X\to B\) its map, and \(t\) the specified cusp coordinate.
Put \(W=f^{-1}(0)\), choose the original \(x_0\in W\), and retain the
ordered quotient lattice basis induced by
\((\widehat\gamma,\widehat u)\), with kernel generated by
\((\widehat w,\widehat\delta)\). Thus the quotient torus is
\(\overline T^2=\mathbb R^2/\mathbb Z^2\). The original periods are
\begin{equation}
 Z(s)=sB_0+C(t),\qquad e^{2\pi i s}=t,\qquad
 B_0=\begin{pmatrix}0&1\\-1&0\end{pmatrix},\qquad
 C(t)=\begin{pmatrix}6\mu(t)&h(t)\\b(t)-h(t)&\mu(t)\end{pmatrix}.
 \label{wb:cusp-null:eq:periods}
\end{equation}
Every entry of \(C\) is holomorphic at \(t=0\).
The dense torus of the original toric cusp has coordinates
\((x_1,x_2,t)\), with
\(x=\exp(2\pi i\zeta)\). The residual lattice action is exactly
\begin{equation}
 \Psi_\lambda(x,t)=
 \bigl(t^{B_0\lambda}\exp(2\pi iC(t)\lambda)x,t\bigr),
 \qquad\lambda\in\mathbb Z^2.
 \label{wb:cusp-null:eq:action}
\end{equation}
It extends freely and properly discontinuously across the original
unimodular toric filling. In particular the quotient projection is a local
biholomorphism at every toric fixed point. These are the reconstructed
data and proved local properties of the original source
(the reconstructed original manuscript, \S4, source pp.~24--29; its complete cusp data are displayed here).

Let
\[
 K^{(1)}=(\mathbb R/\mathbb Z\times\{0\})\cup(\{0\}\times\mathbb R/\mathbb Z),\qquad
 c:\overline T^2\longrightarrow \overline T^2/K^{(1)}\cong S^2
\]
be the specified degree-one cell quotient. Its restriction to the open
two-cell is taken with the smooth coordinates used in the original
regular-fibre description. Both radial endpoints and the complete
basepoint cylinder are collapsed in
\[
 \Sigma_r(S^2,*)=
 \frac{[0,\epsilon]\times S^2}
 {(\{0,\epsilon\}\times S^2)\cup([0,\epsilon]\times\{*\})}
 \cong S^3.
\]
The original based map is
\begin{equation}
 g(x)=[|t|,c(\overline\pi_f(x))]\quad(0<|t|<\epsilon),\qquad
 g(W)=g(X\setminus\operatorname{int}N_0)=*.
 \label{wb:cusp-null:eq:g}
\end{equation}

\begin{theorem}[The original one-copy cusp map is null-homotopic]
\label{wb:cusp-null:thm}
\StatementRecord{wb:cusp-null:thm}{proved}{Historical events: GPT-6 Astra (OpenAI), ChatGPT Pro mode, effort Pro, S6\_09\_PAPERS, 17 November 2026: targeted proof check. Also checked in this writing. Pre-continuation writing events: GPT-6 Astra (OpenAI), ChatGPT Pro mode, effort Pro, S6\_10\_PAPERS, 18 November 2026: targeted proof check. Continuation metadata events: GPT-6 Astra (OpenAI), ChatGPT Pro mode, effort Pro, S6\_10\_PAPERS continuation, 19 November 2026: used and rechecked in this run; status-vocabulary reconciliation; metadata reconciliation only; no new mathematical check. \emph{Pre-continuation recorded proof scope.} GPT-6 Astra (OpenAI, ChatGPT Pro mode), S6\_09\_PAPERS, effort Pro, 2026-11-17; also checked in this writing by GPT-6 Astra (OpenAI), ChatGPT Pro mode, effort Pro, 18 November 2026}{the hypotheses, notation and definitions stated here; \cref{wb:cusp-null:eq:g}; \cref{wb:cusp-null:eq:moduli}; \cref{wb:cusp-null:eq:chart}; \cref{wb:cusp-null:eq:J}; \cref{wb:cusp-null:eq:spatial-bound}; \cref{wb:cusp-null:eq:bordism}; \cref{wb:cusp-null:eq:P}; \cite[\S7, Theorem B and Lemma 4,
pp.~43, 48--50]{CuspNullMilnor1965}}

The based map \(g:(X,x_0)\to(S^3,*)\) in \eqref{wb:cusp-null:eq:g} is based
null-homotopic. More precisely, a regular inverse image with its actual
ordered normal framing induced by \((r,c_1,c_2)\) is framed nullbordant
inside one original cusp chart times \([0,1]\), disjoint from
\(\{x_0\}\times[0,1]\).
\end{theorem}



\begin{proof}
\begin{enumerate}[label=\arabic*.,leftmargin=*]
\item The proof constructs the complete inverse image, compares its actual
framing with an explicitly extendible framing, and gives the extension.

\item \paragraph{1. The exact complete inverse image.}
Retain the lift
\[
 u^0=\begin{pmatrix}-1/4\\1/4\end{pmatrix},\qquad
 \overline u^0=(3/4,1/4)\pmod{\mathbb Z^2},\qquad
 B_0u^0=\begin{pmatrix}1/4\\1/4\end{pmatrix}.
\]
Thus \(\overline u^0\notin K^{(1)}\), and \(B_0u^0\) lies in the
interior of the original triangle \(\operatorname{conv}(0,e_1,e_2)\).
Use its original unimodular chart
\begin{equation}
 z_0=t/(x_1x_2),\qquad z_1=x_1,\qquad z_2=x_2,
 \qquad t=z_0z_1z_2.
 \label{wb:cusp-null:eq:chart}
\end{equation}
Set
\[
 \ell=\log|t|<0,\qquad R(t)=-2\pi\operatorname{Im}C(t),\qquad
 A(t)=\ell B_0+R(t).
\]
For a representative \(\zeta=Z(s)u+k\), where \(u,k\in\mathbb R^2\),
\(\operatorname{Im}s=-\ell/(2\pi)\) gives the exact identity
\begin{equation}
 \begin{pmatrix}\log|x_1|\\\log|x_2|\end{pmatrix}
 =A(t)u.
 \label{wb:cusp-null:eq:moduli}
\end{equation}
The quotient value is \(u\bmod\mathbb Z^2\).
Since \(B_0\) is orthogonal, for all \(v\in\mathbb R^2\),
\begin{equation}
 \|A(t)v\|\geq (|\ell|-\|R(t)\|)\|v\|.
 \label{wb:cusp-null:eq:invertibility}
\end{equation}
The boundedness of \(R\) therefore makes \(A(t)\) invertible on a
sufficiently small punctured disc.

\item For any point of the fibre quotient over \(\overline u^0\), subtracting
the appropriate original period, equivalently applying
\(\Psi_{-\lambda}\), gives a representative with the exact lift
\(u=u^0\). This representative is unique: another residual translate
has logarithmic moduli differing by \(A(t)\lambda\), which vanishes
only for \(\lambda=0\). Conversely every dense-torus point satisfying
\(\log|x|=A(t)u^0\) has this quotient value, since the remaining real
parts of \(\zeta\) are exactly the kernel coordinates \(k\).
For each fixed \(t\), varying \(k\bmod\mathbb Z^2\) gives all independent
phases of \(x_1,x_2\).

\item The following change of coordinates uses the full original period units:
\begin{align}
 h_j^0(t)&=\exp\bigl(-2\pi i(C(t)u^0)_j\bigr),\qquad j=1,2,\\
 \mathcal F(z_0,z_1,z_2)
 &=(w_0,w_1,w_2)
 =\left(\frac{z_0}{h_1^0(t)h_2^0(t)},\,
 h_1^0(t)z_1,\,h_2^0(t)z_2\right).
 \label{wb:cusp-null:eq:straightening}
\end{align}
The functions \(h_j^0\) are nowhere-zero holomorphic units, including
at \(t=0\). The product \(w_0w_1w_2\) is exactly \(t\).
The inverse of \(\mathcal F\) is
\[
 (w_0,w_1,w_2)\longmapsto
 \left(w_0h_1^0(t)h_2^0(t),\,
 \frac{w_1}{h_1^0(t)},\,\frac{w_2}{h_2^0(t)}\right),
 \qquad t=w_0w_1w_2.
\]
Thus this is a biholomorphism of the original chart restricted to the
cusp disc, not a replacement of \(C(t)\).

\item For \(r_0\in(0,\epsilon)\), the entire inverse image of
\(y_0=[r_0,c(\overline u^0)]\) becomes exactly
\begin{equation}
 T_{r_0}=
 \{|w_0|=\rho_0,\ |w_1|=\rho_1,\ |w_2|=\rho_2\},\qquad
 (\rho_0,\rho_1,\rho_2)=
 (r_0^{1/2},r_0^{1/4},r_0^{1/4}).
 \label{wb:cusp-null:eq:torus}
\end{equation}
Indeed, \(\log|h^0(t)|=2\pi\operatorname{Im}C(t)u^0=-R(t)u^0\),
so \eqref{wb:cusp-null:eq:moduli} gives \(\log|w_j|=\ell(B_0u^0)_j\) for
\(j=1,2\); the product identity gives the exponent \(1/2\) for
\(w_0\). Conversely these moduli recover \(u=u^0\) by invertibility
of \(A(t)\). The angular map is explicitly
\phantomsection\label{qa:III:layout:1}
\[
\begin{gathered}
 (\arg t,\arg w_1,\arg w_2)\longleftrightarrow
 (\arg w_0,\arg w_1,\arg w_2),\\
 \arg w_0=\arg t-\arg w_1-\arg w_2\pmod{2\pi},
\end{gathered}
\]
with inverse \(\arg t=\arg w_0+\arg w_1+\arg w_2\).
Its integral matrix has determinant one. This includes the complete
angular mapping torus and all kernel phases, with their original signs.

\item Let \(o\) be the origin of \eqref{wb:cusp-null:eq:chart}. The original free properly
discontinuous action supplies a neighborhood \(V\) of \(o\) on which
the quotient projection is injective. To see the local injectivity
directly, first choose a relatively compact neighborhood; proper
discontinuity leaves only finitely many translates that can intersect it.
None of those nonidentity translates fixes \(o\). Shrinking finitely
many times separates the neighborhood from all of them. The quotient
map is a local biholomorphism, so its restriction to \(V\) is a chart
of the original source.
Choose a polydisc about \(0\) whose closure lies in \(\mathcal F(V)\).
For sufficiently small \(r_0\), both \(T_{r_0}\) and
\begin{equation}
 \{|w_0|\leq\rho_0,\ |w_1|=\rho_1,\ |w_2|=\rho_2\}
 \label{wb:cusp-null:eq:spatial-bound}
\end{equation}
lie in this polydisc. This proves that the whole inverse image embeds
in one chart and supplies space for the bounding manifold below.

\item \paragraph{2. The actual normal framing, including its angle dependence.}
Put \(L_j=\log|w_j|\), and use the unique real lift of the quotient
coordinate near \(u^0\). From the exact change of coordinates,
\begin{equation}
 \begin{pmatrix}L_1\\L_2\end{pmatrix}
 =\ell B_0u^0+A(t)(u-u^0),\qquad
 L_0=\ell-L_1-L_2.
 \label{wb:cusp-null:eq:nearby-moduli}
\end{equation}
Differentiating on \(u=u^0\) annihilates the term
\((dA(t))(u-u^0)\). Hence no derivative of \(C(t)\) has been omitted:
the exact differential there is
\begin{equation}
 \begin{pmatrix}dL_0\\dL_1\\dL_2\end{pmatrix}
 =J(t)\begin{pmatrix}d\ell\\du_1\\du_2\end{pmatrix},\quad
 J(t)=
 \begin{pmatrix}
 1/2&-A_{11}-A_{21}&-A_{12}-A_{22}\\
 1/4&A_{11}&A_{12}\\
 1/4&A_{21}&A_{22}
 \end{pmatrix}.
 \label{wb:cusp-null:eq:J}
\end{equation}
Adding the second and third rows to the first gives \((1,0,0)\),
and expansion in that row proves \(\det J(t)=\det A(t)\).

\item Let \(D=Dc_{\overline u^0}\) be the actual two-by-two derivative in
the ordered open-cell target coordinates \((c_1,c_2)\). It is a fixed
invertible matrix; all its entries and signs are retained. Let
\[
 H=\begin{pmatrix}r_0&0&0\\0&D_{11}&D_{12}\\0&D_{21}&D_{22}\end{pmatrix}.
\]
Let \(n_j\) be the real radial vector \(w_j\) in the \(j\)-th complex
coordinate, with other components zero. Then \(dL_i(n_j)=\delta_{ij}\).
Writing \(N=(n_0,n_1,n_2)\) as the matrix with these vector columns, the actual
\(dg\)-dual normal framing is, in the normal quotient bundle,
\begin{equation}
 \mathfrak n(t)=N J(t)H^{-1}.
 \label{wb:cusp-null:eq:actual-frame}
\end{equation}
Indeed, \(dg(N)=HJ(t)^{-1}\), so multiplication by \(J(t)H^{-1}\)
gives the identity. In the Euclidean metric of this chart the radial
vectors are perpendicular to the product torus; a different source
metric changes their representatives by the canonical projection to its
normal bundle and preserves this normal-quotient framing.

\item Put \(\ell_0=\log r_0\). Decrease \(r_0\) further, if necessary, so
that \(|\ell_0|>2\sup_{|t|\leq r_0}\|R(t)\|\). For
\(0\leq\sigma\leq1\), define only a homotopy of frames by
\[
 A_\sigma(t)=\ell_0B_0+\sigma R(t),
\]
and let \(J_\sigma(t)\) be \eqref{wb:cusp-null:eq:J} with \(A_\sigma\) in place
of \(A\). For every \(v\in\mathbb R^2\),
\[
 \|A_\sigma(t)v\|\geq
 (|\ell_0|-\|R(t)\|)\|v\|>0\quad(v\ne0).
\]
Thus all matrices are invertible. At \(\sigma=0\),
\begin{equation}
 J_0=
 \begin{pmatrix}
 1/2&\ell_0&-\ell_0\\
 1/4&0&\ell_0\\
 1/4&-\ell_0&0
 \end{pmatrix},\qquad \det J_0=\ell_0^2>0.
 \label{wb:cusp-null:eq:J0}
\end{equation}
Continuity also proves \(\det J_\sigma(t)>0\) for every \(\sigma,t\).
The homotopy \(N J_\sigma(t)H^{-1}\), run from \(\sigma=1\) to
\(\sigma=0\), compares the actual frame with the exact constant-matrix
frame \(N P\), where
\begin{equation}
 P=J_0H^{-1}.
 \label{wb:cusp-null:eq:P}
\end{equation}
The original embedded torus and original functions \(C(t)\) have not
changed. A framing homotopy is a framed cylinder: choose a smooth
parameter function constant near its endpoints and use the displayed
frames on \(T_{r_0}\times[0,1]\).

\item \paragraph{3. A collared framed nullbordism with the exact constant matrix.}
Use a new time coordinate \(v\in[0,1]\). Define
\begin{equation}
 q(v)=
 \begin{cases}
 0,&0\leq v\leq1/4,\\
 \rho_0^2\exp\bigl(4-1/(v-1/4)\bigr),&1/4<v\leq1.
 \end{cases}
 \label{wb:cusp-null:eq:q}
\end{equation}
This function is smooth: all right derivatives of
\(\exp(-1/(v-1/4))\) at \(1/4\) vanish. It is strictly increasing
for \(v>1/4\), satisfies \(q(1/2)=\rho_0^2\), and
\(q'(1/2)=16\rho_0^2>0\).
In the chart times \([0,1]\), put
\begin{equation}
 \mathcal W=
 \{|w_0|^2+q(v)=\rho_0^2,\quad
       |w_1|^2=\rho_1^2,\quad|w_2|^2=\rho_2^2\}.
 \label{wb:cusp-null:eq:bordism}
\end{equation}
No point has \(v>1/2\). The three defining differentials are independent:
the latter two have nonzero gradients in their separate complex factors;
the first has nonzero \(w_0\)-gradient whenever \(w_0\ne0\), and at
\(w_0=0\) it has the nonzero \(v\)-derivative \(q'(1/2)\).
Thus \(\mathcal W\) is a smooth compact four-manifold. Its only
boundary is \(T_{r_0}\times\{0\}\), and on \(0\leq v\leq1/4\)
it is the exact product collar \(T_{r_0}\times[0,1/4]\).
The north locus \(w_0=0,v=1/2\) is an interior cap: the implicit
function theorem expresses \(v\) smoothly in terms of
\((\operatorname{Re}w_0,\operatorname{Im}w_0)\) there. The first
factor is a disk obtained by rotating its radial profile with this smooth
cap and one boundary circle. Therefore
\(\mathcal W\cong D^2\times T^2\).
Its spatial coordinates lie in \eqref{wb:cusp-null:eq:spatial-bound}, so the complete
manifold lies in the injective original chart.

\item In the Euclidean chart product metric define
\[
 V_0=(w_0,0,0,q'(v)/2),\qquad
 V_1=(0,w_1,0,0),\qquad V_2=(0,0,w_2,0).
\]
These are one half of the gradients of the three defining equations
in \eqref{wb:cusp-null:eq:bordism}. They are mutually orthogonal and never zero.
For example, if \(F_0=|w_0|^2+q(v)-\rho_0^2\), then
\[
 dF_0(V_0)=2|w_0|^2+\frac12q'(v)^2>0
\]
at every point of \(\mathcal W\). On the entire collar \(q'=0\), so
\((V_0,V_1,V_2)=(N,0)\) exactly. Consequently
\begin{equation}
 (V_0,V_1,V_2)P
 \label{wb:cusp-null:eq:extended-frame}
\end{equation}
is a framing of the normal three-plane bundle of \(\mathcal W\),
and its restriction to the collar is exactly the framing \(NP\) from
\eqref{wb:cusp-null:eq:P}, with every original target-coordinate sign retained.
Concatenate the framing cylinder of the preceding paragraph with this
nullbordism, rescaling the two interval lengths. Their framings agree
exactly on the common product collar, so the concatenation is a smooth
framed nullbordism of the actual inverse image.

\item \paragraph{4. The based map represented by this framed inverse image.}
Choose the injective chart to exclude \(x_0\) if its center is not
\(x_0\). If its center is \(x_0\), the constructed nullbordism still
excludes \(x_0\), since \(|w_1|=\rho_1>0\) and
\(|w_2|=\rho_2>0\) everywhere on it. Compactness therefore gives
tubular neighborhoods for the framed cylinder and cap disjoint from
\(\{x_0\}\times[0,1]\).

\item The original continuous quotient map is smooth and a submersion near
its entire inverse image of \(y_0\). First make it constant in a
neighborhood of \(x_0\), by a based homotopy that changes nothing near
that inverse image. Explicitly, take a target chart
\(\kappa:O\to\mathbb R^3\) with \(\kappa(*)=0\), whose image is a ball,
and \(y_0\notin O\). By continuity choose a neighborhood \(U\) of
\(x_0\) with \(g(U)\subset O\), and a smooth bump function
\(\beta\) supported in \(U\) and equal to one near \(x_0\).
For \(x\in U\), replace \(g(x)\) during this homotopy by
\(\kappa^{-1}((1-a\beta(x))\kappa(g(x)))\), \(0\leq a\leq1\),
and leave \(g\) fixed off \(U\). Convexity of the target chart image
makes the formula well-defined, the vanishing of \(\beta\) near the
boundary of its support proves continuity of the pasted formula, and
\(\kappa(g(x_0))=0\) makes the homotopy based. It never introduces
an inverse image of \(y_0\). Its endpoint is smooth near \(x_0\) and
agrees with \(g\) near the entire old inverse image.

\item Denote the endpoint of the preceding homotopy by \(g_b\). Choose one
closed tubular neighborhood \(U_1\) of the inverse image, contained in
the smooth submersion region and with the entire inverse image contained
in \(\operatorname{int}U_1\). The compact set
\(X\setminus\operatorname{int}U_1\) has image under \(g_b\) missing
\(y_0\). In a fixed target metric put
\[
 d_0=\min_{x\in X\setminus\operatorname{int}U_1}
       d(g_b(x),y_0)>0.
\]
Relative smooth approximation, fixed on the whole closed set \(U_1\)
and on a neighborhood of \(x_0\), gives a smooth based representative
\(g_{\mathrm{sm}}\) whose error on
\(X\setminus\operatorname{int}U_1\) is less than \(d_0/2\).
One explicit relative approximation is obtained by embedding $S^3$ in $\mathbb R^4$, smoothing the coordinate functions by convolution in finitely many source charts and a partition of unity, and interpolating this approximation with $g_b$ using a smooth cutoff equal to one on the protected neighborhoods and supported where $g_b$ is smooth. Radial projection $v\mapsto v/\|v\|$ returns the result to $S^3$. A sufficiently small uniform error keeps $v$ nonzero, fixes the protected values exactly, and has the required target error; short target geodesics give the relative homotopy. It therefore has the same inverse image and differential on \(U_1\),
and \(d(g_{\mathrm{sm}}(x),y_0)>d_0/2\) on its complement. Thus there
is no uncontrolled intervening region and no additional inverse image.
Make the approximation also closer than a fixed target injectivity
radius; its short geodesics give the based homotopy.

\item The relative Pontryagin construction now applies to the explicit
framed nullbordism. It produces a based homotopy to a map missing
\(y_0\); stereographic contraction of \(S^3\setminus\{y_0\}\cong\mathbb R^3\)
to the image of \(*\) is based and contracts that last map.
For completeness, the dependence on the framed inverse image in this
argument is precisely the local product-neighborhood construction and
the interpolation lemma of \cite[\S7, Theorem B and Lemma 4,
pp.~43, 48--50]{CuspNullMilnor1965}. Their relative versions fix the source
basepoint because all tubular modifications are away from it and
stereographic interpolation between equal values keeps that value
fixed. A framed cylinder realizes a framing homotopy; the empty
Pontryagin manifold realizes a map missing the chosen regular value.
This proves the claimed based null-homotopy of the exact map \(g\).
\end{enumerate}
\end{proof}

\begin{corollary}\label{wb:cusp_s12_carrier_homotopy:L567}
\StatementRecord{wb:cusp_s12_carrier_homotopy:L567}{proved}{Historical events: GPT-6 Astra (OpenAI), ChatGPT Pro mode, effort Pro, S6\_09\_PAPERS, 17 November 2026: targeted proof check. Also checked in this writing. Pre-continuation writing events: GPT-6 Astra (OpenAI), ChatGPT Pro mode, effort Pro, S6\_10\_PAPERS, 18 November 2026: targeted proof check. Continuation metadata events: GPT-6 Astra (OpenAI), ChatGPT Pro mode, effort Pro, S6\_10\_PAPERS continuation, 19 November 2026: used and rechecked in this run; status-vocabulary reconciliation; metadata reconciliation only; no new mathematical check. \emph{Pre-continuation recorded proof scope.} GPT-6 Astra (OpenAI, ChatGPT Pro mode), S6\_09\_PAPERS, effort Pro, 2026-11-17; also checked in this writing by GPT-6 Astra (OpenAI), ChatGPT Pro mode, effort Pro, 18 November 2026}{the hypotheses, notation and definitions stated here; the explicit constructions and calculations in the proof}

Under the original fixed based identification \(X\cong S^6\), the
one-copy coefficient is
\[
 [g]=0\in\pi_6(S^3)\cong\mathbb Z/12,\qquad a=0\pmod{12},\qquad
 [E^6g]=0\in\pi_{12}(S^9)\cong\mathbb Z/24.
\]
The previously defined homomorphism
\(\Theta_g:\mathbb Z/12\to\pi_6(S^3)\), \(\Theta_g([n])=n[g]\), and
its composite with the marked peripheral character are zero.
\end{corollary}


\begin{proof}
\begin{enumerate}[label=\arabic*.,leftmargin=*]
\item The theorem states that the actual based map is null-homotopic.
Precomposition with the fixed based identification preserves a based
null-homotopy. Suspension preserves it as well. Thus the unique
coefficient of \([g]\) in any generator of \(\mathbb Z/12\) is zero.
For every \([n]\), \(\Theta_g([n])=n\cdot0=0\), and every
composite with this homomorphism is zero.
\end{enumerate}
\end{proof}

\paragraph{Exact scope.}
The proof computes this sphere-valued postcomposition. It does not
compute the homotopy class of the earlier conical-target map
\(q_2:X\to\operatorname{Susp}(S^1\times\overline T^2)\), and it
does not assert a null-homotopy for that different target. Point-set
surjectivity of \(g\) and all of its stated metric restrictions remain
properties of the original map: a homotopy-class computation does not
alter the defining formula \eqref{wb:cusp-null:eq:g}.
The retained-angle companion is proved separately in
Theorem~\ref{maps:src:B.3.1}: $p_3q_2:X\to S^4$ represents
$\eta_4\eta_5\ne0$, so $q_2$ is nonnull. Its Hopf smash
represents $12\nu_6\ne0$ in $\pi_9(S^6)$. These positive
computations use the actual four-framed inverse torus of that different
postcomposition; they do not reverse or contradict the preceding
null-homotopy through $\widehat p_2$.




\subsection{Fixed-target metric realization and its sharp dimensional scope}
\label{wb:sphere-metric-limit-note}

This note concerns the codomains of the three based maps
\[
 q_2^{\wedge 2}:X^{\wedge 2}\cong S^{12}
   \longrightarrow Z_{12}:=\mathcal A^{\wedge 2},
\]
\[
 q_{1,\mathrm{cusp}}^{\wedge 4}:X^{\wedge 4}\cong S^{24}
   \longrightarrow Z_{24}^{\mathcal B}:=\mathcal B^{\wedge 4},
 \qquad
 q_1^{\wedge 4}:X^{\wedge 4}\cong S^{24}
   \longrightarrow Z_{24}^{\mathcal C}:=\mathcal C^{\wedge 4}.
\]
In particular, \(S^{12}\) and \(S^{24}\) are the sources, not the spaces
in which the metric limits occur.  We construct one fixed compatible
length metric on each codomain which realizes the native
four-dimensional flat limits on the target slices.  The realization is
sharp: after the same rescaling, the full ambient limits have dimensions
\(8\) and \(12\), and the collapsed smash base point goes to infinity.

Put
\[
 r_T=e^{-2\pi T},\qquad
 J=\begin{pmatrix}0&1\\-1&0\end{pmatrix},\qquad
 Y_T=TJ+E_T,
\]
where \(E_T\) is the bounded matrix in the radial cusp expansion.  Choose
and fix \(T_0\) and \(C>0\) so that
\(\|E_T\|_{\mathrm{op}}\leq C\) for every \(T\geq T_0\), in the fixed
Euclidean coordinates used below.  This notation avoids confusing the
error matrix with the quotient circle \(C_T\) used below.  On the
fixed torus
\[
 \overline T^2=(\overline\Lambda\otimes\mathbb R)/\overline\Lambda
              \cong\mathbb R^2/\mathbb Z^2
\]
let
\[
 h_T:=\iota_T^*g_{Q_T}=dx^{\mathsf T}Y_T^{\mathsf T}Y_Tdx.
\]
Thus
\begin{equation}
 T^{-2}h_T\longrightarrow h_\infty:=dx^{\mathsf T}dx
 \label{wb:sphere-metric-normalized-Q-limit}
\end{equation}
in every \(C^k\) norm in the fixed flat coordinates.  Indeed, the cusp
coefficients are holomorphic in \(t=e^{-2\pi T}\), so their \(T\)-derivatives
are exponentially small, while \(T^{-1}Y_T=J+T^{-1}E_T\).

Use normalized angular coordinates
\[
 u,\theta\in\mathbb R/\mathbb Z;
\]
\(u\) is the argument of the base coordinate and \(\theta\) is the value of
the invariant character.  Increase this same \(T_0\) once more, if
necessary, so that all radial cusp formulas used below hold for
\(T\geq T_0\) and \(e^{-2\pi T_0}<\epsilon\), and put
\[
 r_0=e^{-2\pi T_0}.
\]
For \(0<r\leq r_0\), write
\[
 T(r):=-\frac{1}{2\pi}\log r.
\]

\paragraph{Target-slice notation.}
We use the embeddings
\(j_T,k_T^{\mathcal B},k_T^{\mathcal C}\) defined in
\cref{wb:cusp-cone-target-slice-embeddings}; the normalized angle \(u=0\) used
below is the same positive radial base angle denoted by \(1\in S^1\) there.

\begin{theorem}[Topology-forced collapse for every fixed compatible metric]
\label{wb:sphere-metric-uniform-collapse}
\StatementRecord{wb:sphere-metric-uniform-collapse}{proved}{Historical events: GPT-6 Astra (OpenAI), ChatGPT Pro mode, effort Pro, S6\_09\_PAPERS, 17 November 2026: targeted proof check. Also checked in this writing. Pre-continuation writing events: GPT-6 Astra (OpenAI), ChatGPT Pro mode, effort Pro, S6\_10\_PAPERS, 18 November 2026: targeted proof check. Continuation metadata events: GPT-6 Astra (OpenAI), ChatGPT Pro mode, effort Pro, S6\_10\_PAPERS continuation, 19 November 2026: used and rechecked in this run; status-vocabulary reconciliation; metadata reconciliation only; no new mathematical check. \emph{Pre-continuation recorded proof scope.} /root, effort Pro, 2026-11-17; also checked in this writing by GPT-6 Astra (OpenAI), ChatGPT Pro mode, effort Pro, 18 November 2026}{the hypotheses, notation and definitions stated here; the explicit constructions and calculations in the proof}

Let \(D_{12},D_{24}^{\mathcal B},D_{24}^{\mathcal C}\) be arbitrary fixed
metrics compatible with the respective compact target topologies, and let
\(*\) denote the corresponding smash base point.  Then
\begin{align}
 d_H^{D_{12}}\bigl(j_T(Q_T^2),\{*\}\bigr)&\longrightarrow0,\\
 d_H^{D_{24}^{\mathcal B}}
  \bigl(k_T^{\mathcal B}(C_T^4),\{*\}\bigr)&\longrightarrow0,\\
 d_H^{D_{24}^{\mathcal C}}
  \bigl(k_T^{\mathcal C}(C_T^4),\{*\}\bigr)&\longrightarrow0.
 \label{wb:sphere-metric-Hausdorff-collapse}
\end{align}
Consequently, with the metrics restricted from any one fixed ambient
target, all three target slices converge in Gromov--Hausdorff topology to
a point, not to \(\mathbb R^4\).
\end{theorem}



\begin{proof}
\begin{enumerate}[label=\arabic*.,leftmargin=*]
\item We prove the first assertion; the other two are identical.  Let
\[
 G:[0,r_0]^2\times(T^3)^2\longrightarrow Z_{12}
\]
be the product of the two inner-cone quotient maps followed by the smash
quotient.  It is continuous and
\(G(0,0,z_1,z_2)=*\) for every \((z_1,z_2)\in(T^3)^2\).  Since the domain
is compact, uniform continuity with respect to \(D_{12}\) gives
\[
 \sup_{(z_1,z_2)\in(T^3)^2}
 D_{12}\bigl(G(r,r,z_1,z_2),*\bigr)\longrightarrow0.
\]
The set \(j_T(Q_T^2)\) is contained in such a cross-section, which proves
the first Hausdorff convergence.  Replacing \(T^3\) by \(T^2\) and taking
four factors proves the other two.  Hausdorff convergence to a singleton
implies Gromov--Hausdorff convergence to that singleton.
\end{enumerate}
\end{proof}

The preceding theorem is unrescaled.  It does not prevent a blow-up of one
fixed metric.  The following construction gives the strongest natural
positive statement.

\begin{theorem}[One fixed cusp-adapted metric realizes the vertical
\(\mathbb R^4\) limits]
\label{wb:sphere-metric-fixed-realization}
\StatementRecord{wb:sphere-metric-fixed-realization}{proved}{Historical events: GPT-6 Astra (OpenAI), ChatGPT Pro mode, effort Pro, S6\_09\_PAPERS, 17 November 2026: targeted proof check. Also checked in this writing. Pre-continuation writing events: GPT-6 Astra (OpenAI), ChatGPT Pro mode, effort Pro, S6\_10\_PAPERS, 18 November 2026: targeted proof check. Continuation metadata events: GPT-6 Astra (OpenAI), ChatGPT Pro mode, effort Pro, S6\_10\_PAPERS continuation, 19 November 2026: used and rechecked in this run; status-vocabulary reconciliation; metadata reconciliation only; no new mathematical check. \emph{Pre-continuation recorded proof scope.} /root, effort Pro, 2026-11-17; also checked in this writing by GPT-6 Astra (OpenAI), ChatGPT Pro mode, effort Pro, 18 November 2026}{the hypotheses, notation and definitions stated here; \cref{wb:sphere-metric-normalized-Q-limit}}

There are fixed compatible length metrics
\[
 D_{12}\text{ on }Z_{12},\qquad
 D_{24}^{\mathcal B}\text{ on }Z_{24}^{\mathcal B},\qquad
 D_{24}^{\mathcal C}\text{ on }Z_{24}^{\mathcal C}
\]
with the following properties.

On the inner cone of \(\mathcal A\), the metric is induced by
\begin{equation}
 g_{\mathcal A}
 =dr^2+r^2du^2+\frac{r^2}{T(r)^2}h_{T(r)}.
 \label{wb:sphere-metric-A-fixed-metric}
\end{equation}
On the cusp cone of either \(\mathcal B\) or \(\mathcal C\), it is induced
by
\begin{equation}
 g_{\mathcal D}=dr^2+r^2(du^2+d\theta^2),
 \qquad\mathcal D\in\{\mathcal B,\mathcal C\}.
 \label{wb:sphere-metric-BC-fixed-metric}
\end{equation}
The metrics on the smash powers are the metric quotients of the product
length metrics after the closed fat wedge is collapsed.  Write
\(g_{Z_{12}}\) and \(g_{Z_{24}^{\mathcal D}}\) for the resulting smooth
Riemannian tensors on the open regular strata.

Let
\[
 p_T^{12}:=j_T(0,0),\qquad
 p_T^{24,\mathcal D}:=k_T^{\mathcal D}(0,0,0,0).
\]
Put
\[
 \Sigma_T^{12}:=j_T(Q_T^2),\qquad
 \Sigma_T^{24,\mathcal D}:=k_T^{\mathcal D}(C_T^4).
\]
Let \(d_{T,12}^{\mathrm{slice}}\) and
\(d_{T,24,\mathcal D}^{\mathrm{slice}}\) be the intrinsic Riemannian path
metrics of the first fundamental forms on these embedded slices.  These
must be distinguished from the restricted ambient distance functions
\[
 d_{T,12}^{\mathrm{amb}}
 :=D_{12}|_{\Sigma_T^{12}\times\Sigma_T^{12}},\qquad
 d_{T,24,\mathcal D}^{\mathrm{amb}}
 :=D_{24}^{\mathcal D}|_{\Sigma_T^{24,\mathcal D}
                         \times\Sigma_T^{24,\mathcal D}}.
\]
In general the latter can be smaller, because an ambient path is allowed
to leave the slice.

Then the induced first fundamental forms on the target slices satisfy the
exact identities
\begin{equation}
 \left(\frac{T}{r_T}\right)^2
  j_T^*g_{Z_{12}}
 =g_{Q_T}\oplus g_{Q_T},
 \label{wb:sphere-metric-A-slice-exact}
\end{equation}
and
\begin{equation}
 \left(\frac{\ell_T}{r_T}\right)^2
  (k_T^{\mathcal D})^*g_{Z_{24}^{\mathcal D}}
 =g_{C_T}^{\oplus4}.
 \label{wb:sphere-metric-BC-slice-exact}
\end{equation}
Consequently,
\begin{align}
 \left(\Sigma_T^{12},\frac{T}{r_T}d_{T,12}^{\mathrm{slice}},
       p_T^{12}\right)
 &\longrightarrow(\mathbb R^4,0),\\
 \left(\Sigma_T^{24,\mathcal D},
       \frac{\ell_T}{r_T}d_{T,24,\mathcal D}^{\mathrm{slice}},
       p_T^{24,\mathcal D}\right)
 &\longrightarrow(\mathbb R^4,0)
 \label{wb:sphere-metric-slice-fixed-blowups}
\end{align}
in pointed local-isometric and pointed Gromov--Hausdorff topology, and
smoothly on the exhausting Euclidean slice balls.

The full ambient pointed blow-ups are instead
\begin{align}
 \left(Z_{12},\frac{T}{r_T}D_{12},p_T^{12}\right)
 &\longrightarrow(\mathbb R^8,0),\\
 \left(Z_{24}^{\mathcal D},
       \frac{\ell_T}{r_T}D_{24}^{\mathcal D},
       p_T^{24,\mathcal D}\right)
 &\longrightarrow(\mathbb R^{12},0).
 \label{wb:sphere-metric-ambient-fixed-blowups}
\end{align}
as pointed metric spaces in pointed Gromov--Hausdorff topology.  More
precisely, for every fixed radius the corresponding rescaled target ball is
eventually contained in an open regular-stratum coordinate neighborhood,
and the pulled-back Riemannian tensors converge there in every \(C^k\) norm.
Thus smooth convergence is asserted on an exhaustion by regular-stratum
neighborhoods, not globally on the singular compact targets.  Under these
limits, the slice \(\mathbb R^4\) is respectively the coordinate subspace
obtained by fixing the two radial and two base-angle coordinates, or by
fixing the four radial and four base-angle coordinates.
\end{theorem}



\begin{proof}
\begin{enumerate}[label=\arabic*.,leftmargin=*]
\item Equation \eqref{wb:sphere-metric-normalized-Q-limit} shows that the metrics on
the compact cross-sections of \eqref{wb:sphere-metric-A-fixed-metric} remain
uniformly bounded and converge to a positive definite flat metric.
Therefore the cross-section diameter is \(O(r)\).  Moreover the radial
coordinate is \(1\)-Lipschitz and a radial segment from level \(r\) to the
vertex has length \(r\).  The length completion consequently adds exactly
one point at \(r=0\), and its topology is the prescribed cone topology.
The same assertions are immediate for
\eqref{wb:sphere-metric-BC-fixed-metric}.  Interpolation on a closed annulus
away from the vertex extends these metrics to fixed compatible length
metrics on the compact spaces \(\mathcal A,\mathcal B,\mathcal C\).

\item If a compact metric space \((E,d)\) and a closed subset \(W\subset E\) are
given, the metric quotient which collapses \(W\) to one point is
\[
 d_{E/W}([x],[y])
 =\min\{d(x,y),d(x,W)+d(y,W)\}.
\]
Applying this to the fat wedges in the product length spaces gives the
claimed fixed compatible metrics on the smash targets.

\item On the \(r_T\)-slice of \(\mathcal A\), with \(u=0\), the first fundamental
form is
\[
 \frac{r_T^2}{T^2}h_T.
\]
The definition \(h_T=\iota_T^*g_{Q_T}\) proves
\eqref{wb:sphere-metric-A-slice-exact}.  On the \(r_T\)-slice of
\(\mathcal B\) or \(\mathcal C\), with \(u=0\), the form is
\(r_T^2d\theta^2\).  Under
\(\vartheta_T:[\theta]\mapsto[\ell_T\theta]\), the native quotient-circle
metric is \(\ell_T^2d\theta^2\).  This proves
\eqref{wb:sphere-metric-BC-slice-exact}.  The native systole estimates
\[
 \operatorname{sys}(Q_T^2)\geq T-C,
 \qquad
 \operatorname{sys}(C_T^4)=\ell_T
\]
then prove \eqref{wb:sphere-metric-slice-fixed-blowups}: every prescribed
radius ball is eventually exactly a Euclidean four-ball.

\item It remains to compute the ambient limits.  In one \(\mathcal A\)-factor,
put
\[
 \lambda_T=\frac{T}{r_T},\qquad
 s=\lambda_T(r-r_T),\qquad
 \eta=Tu,\qquad z=Y_Tx
\]
on the universal cover of the torus directions.  On every fixed bounded
set in \((s,\eta,z)\), one has
\[
 \frac r{r_T}=1+\frac sT,
 \qquad
 T(r)=T+O(T^{-1}),
 \qquad
 Y_{T(r)}Y_T^{-1}=I+O(T^{-2}),
\]
with the same estimates after any fixed number of derivatives.  Thus
\(\lambda_T^2g_{\mathcal A}\) converges in these coordinates to
\[
 ds^2+d\eta^2+dz_1^2+dz_2^2.
\]
The inward radial boundary is at distance \(T\), the \(u\)-period is \(T\),
and the \(z\)-period lattice is \(Y_T\mathbb Z^2\), whose systole is at
least \(T-C\).  The outward edge of the fixed cusp chart is also at scaled
distance
\[
 \frac{T}{r_T}(r_0-r_T)
 =T\left(\frac{r_0}{r_T}-1\right)\longrightarrow\infty.
\]
Thus the inward vertex, all periods, and the outward edge recede to
infinity.  A single factor converges on exhausting regular-stratum
coordinate neighborhoods to \(\mathbb R^4\), and two factors converge to
\(\mathbb R^8\).

\item For one \(\mathcal B\)- or \(\mathcal C\)-factor, put
\[
 \lambda_T=\frac{\ell_T}{r_T},\qquad
 s=\lambda_T(r-r_T),\qquad
 \eta=\ell_Tu,\qquad \zeta=\ell_T\theta.
\]
Then
\[
 \lambda_T^2g_{\mathcal D}
 =ds^2+\left(\frac r{r_T}\right)^2
       (d\eta^2+d\zeta^2),
\]
which converges smoothly on bounded coordinate sets to the Euclidean metric
on \(\mathbb R^3\).  The two circle periods and the inward radial boundary
are all at scaled distance \(\ell_T\to\infty\), while the outward edge is
at scaled distance
\[
 \frac{\ell_T}{r_T}(r_0-r_T)
 =\ell_T\left(\frac{r_0}{r_T}-1\right)\longrightarrow\infty.
\]
Four factors therefore converge on exhausting regular-stratum coordinate
neighborhoods to \(\mathbb R^{12}\).

\item Finally, a level-\(r_T\) factor point has distance exactly \(r_T\) from its
cone vertex.  Indeed, the radial segment has that length.  A competing
curve contained in the inner chart has length at least \(r_T\), while one
which leaves has inner-chart portions of total length at least
\((r_0-r_T)+r_0>r_T\).  The Riemannian-\(\ell^2\)-product distance from a
synchronized point to the fat wedge is the minimum of the factor
distances, and hence is also \(r_T\).

\item For the chosen scale put
\[
 b_T=\lambda_Tr_T,\qquad a_T=\lambda_T(r_0-r_T),
\]
and let \(\sigma_T\) be the least scaled period in a factor.  Thus
\(\sigma_T\geq T-C\) in the \(\mathcal A\) calculation and
\(\sigma_T=\ell_T\) in the \(\mathcal B,\mathcal C\) calculation.  The
preceding computations give
\[
 a_T\longrightarrow\infty,\qquad
 b_T\longrightarrow\infty,\qquad
 \sigma_T\longrightarrow\infty.
\]
Fix \(R\).  A curve of rescaled length at most \(R\) from \(p_T\) meets
neither the fat wedge nor \(r=r_0\) for all large \(T\), so the
radius-\(R\) target ball lies in the open regular stratum.  If \(x,y\) lie
in that ball, a competitor through the fat wedge has length at least
\[
 (b_T-R)+(b_T-R)>2R,
\]
and a competitor which exits the inner chart and returns has length at
least
\[
 2(a_T-R)>2R.
\]
Every remaining competitor of length at most \(2R\) stays in
\(r/r_T=1+o_R(1)\).  After lifting it to the angular and torus covers, a
nonzero deck translation has length at least
\((1-o_R(1))(\sigma_T-2R)>2R\) for all large \(T\).  Since the path through
\(p_T\) has length at most \(2R\), none of these three kinds of global
shortcut computes the distance between \(x\) and \(y\).  Thus the metric
on every fixed target ball is computed entirely in one lifted
regular-stratum chart.  The \(C^\infty\) convergence of the pulled-back
tensors on these exhausting charts implies the pointed
Gromov--Hausdorff limits in
\eqref{wb:sphere-metric-ambient-fixed-blowups}.
\end{enumerate}
\end{proof}

\begin{theorem}[Complete intrinsic scalar-rescaling trichotomy]
\label{wb:sphere-metric-scaling-trichotomy}
\StatementRecord{wb:sphere-metric-scaling-trichotomy}{proved}{Historical events: GPT-6 Astra (OpenAI), ChatGPT Pro mode, effort Pro, S6\_09\_PAPERS, 17 November 2026: targeted proof check. Also checked in this writing. Pre-continuation writing events: GPT-6 Astra (OpenAI), ChatGPT Pro mode, effort Pro, S6\_10\_PAPERS, 18 November 2026: targeted proof check. Continuation metadata events: GPT-6 Astra (OpenAI), ChatGPT Pro mode, effort Pro, S6\_10\_PAPERS continuation, 19 November 2026: used and rechecked in this run; status-vocabulary reconciliation; metadata reconciliation only; no new mathematical check. \emph{Pre-continuation recorded proof scope.} /root, effort Pro, 2026-11-17; also checked in this writing by GPT-6 Astra (OpenAI), ChatGPT Pro mode, effort Pro, 18 November 2026}{the hypotheses, notation and definitions stated here; \cref{wb:sphere-metric-fixed-realization}; \cref{wb:sphere-metric-outward-shortcut-exclusion}}

For the fixed metrics of
\cref{wb:sphere-metric-fixed-realization}, let \(\lambda_T>0\) be arbitrary
and put
\[
 \beta_T:=\lambda_Tr_T.
\]
Choose one of the following triples, retaining the indicated ambient
metric:
\[
 \begin{split}
  (\Sigma_T,d_T^{\mathrm{slice}},d_T^{\mathrm{amb}},p_T;D)
  &=(\Sigma_T^{12},d_{T,12}^{\mathrm{slice}},
      d_{T,12}^{\mathrm{amb}},p_T^{12};D_{12}),\\
  (\Sigma_T,d_T^{\mathrm{slice}},d_T^{\mathrm{amb}},p_T;D)
  &=(\Sigma_T^{24,\mathcal D},
      d_{T,24,\mathcal D}^{\mathrm{slice}},
      d_{T,24,\mathcal D}^{\mathrm{amb}},
      p_T^{24,\mathcal D};D_{24}^{\mathcal D}),
      \qquad \mathcal D\in\{\mathcal B,\mathcal C\}.
 \end{split}
\]
After passing to a subsequence so that \(\beta_T\) has a limit in
\([0,\infty]\), the intrinsic pointed slice limit is
\begin{equation}
 (\Sigma_T,\lambda_Td_T^{\mathrm{slice}},p_T)
 \longrightarrow
 \begin{cases}
  \{\mathrm{pt}\},&\beta_T\longrightarrow0,\\
  ((\mathbb R/L\mathbb Z)^4,0),
      &\beta_T\longrightarrow L\in(0,\infty),\\
  (\mathbb R^4,0),&\beta_T\longrightarrow\infty.
 \end{cases}
 \label{wb:sphere-metric-scaling-trichotomy-result}
\end{equation}
The convergence in
\eqref{wb:sphere-metric-scaling-trichotomy-result} is pointed
Gromov--Hausdorff in all three cases.  In the finite positive case the flat
metrics converge smoothly under the fixed marked-torus identifications; in
the unrolling case the convergence is pointed local-isometric and smooth on
the exhausting Euclidean balls.
In the middle line the lattice is orthogonally equivalent to the square
lattice \(L\mathbb Z^4\).  At the same time,
\begin{equation}
 \lambda_TD(p_T,*)=\beta_T
 \label{wb:sphere-metric-vertex-distance}
\end{equation}
for every one of the displayed choices.

Only in the unrolling regime \(\beta_T\to\infty\) do we compare the
intrinsic metric with the restricted ambient distance.  In that regime, for
every \(R\in(0,\infty)\),
\begin{equation}
 \sup_{x,y\in
  \overline B_{\lambda_Td_T^{\mathrm{slice}}}(p_T,R)}
 \left|
  \lambda_Td_T^{\mathrm{amb}}(x,y)
  -\lambda_Td_T^{\mathrm{slice}}(x,y)
 \right|\longrightarrow0.
 \label{wb:sphere-metric-intrinsic-ambient-agreement}
\end{equation}
Moreover, there are \(\varepsilon_{T,R}\to0\) such that
\begin{equation}
 \overline B_{\lambda_Td_T^{\mathrm{slice}}}(p_T,R)
 \subset
 \overline B_{\lambda_Td_T^{\mathrm{amb}}}(p_T,R)
 \subset
 \overline B_{\lambda_Td_T^{\mathrm{slice}}}
             (p_T,(1+\varepsilon_{T,R})R).
 \label{wb:sphere-metric-intrinsic-ambient-ball-comparison}
\end{equation}
Thus the restricted ambient metric has the same pointed
\(\mathbb R^4\) limit on these fixed slice balls.  No agreement is asserted
in the regimes \(\beta_T\to0\) or \(\beta_T\to L\in(0,\infty)\), where
radial or fat-wedge shortcuts need not disappear.  Every rescaling which
unrolls the intrinsic slice sends the cone/fat-wedge base point to infinity,
so no such limit is a tangent cone based at that point.

In the same regime \(\beta_T\to\infty\), the corresponding full ambient
bubbles are
\begin{align}
 (Z_{12},\lambda_TD_{12},p_T^{12})
 &\longrightarrow(\mathbb R^8,0),\\
 (Z_{24}^{\mathcal D},\lambda_TD_{24}^{\mathcal D},
  p_T^{24,\mathcal D})
 &\longrightarrow(\mathbb R^{12},0),
 \qquad\mathcal D\in\{\mathcal B,\mathcal C\}.
 \label{wb:sphere-metric-arbitrary-unrolling-ambient-bubbles}
\end{align}
These are pointed Gromov--Hausdorff limits of the singular metric targets.
The convergence of tensors is \(C^\infty\) only under coordinate
identifications on exhausting open regular-stratum neighborhoods; no
global smooth convergence of the singular targets is claimed.
\end{theorem}



\begin{proof}
\begin{enumerate}[label=\arabic*.,leftmargin=*]
\item For the \(Q_T^2\) route, the rescaled period matrix in each two-dimensional
factor is
\[
 \frac{\lambda_Tr_T}{T}Y_T
 =\beta_T\left(J+\frac{E_T}{T}\right).
\]
If \(\beta_T\to L\in(0,\infty)\), this lattice converges to
\(LJ\mathbb Z^2\); two factors give the square four-torus in
\eqref{wb:sphere-metric-scaling-trichotomy-result}.  If
\(\beta_T\to0\), the lengths of a lattice basis tend to zero, so the
diameter tends to zero.  If \(\beta_T\to\infty\), the reverse triangle
inequality gives systole
\[
 \frac{\beta_T}{T}(T-C)
 =\beta_T\left(1-\frac CT\right)\longrightarrow\infty,
\]
so every fixed-radius ball is eventually Euclidean.

\item For a \(C_T^4\) slice, the four periods after rescaling have the common
length
\[
 \lambda_Tr_T=\beta_T,
\]
which gives the same three intrinsic cases directly.

\item Consider first one factor.  A curve from a level-\(r_T\) point to the cone
vertex which remains in the inner chart has length at least \(r_T\), because
the \(dr^2\)-term makes \(r\) \(1\)-Lipschitz there.  If it leaves the
chart, it has a first crossing of \(r=r_0\) and, before ending at the
vertex, a last re-entry through \(r=r_0\); its two inner-chart portions
have total length at least
\((r_0-r_T)+r_0>r_T\).  The radial segment has length \(r_T\), so the
factor distance to its based vertex is exactly \(r_T\).  The
Riemannian-\(\ell^2\)-product distance from every synchronized slice point
to the fat wedge is the minimum of these factor distances and hence is
also \(r_T\).  The metric-quotient formula in the preceding proof therefore
gives \eqref{wb:sphere-metric-vertex-distance}.

\item It remains to prove the restricted-ambient comparison, so assume
\(\beta_T\to\infty\) and fix \(R\).  Two points in
\(\overline B_{\lambda_Td_T^{\mathrm{slice}}}(p_T,R)\) have intrinsic
distance at most \(2R/\lambda_T\).  The slice path is an ambient competitor,
so
\[
 d_T^{\mathrm{amb}}(x,y)\leq d_T^{\mathrm{slice}}(x,y).
\]
The quotient shortcut through the fat wedge has rescaled length
\(2\lambda_Tr_T=2\beta_T\), and hence is longer than \(2R\) for all large
\(T\).  A path which leaves the fixed cusp chart must reach \(r=r_0\).
Its outward cost is controlled by
\begin{equation}
 \lambda_T(r_0-r_T)
 =\beta_T\left(\frac{r_0}{r_T}-1\right)\longrightarrow\infty,
 \label{wb:sphere-metric-outward-shortcut-exclusion}
\end{equation}
so no path of rescaled length at most \(2R\) can use the interpolating
annulus or the global part of a factor.

\item Consider now an ambient competitor of length at most
\(2R/\lambda_T+o(\lambda_T^{-1})\) which stays in the cusp chart.  In every
factor its radial coordinate satisfies
\[
 |r-r_T|\leq \frac{2R+o(1)}{\lambda_T},
 \qquad
 \frac r{r_T}=1+O_R(\beta_T^{-1}).
\]
On this tube, the horizontal tensors in the directions tangent to the slice
are uniformly \((1+\varepsilon_{T,R})\)-bilipschitz to their values at
\(r=r_T\), where \(\varepsilon_{T,R}\to0\).  For
\(\mathcal B\) and \(\mathcal C\) this follows directly from
\((r/r_T)^2=1+O_R(\beta_T^{-1})\).  For \(\mathcal A\), the same conclusion
follows, without changing the fixed coordinates or the matrix \(Y_T\), from
\[
 \begin{aligned}
  T(r)-T&=O_R(\beta_T^{-1}),\\
  \frac{T(r)}T&=1+O_R((\beta_TT)^{-1}),\\
  Y_{T(r)}Y_T^{-1}&=I+O_R((\beta_TT)^{-1}).
 \end{aligned}
\]
Projecting the competitor to the torus or character directions, and
discarding its radial and base-angle components, produces a path between
the same slice endpoints.  The bilipschitz estimate therefore gives
\[
 d_T^{\mathrm{slice}}(x,y)
 \leq(1+\varepsilon_{T,R})d_T^{\mathrm{amb}}(x,y).
\]
Together with the opposite inequality and
\(\lambda_Td_T^{\mathrm{slice}}(x,y)\leq2R\), this proves
\eqref{wb:sphere-metric-intrinsic-ambient-agreement}.  The first inclusion in
\eqref{wb:sphere-metric-intrinsic-ambient-ball-comparison} follows from
\(d_T^{\mathrm{amb}}\leq d_T^{\mathrm{slice}}\).  For a point \(x\) in
the middle ball, apply the same fat-wedge and outward exclusions to a
near-minimizing ambient path from \(p_T\) to \(x\), whose rescaled length is
at most \(R+o(1)\), and project that path to the slice directions.  This
gives
\[
 \lambda_Td_T^{\mathrm{slice}}(p_T,x)
 \leq(1+\varepsilon_{T,R})R
\]
after enlarging \(\varepsilon_{T,R}\) by a quantity tending to zero, and
proves the second inclusion.

\item For completeness, the full ambient bubbles at an arbitrary unrolling scale
are obtained without changing any coordinate or orientation.  In one
\(\mathcal A\)-factor use
\[
 s=\lambda_T(r-r_T),\qquad
 \eta=\beta_Tu,\qquad
 z=\frac{\beta_T}{T}Y_Tx.
\]
Writing \(A_T(r)=Y_{T(r)}Y_T^{-1}\), one obtains the exact expression
\[
 \lambda_T^2g_{\mathcal A}
 =ds^2+\left(\frac r{r_T}\right)^2d\eta^2
  +\left(\frac r{r_T}\right)^2
   \left(\frac{T}{T(r)}\right)^2
   dz^{\mathsf T}A_T(r)^{\mathsf T}A_T(r)dz.
\]
On every fixed bounded coordinate set,
\[
 \frac r{r_T}=1+O_R(\beta_T^{-1}),\qquad
 T(r)-T=O_R(\beta_T^{-1}),\qquad
 A_T(r)=I+O_R((\beta_TT)^{-1}),
\]
with the same estimates after every fixed number of derivatives.  Hence
the tensor converges smoothly to the Euclidean metric on \(\mathbb R^4\).
The \(\eta\)-period is \(\beta_T\), and the \(z\)-period lattice is
\[
 \frac{\beta_T}{T}Y_T\mathbb Z^2,
 \qquad
 \operatorname{sys}\left(
   \frac{\beta_T}{T}Y_T\mathbb Z^2\right)
 \geq\beta_T\left(1-\frac CT\right)\longrightarrow\infty.
\]
In one \(\mathcal B\)- or \(\mathcal C\)-factor use
\[
 s=\lambda_T(r-r_T),\qquad
 \eta=\beta_Tu,\qquad \zeta=\beta_T\theta.
\]
Then
\[
 \lambda_T^2g_{\mathcal D}
 =ds^2+\left(\frac r{r_T}\right)^2
       (d\eta^2+d\zeta^2),
\]
and both angular periods are \(\beta_T\to\infty\).

\item In all three cases the inward boundary and the fat wedge are at scaled
distance \(\beta_T\), while a lift starting at radius \(r_T\) incurs the
one-way radial cost in
\eqref{wb:sphere-metric-outward-shortcut-exclusion} before its first outward
crossing of \(r=r_0\).  Fix \(R\).  For large \(T\), all these quantities
and all nonzero period lengths exceed \(4R\).
If \(x,y\) lie in the ambient rescaled radius-\(R\) ball, their distances
from the fat wedge are at least \(\beta_T-R\); hence a quotient path through
the smash base point has length at least \(2(\beta_T-R)>2R\).  Every
remaining competitor lifts to the complement of the fat wedge in the
product.  If such a lift exits the inner product chart, then some factor
has a first crossing of \(r=r_0\); the radial variation before that crossing
is at least \(\lambda_T(r_0-r_T)-R\), and returning to the lift of \(y\)
incurs the same lower bound.  Its length is therefore greater than \(2R\).
Finally, a path of length at most \(2R\) remains in the tube
\(r/r_T=1+O_R(\beta_T^{-1})\); there every nonzero deck translation has
length tending to infinity, so such a path cannot wrap around an angular
or torus period.  The displayed coordinate neighborhoods therefore exhaust
the pointed balls without quotient or global shortcuts.
Their smooth tensor convergence proves the pointed Gromov--Hausdorff
limits in
\eqref{wb:sphere-metric-arbitrary-unrolling-ambient-bubbles}, with smooth
convergence confined to this regular-stratum exhaustion.
\end{enumerate}
\end{proof}

\begin{corollary}[Unit-volume normalization does not give
\(\mathbb R^4\)]
\label{wb:sphere-metric-unit-volume}
\StatementRecord{wb:sphere-metric-unit-volume}{proved}{Historical events: GPT-6 Astra (OpenAI), ChatGPT Pro mode, effort Pro, S6\_09\_PAPERS, 17 November 2026: targeted proof check. Also checked in this writing. Pre-continuation writing events: GPT-6 Astra (OpenAI), ChatGPT Pro mode, effort Pro, S6\_10\_PAPERS, 18 November 2026: targeted proof check. Continuation metadata events: GPT-6 Astra (OpenAI), ChatGPT Pro mode, effort Pro, S6\_10\_PAPERS continuation, 19 November 2026: used and rechecked in this run; status-vocabulary reconciliation; metadata reconciliation only; no new mathematical check. \emph{Pre-continuation recorded proof scope.} /root, effort Pro, 2026-11-17; also checked in this writing by GPT-6 Astra (OpenAI), ChatGPT Pro mode, effort Pro, 18 November 2026}{the hypotheses, notation and definitions stated here; the argument and qualifications immediately following}

The four-volume of \(Q_T^2\) is \((\det Y_T)^2\sim T^4\).  Its unit-volume
length factor is therefore \((\det Y_T)^{-1/2}\sim T^{-1}\), and its
unit-volume limit is a square flat four-torus, not \(\mathbb R^4\).
Likewise, the four-volume of \(C_T^4\) is \(\ell_T^4\), and its unit-volume
normalization is exactly the square unit four-torus.
\end{corollary}



\begin{proposition}[What vertex tangents and compact asymptotic cones can
mean]
\label{wb:sphere-metric-tangent-scope}
\StatementRecord{wb:sphere-metric-tangent-scope}{proved}{Historical events: GPT-6 Astra (OpenAI), ChatGPT Pro mode, effort Pro, S6\_09\_PAPERS, 17 November 2026: targeted proof check. Also checked in this writing. Pre-continuation writing events: GPT-6 Astra (OpenAI), ChatGPT Pro mode, effort Pro, S6\_10\_PAPERS, 18 November 2026: targeted proof check. Continuation metadata events: GPT-6 Astra (OpenAI), ChatGPT Pro mode, effort Pro, S6\_10\_PAPERS continuation, 19 November 2026: used and rechecked in this run; status-vocabulary reconciliation; metadata reconciliation only; no new mathematical check. \emph{Pre-continuation recorded proof scope.} /root, effort Pro, 2026-11-17; also checked in this writing by GPT-6 Astra (OpenAI), ChatGPT Pro mode, effort Pro, 18 November 2026}{the hypotheses, notation and definitions stated here; \cref{wb:sphere-metric-fixed-realization}; \cref{wb:sphere-metric-C-boundary-quotient-homeomorphism}; \cref{wb:cusp-cone-S24-route}; \cref{wb:sphere-metric-boundary-link-local-shift}}

The compact target topology alone determines no metric tangent cone.
For an exact conic length metric built from fixed piecewise-Riemannian
metrics on the links, however, the vertex links have topological and
Hausdorff dimensions \(7\) for \(Z_{12}\) and \(11\) for both
\(Z_{24}\) targets.  The resulting vertex cones have topological and
Hausdorff dimensions \(8\) and \(12\), and are not \(\mathbb R^4\).

More precisely, writing \(L_3=T^3\) and \(L_2=T^2\), one has the based
topological models
\[
 Z_{12}\cong\operatorname{Susp}(L_3*L_3),
 \qquad
 Z_{24}^{\mathcal B}
 \cong\operatorname{Susp}(L_2*L_2*L_2*L_2).
\]
At the suspension base point their conic links are respectively
\(L_3*L_3\) and \(L_2^{*4}\).  Also
\[
 (\mathcal C,c_0)
 \cong
 (\mathcal M/\partial\mathcal M,[\partial\mathcal M]),
 \qquad
 (Z_{24}^{\mathcal C},*)
 \cong
 (\mathcal M^4/\partial(\mathcal M^4),
  [\partial(\mathcal M^4)]),
\]
so an exact conic metric at the collapsed boundary has link
\(\partial(\mathcal M^4)\).  At the collapsed point of
\(Z_{24}^{\mathcal C}\), the complete integral local-homology calculation is
\begin{equation}
 H_i\bigl(Z_{24}^{\mathcal C},
          Z_{24}^{\mathcal C}\setminus\{*\};\mathbb Z\bigr)
 \cong
 \begin{cases}
  \mathbb Z^{\binom4{i-1}},&2\leq i\leq5,\\
  0,&i=6,7,\\
  \mathbb Z^{\binom4{i-8}},&8\leq i\leq12,\\
  0,&\text{otherwise}.
 \end{cases}
 \label{wb:sphere-metric-C-smash-local-homology}
\end{equation}
In particular, its local homology is not concentrated in top degree.  For
the first two targets the top local homology occurs in degree \(8\) and
\(12\), respectively; for the third target
\eqref{wb:sphere-metric-C-smash-local-homology} includes a top
\(\mathbb Z\) in degree \(12\).  None has the local homology of
\(\mathbb R^4\).

Finally, the standard asymptotic cone obtained by blowing down any fixed
compact metric target is a point.  Indeed, if \(R_n\to\infty\), then
\[
 \operatorname{diam}(Z,R_n^{-1}D)
 =R_n^{-1}\operatorname{diam}(Z,D)\longrightarrow0.
\]
Thus the limits in \cref{wb:sphere-metric-fixed-realization} are
moving-base-point blow-ups, often called bubbles.  They are neither
vertex-based tangent cones nor asymptotic cones of the compact targets.
\end{proposition}



\begin{proof}
\begin{enumerate}[label=\arabic*.,leftmargin=*]
\item For compact spaces \(V_i\) with closed subspaces \(A_i\), there is a
canonical homeomorphism
\[
 (V_1/A_1)\wedge\cdots\wedge(V_m/A_m)
 \cong
 \frac{V_1\times\cdots\times V_m}
      {\bigcup_iV_1\times\cdots\times A_i\times\cdots\times V_m}.
\]
First reverse the suspension coordinate, $r\mapsto\epsilon-r$, so that the chosen inner basepoint is carried to the outer suspension vertex. Take $V_i=\operatorname{Cone}(L_i)$ and $A_i=L_i$, its outer boundary.
The actual map from the product cones to the cone on the iterated join uses radial coordinate $\rho=\max_i r_i$ and join weights $t_i=r_i/\sum_jr_j$ away from the common vertex. Its inverse is $r_i=\rho t_i/\max_jt_j$; zero weights discard exactly the irrelevant cone-link coordinates. These maps extend continuously at the vertex, and the union of outer faces goes to $\rho=1$.
Collapsing that link gives its suspension.  This proves the first two
models and the link dimensions
\[
 \dim(T^3*T^3)=3+3+1=7,
 \qquad
 \dim((T^2)^{*4})=4\cdot2+3=11.
\]

\item We next prove the based homeomorphism for \(\mathcal C\), rather than
identifying the two presentations without a map.  Put
\(L=\partial\mathcal M\), write the attached cone as
\[
 \operatorname{Cone}(L)
 =L\times[0,1]/(L\times\{0\}),
\]
with base \(L\times\{1\}\), and choose a collar
\(\kappa:L\times[0,1]\hookrightarrow\mathcal M\) satisfying
\(\kappa(x,0)=x\).  There is a homeomorphism
\[
 \Phi:\mathcal C\setminus\{c_0\}
       \longrightarrow\operatorname{int}\mathcal M
\]
given on the punctured cone, the inner collar, and the remaining core by
\[
 \Phi([x,t])=\kappa(x,t/e)\quad(0<t\leq1),\qquad
 \Phi(\kappa(x,s))=\kappa(x,e^{s-1})\quad(0\leq s\leq1),
\]
\[
 \Phi(q)=q
 \quad\text{for}\quad
 q\in\mathcal M\setminus\kappa(L\times[0,1)).
\]
The first two formulas agree at the glued points
\([x,1]=\kappa(x,0)\), and the second and third agree at
\(\kappa(x,1)\).  Their collar-coordinate images are respectively
\(L\times(0,e^{-1}]\) and \(L\times[e^{-1},1]\); hence the displayed
piecewise map and its piecewise inverse are continuous.

\item The compact Hausdorff spaces \(\mathcal C\) and
\(\mathcal M/L\) are both one-point compactifications of
\(\operatorname{int}\mathcal M\): for \(\mathcal C\) use \(\Phi\), and for
\(\mathcal M/L\) use the quotient map on
\(\mathcal M\setminus L=\operatorname{int}\mathcal M\).  Uniqueness of the
one-point compactification therefore extends \(\Phi\) to a based
homeomorphism
\begin{equation}
 (\mathcal C,c_0)
 \xrightarrow{\ \cong\ }
 (\mathcal M/\partial\mathcal M,[\partial\mathcal M]).
 \label{wb:sphere-metric-C-boundary-quotient-homeomorphism}
\end{equation}
Only after establishing
\eqref{wb:sphere-metric-C-boundary-quotient-homeomorphism} do we apply the
quotient identity at the start of the proof.  It gives
\[
 \mathcal C^{\wedge4}
 \cong
 \frac{\mathcal M^4}
 {\displaystyle\bigcup_{a=1}^4
  \mathcal M^{a-1}\times\partial\mathcal M
  \times\mathcal M^{4-a}}
 =\mathcal M^4/\partial(\mathcal M^4),
\]
and carries the smash base point to the collapsed boundary.

\item Put \(N=\mathcal M^4\).  Rounding its product corners preserves the
homeomorphism type of the pair \((N,\partial N)\) and makes
\(\partial N\) an eleven-dimensional closed manifold.  A boundary collar
descends in \(N/\partial N\) to a cone neighborhood of the collapsed
point, so the local link is exactly \(\partial N\) and
\begin{equation}
 H_i(N/\partial N,(N/\partial N)\setminus\{*\};\mathbb Z)
 \cong\widetilde H_{i-1}(\partial N;\mathbb Z).
 \label{wb:sphere-metric-boundary-link-local-shift}
\end{equation}

\item The integral homology calculation in the proof of
\cref{wb:cusp-cone-S24-route} for the specific compact orientable
three-manifold \(\mathcal M\) gives
\[
 H_j(\mathcal M;\mathbb Z)\cong
 \begin{cases}
  \mathbb Z,&j=0,1,\\
  0,&\text{otherwise}.
 \end{cases}
\]
All groups are free, so the integral Kunneth theorem gives
\[
 H_j(N;\mathbb Z)\cong
 \begin{cases}
  \mathbb Z^{\binom4j},&0\leq j\leq4,\\
  0,&\text{otherwise}.
 \end{cases}
\]
Poincare--Lefschetz duality, with the product orientation retained, gives
\[
 H_j(N,\partial N;\mathbb Z)
 \cong H^{12-j}(N;\mathbb Z)
 \cong
 \begin{cases}
  \mathbb Z^{\binom4{j-8}},&8\leq j\leq12,\\
  0,&\text{otherwise}.
 \end{cases}
\]
The long exact sequence of \((N,\partial N)\) now has disjoint absolute
and relative degree ranges.  It gives
\[
 \widetilde H_j(\partial N;\mathbb Z)\cong
 \begin{cases}
  \mathbb Z^{\binom4j},&1\leq j\leq4,\\
  0,&j=5,6,\\
  \mathbb Z^{\binom4{j-7}},&7\leq j\leq11,\\
  0,&\text{otherwise}.
 \end{cases}
\]
Substitution in
\eqref{wb:sphere-metric-boundary-link-local-shift} proves
\eqref{wb:sphere-metric-C-smash-local-homology}.

\item To see explicitly why topology alone does not fix the tangent metric, compare the cone tensors $dr^2+r^2h$ and $dr^2+4r^2h$ on the same link of dimension $d$. The identity is a homeomorphism, while their unit-ball Hausdorff volumes differ by the factor $2^d$, so the marked metric cones are not isometric. An exact metric cone is invariant under every radial blow-up at its vertex,
and its Hausdorff dimension is one plus that of its fixed
piecewise-Riemannian link.  For the first target,
\[
 H_8\bigl(\operatorname{Cone}(T^3*T^3),T^3*T^3;\mathbb Z\bigr)
 \cong\mathbb Z;
\]
the analogous top local group for the \(\mathcal B\) target is
\(\mathbb Z\) in degree \(12\), while the full \(\mathcal C\) calculation
is \eqref{wb:sphere-metric-C-smash-local-homology}.  Euclidean
\(\mathbb R^4\) has local homology \(\mathbb Z\) only in degree \(4\), so
none of these conic vertex models is \(\mathbb R^4\).  The diameter
calculation proves the final assertion.
\end{enumerate}
\end{proof}

\paragraph{Consequences for the three sphere routes.}
The exact positive statement is that one fixed metric on each compact
codomain contains the native \(\mathbb R^4\) decompactification as a
vertical four-plane in a moving-point Euclidean bubble.  The whole bubbles
are \(\mathbb R^8\) for the \(S^{12}\) route and \(\mathbb R^{12}\) for the
\(S^{24}\) routes.  The factor \(\beta_T=\lambda_Tr_T\) proves that this
dimension discrepancy cannot be removed by a different scalar rescaling
of the same fixed cusp-adapted metric: the slice unrolls only when the
smash base point has already gone to infinite scaled distance.

No metric on a source sphere is supplied by these quotient constructions.
If an independent smooth Riemannian metric is placed on \(S^{12}\) or
\(S^{24}\), its ordinary tangent spaces are \(\mathbb R^{12}\) or
\(\mathbb R^{24}\).  Pulling a codomain metric back through
\(q_2^{\wedge2}\) or \(q_1^{\wedge4}\), on the regular strata where the
maps are differentiable, gives a degenerate quadratic form along the
positive-dimensional quotient kernels, not a Riemannian source metric.

The moving four-sphere carriers
\[
 \mathbb S_{12,T}^4=S_T^2\wedge S_T^2\subset S^6,\qquad
 \mathbb S_{24,T}^4=C_T^{\wedge4}\subset S^8
\]
of \cref{wb:cusp-cone-angle-forgotten-sphere-carriers} carry separately
assigned metrics.  Namely, \(S_T^2\) first receives the quotient length
metric obtained from
\((\overline T^2,h_T)\cong(Q_T,g_{Q_T})\) by collapsing
\(K^{(1)}\subset\overline T^2\), equivalently by collapsing
\(\iota_T(K^{(1)})\subset Q_T\).  Then
\(\mathbb S_{12,T}^4\) receives the quotient length metric of the
Riemannian \(\ell^2\)-product \((S_T^2)^2\) after its fat wedge is
collapsed.  Likewise, \(\mathbb S_{24,T}^4\) receives the quotient length
metric of the Riemannian \(\ell^2\)-product \(C_T^4\) after its fat wedge is
collapsed.  These metrics are assigned directly to the moving carriers;
they are unrescaled intrinsic quotient metrics, not restrictions of
\(D_{12}\) or \(D_{24}^{\mathcal B}\), whose domains are the conical
targets, and no fixed ambient metric on \(S^6\) or \(S^8\) is used to
define them.  Their Euclidean balls are
based at the midpoint classes
\[
 \begin{aligned}
  y_*&=[(1/2,1/2)]\in\overline T^2,\qquad z_T=c(y_*),\\
  \zeta_{12,T}&=[z_T,z_T],\\
  \zeta_{24,T}
  &=[(\ell_T/2,\ell_T/2,\ell_T/2,\ell_T/2)].
 \end{aligned}
\]
not at the images of \(p_T^{12}\) or \(p_T^{24,\mathcal B}\), which belong
to the subsequently collapsed loci.  In particular, the trichotomy in
\cref{wb:sphere-metric-scaling-trichotomy} is the trichotomy of the
uncollapsed intrinsic flat-torus slices; it is not a finite-\(\beta_T\)
description of the singular quotient-metric moving four-spheres.  Indeed,
the restriction of any one compatible metric on the compact sphere
\(S^6\) or \(S^8\) has diameter bounded by the diameter of that sphere.  By
contrast, in the separately assigned moving metrics the midpoint distances
to the collapsed loci are at least \((T-C)/2\) and exactly
\(\ell_T/2\), respectively, and hence tend to infinity.  Thus the assigned
unrescaled metrics cannot themselves be restrictions of one fixed ambient
metric.  This does not preclude a separately proved realization after a
scalar blow-up of another fixed quotient metric.

Finally, the rank obstruction for a full-base torus quotient must retain
its actual metric hypothesis: \emph{quotient systole tends to infinity}
means the quotient norm induced by the unrescaled Euclidean fibre metric.
Without that qualification the assertion is false, since an arbitrary
\(T\)-dependent scalar multiplication can make the systole of any fixed
nonzero quotient tend to infinity.

\subsection{Expanded proof of the retained-angle cusp class}
\label{wb:h24:ret:section}

\paragraph{The exact object.}
We use the fixed based smooth identification of the family member $X$ with
$S^6$, its cusp filling, and the complete original toric coordinates of
\cref{mapsI:src:A.1.2}. The native quotient and compact extension are
\cref{wb:cusp-cone-native-descent} and \cref{wb:cusp-cone-S12-route}.
The local parameter $t$ in the following formulas is the cusp parameter
$q$ of the family convention; $B_0=J$ in the radial metric notation.
All original period terms and native lattices are retained. The
calculation is independent of the angle-forgetting map to $S^3$.

\paragraph{Original coordinates and the specified sphere-valued map}

The ordered quotient lattice is
\[
 \overline\Lambda=\mathbb Z\langle\overline{\widehat\gamma},
              \overline{\widehat u}\rangle,
 \qquad
 K_\Lambda=\mathbb Z\langle\widehat w,\widehat\delta\rangle,
 \qquad
 \overline T^2=\mathbb R^2/\mathbb Z^2.
\]
In these ordered bases the original period data are
\begin{equation}
 B_0=\begin{pmatrix}0&1\\-1&0\end{pmatrix},\qquad
 Z(s)=sB_0+C(t),\qquad t=e^{2\pi i s},\qquad
 C(t)=\begin{pmatrix}6\mu(t)&h_{\rm per}(t)\\
 b(t)-h_{\rm per}(t)&\mu(t)\end{pmatrix}.
 \label{wb:h24:ret:eq:original-periods}
\end{equation}
Here \(b=\beta+\tau\), \(s=\tau-h_{\rm per}(t)\), and every entry of
\(C\) is holomorphic at \(t=0\), exactly as in the original source.
The subscript in \(h_{\rm per}\) distinguishes that period function
from the sphere-valued map below.  Write
\begin{equation}
 R(t)=-2\pi\operatorname{Im}C(t),\qquad
 \ell=\log|t|,\qquad A(t)=\ell B_0+R(t).
 \label{wb:h24:ret:eq:A}
\end{equation}
Choose a closed cusp disc \(|t|\leq\epsilon\) inside the original
coordinate disc and small enough for the quotient construction.  On a
slightly larger closed disc put \(M=\sup\|R(t)\|<\infty\).  We will
choose a regular value with \(0<r_0<\epsilon\) and
\(|\log r_0|>M\).  This chooses a point of the already specified map;
it does not change that map or the cutoff radius \(\epsilon\).

The original covering model on its dense torus is
\begin{equation}
 x_j=e^{2\pi i\zeta_j},\qquad
 \Psi_n(x,t)=
 \bigl(e^{2\pi i C(t)n}t^{B_0n}x,t\bigr),\qquad n\in\mathbb Z^2.
 \label{wb:h24:ret:eq:deck}
\end{equation}
For a branch of \(s\), writing
\(\zeta=Z(s)u+k\) with \(u,k\in\mathbb R^2\) gives
\begin{equation}
 \log|x|=A(t)u.
 \label{wb:h24:ret:eq:logx}
\end{equation}
Indeed, \(\operatorname{Im}s=-\ell/(2\pi)\), so
\(-2\pi\operatorname{Im}(Z(s)u+k)=\ell B_0u+R(t)u\).
Since \(B_0\) is orthogonal,
\begin{equation}
 \bigl|\bigl(\ell B_0+vR(t)\bigr)\xi\bigr|
 \geq (|\ell|-vM)|\xi|>0
 \quad(0\leq v\leq1,\ \xi\ne0)
 \label{wb:h24:ret:eq:invertibility}
\end{equation}
whenever \(|\ell|>M\).  Thus \(A(t)\) is invertible there.  Equation
\eqref{wb:h24:ret:eq:deck} changes \(A(t)^{-1}\log|x|\) by \(n\), so the native
quotient is exactly
\begin{equation}
 \overline\pi_f([x,t])=[A(t)^{-1}\log|x|]\in\mathbb R^2/\mathbb Z^2.
 \label{wb:h24:ret:eq:native}
\end{equation}
This also follows from the flat marking in \cref{wb:cusp-cone-native-descent}.

Let
\[
 T^3=S^1_\phi\times\overline T^2,
 \qquad S^1_\phi=\mathbb R/(2\pi\mathbb Z),
\]
with coordinate basepoints \(\phi=0\), \(u_1=0\), \(u_2=0\).
Its coordinate two-skeleton is the following specified closed subset:
\begin{equation}
 K^{(2)}=
 \bigl(\{0\}\times\overline T^2\bigr)
 \cup\bigl(S^1_\phi\times\{u_1=0\}\bigr)
 \cup\bigl(S^1_\phi\times\{u_2=0\}\bigr).
 \label{wb:h24:ret:eq:two-skeleton}
\end{equation}
The product CW structure has one open three-cell, with coordinates
\((\phi,u_1,u_2)\in(0,2\pi)\times(0,1)^2\).  Consequently
\[
 c_3:T^3\longrightarrow T^3/K^{(2)}\cong S^3
\]
is the degree-one quotient, where the orientation of \(S^3\) is fixed
by the displayed coordinate order.  One explicit identification of its
open cell with \(\mathbb R^3\) is
\[
 (\phi,u_1,u_2)\longmapsto
 \bigl(-\cot(\phi/2),-\cot(\pi u_1),-\cot(\pi u_2)\bigr).
\]
Each derivative in its own variable is strictly positive, and approach
to any boundary face leaves every compact subset of \(\mathbb R^3\).
Extending by its point at infinity gives the stated based sphere
identification and verifies its continuity on the collapsed skeleton.

The exact conical target and map are
\begin{align}
 \mathcal A&=
 \frac{[0,\epsilon]\times T^3}
 {\{0\}\times T^3\sim v_0,\quad
  \{\epsilon\}\times T^3\sim v_1},
 \label{wb:h24:ret:eq:A-target}\\
 q_2(x)&=[r,e^{i\phi},\overline\pi_f(x)]
       &&(t(f(x))=re^{i\phi},\ 0<r<\epsilon),\notag\\
 q_2(W)&=v_0,\qquad
 q_2\bigl(X\setminus\operatorname{int}N_0\bigr)=v_1.
 \label{wb:h24:ret:eq:q2}
\end{align}
Fix the same source basepoint \(x_{\rm base}\in W\).  Define
\begin{equation}
 p_3:\mathcal A\longrightarrow\Sigma_r(S^3,*),\qquad
 p_3([r,z])=[r,c_3(z)],\qquad
 H=p_3q_2:(X,x_{\rm base})\longrightarrow(S^4,*).
 \label{wb:h24:ret:eq:H}
\end{equation}
The reduced suspension here collapses both endpoint copies of \(S^3\)
and the entire interval over \(*\) to one point.  Thus \(p_3\) sends
both \(v_0\) and \(v_1\) to \(*\) and is well-defined and based.
It is continuous by the defining quotient topology.  An explicit
identification of its open four-cell with \(\mathbb R^4\) is
\begin{equation}
 (r,\phi,u_1,u_2)\longmapsto
 \bigl(-\cot(\pi r/\epsilon),-\cot(\phi/2),
       -\cot(\pi u_1),-\cot(\pi u_2)\bigr).
 \label{wb:h24:ret:eq:four-cell}
\end{equation}
This fixes the target orientation.  Replacing \(r\) locally by
\(\ell=\log r\) preserves this orientation, since \(dr/d\ell=r>0\).
The ordered coordinates in \eqref{wb:h24:ret:eq:four-cell} specify the sphere
identification used throughout this calculation.

\paragraph{The statement proved.} We give the complete expanded proof of \cref{maps:src:B.3.1}, including the full differential, the relative smoothing and the spin bordism calculation.


\paragraph{The entire regular inverse image lies in one original chart}

Choose the actual lift
\begin{equation}
 u_0=(-1/4,1/4)^{\mathsf T},\qquad
 p_0=B_0u_0=(1/4,1/4)^{\mathsf T},\qquad \phi_0=\pi.
 \label{wb:h24:ret:eq:chosen-value}
\end{equation}
Modulo \(\mathbb Z^2\), \(u_0\) is \((3/4,1/4)\) in the open cell.  Put
\(t_0=r_0e^{i\phi_0}\) and
\(y_0=[r_0,c_3(e^{i\phi_0},[u_0])]\).  The numbers \(r_0\) and
\(t_0\) will remain nonzero throughout the framing calculation.

The original triangle \(\{0,e_1,e_2\}\) has toric chart
\begin{equation}
 \mathcal U_\sigma\cong\mathbb C^3,\qquad
 (z_0,z_1,z_2)=(t/(x_1x_2),x_1,x_2),\qquad
 t=z_0z_1z_2.
 \label{wb:h24:ret:eq:z-chart}
\end{equation}
These characters are dual to \((0,0,1),(1,0,1),(0,1,1)\), respectively,
as direct pairing verifies.  This is the chart of
\cref{mapsI:src:A.1.2}, with its original coordinates.

For \(j=1,2\), introduce the nowhere vanishing holomorphic functions
\begin{equation}
 h_j(t)=\exp\bigl(-2\pi i(C(t)u_0)_j\bigr).
 \label{wb:h24:ret:eq:units}
\end{equation}
The following is an exact coordinate map, with the exact inverse
displayed alongside it:
\begin{align}
 F(z)&=(w_0,w_1,w_2)
 =\bigl(z_0/(h_1(t)h_2(t)),\ h_1(t)z_1,\ h_2(t)z_2\bigr),
 \label{wb:h24:ret:eq:F}\\
 F^{-1}(w)&=
 \bigl(w_0h_1(t)h_2(t),\ w_1/h_1(t),\ w_2/h_2(t)\bigr),
 \quad t=w_0w_1w_2.
 \label{wb:h24:ret:eq:Finv}
\end{align}
Both maps are holomorphic wherever \(|t|\) is in the period disc,
both preserve \(t\), and substitution proves that both composites are
the identity.  In particular the maps extend across the origin and
across every coordinate divisor.  No period matrix is replaced by a
different matrix in this coordinate change.

\begin{lemma}\label{wb:h24:ret:chart-lemma}
\StatementRecord{wb:h24:ret:chart-lemma}{proved}{Historical events: GPT-6 Astra (OpenAI), ChatGPT Pro mode, effort Pro, S6\_09\_PAPERS, 17 November 2026: targeted proof check. Also checked in this writing. Pre-continuation writing events: GPT-6 Astra (OpenAI), ChatGPT Pro mode, effort Pro, S6\_10\_PAPERS, 18 November 2026: targeted proof check. Continuation metadata events: GPT-6 Astra (OpenAI), ChatGPT Pro mode, effort Pro, S6\_10\_PAPERS continuation, 19 November 2026: used and rechecked in this run; status-vocabulary reconciliation; metadata reconciliation only; no new mathematical check. \emph{Pre-continuation recorded proof scope.} /root, effort Pro, 2026-11-17; also checked in this writing by GPT-6 Astra (OpenAI), ChatGPT Pro mode, effort Pro, 18 November 2026}{the hypotheses, notation and definitions stated here; \cref{wb:h24:ret:eq:H}; \cref{wb:h24:ret:eq:q2}; \cref{wb:h24:ret:eq:deck}; \cref{wb:h24:ret:eq:z-chart}}

For sufficiently small \(r_0>0\), the whole inverse image
\(M=H^{-1}(y_0)\) is represented injectively inside one contractible
coordinate neighborhood in \(X\).  In its \(w\)-coordinates it is
exactly the embedded two-torus
\begin{equation}
 w(a,b)=
 \bigl(\rho_0e^{i(\phi_0-a-b)},\rho_1e^{ia},\rho_2e^{ib}\bigr),
 \qquad
 \rho_0=r_0^{1/2},\quad \rho_1=\rho_2=r_0^{1/4},
 \quad (a,b)\in\bigl(\mathbb R/(2\pi\mathbb Z)\bigr)^2.
 \label{wb:h24:ret:eq:M}
\end{equation}
\end{lemma}


\begin{proof}
\begin{enumerate}[label=\arabic*.,leftmargin=*]
\item The point \(y_0\) lies off the collapsed subset of the target, so
\eqref{wb:h24:ret:eq:H} and \eqref{wb:h24:ret:eq:q2} imply
\[
 M=\{x:t(f(x))=t_0,\ \overline\pi_f(x)=[u_0]\}.
\]
Any lift of such a point to \((\mathbb C^*)^2\times\{t_0\}\) has
\(A(t_0)^{-1}\log|x|=u_0+n\) for one \(n\in\mathbb Z^2\).
Applying the actual transformation \(\Psi_{-n}\) from
\eqref{wb:h24:ret:eq:deck} produces a lift with
\(\log|x|=A(t_0)u_0\).  If two such lifts represent the same point,
\(\Psi_n\) relates them and changes their lifted quotient coordinate
by \(n\); hence \(n=0\), and the lifts are identical.  This proves
both completeness and injectivity of this chosen representative.

\item Since
\(\log|h_j(t)|=2\pi\operatorname{Im}(C(t)u_0)_j
=-(R(t)u_0)_j\), its \(w\)-coordinates satisfy
\[
 (\log|w_1|,\log|w_2|)^{\mathsf T}
 =A(t_0)u_0-R(t_0)u_0=\ell_0p_0,
 \qquad \ell_0=\log r_0.
\]
Thus \(|w_1|=|w_2|=r_0^{1/4}\), and the unchanged product
\(w_0w_1w_2=t_0\) forces \(|w_0|=r_0^{1/2}\) and the displayed
phase.  Conversely every pair \((a,b)\) in \eqref{wb:h24:ret:eq:M} gives one
such lift.  The original kernel coordinates \(k_j\) had period one
in \(x_j=e^{2\pi i\zeta_j}\); the phase coordinates \(a,b\) have
period \(2\pi\), with no cover or change of lattice.

\item The covering \(Y_\epsilon\to N_0\) is injective on a neighborhood of
the origin in \eqref{wb:h24:ret:eq:z-chart}, because the original action is free
and properly discontinuous.  Choose a ball \(V\) centered at that
origin with closure inside such a neighborhood and the period disc.
Choose a ball \(V_w\) about the origin with closure contained in
\(F(V)\).  Every coordinate in \eqref{wb:h24:ret:eq:M} tends uniformly to zero
as \(r_0\downarrow0\).  Therefore for sufficiently small \(r_0\)
the entire torus lies in \(V_w\).  The inverse image of \(V_w\)
under \(F\), followed by the injective covering projection, is the
required contractible chart in \(X\).
\end{enumerate}
\end{proof}

\paragraph{The full differential and its ordered normal framing}

We calculate in a neighborhood of the compact torus \eqref{wb:h24:ret:eq:M}.
The quotient coordinates admit a unique local lift near \(u_0\).
Set
\[
 v=(\log|w_1|,\log|w_2|)^{\mathsf T},\quad
 \ell=\log|w_0w_1w_2|,\quad
 \phi=\arg(w_0w_1w_2)
\]
using the branch of \(\phi\) near \(\phi_0=\pi\).
Equations \eqref{wb:h24:ret:eq:logx} and \eqref{wb:h24:ret:eq:F} give the exact identity
\begin{equation}
 v=\ell B_0u+R(t)(u-u_0),\qquad
 u=u_0+A(t)^{-1}(v-\ell p_0).
 \label{wb:h24:ret:eq:exact-u}
\end{equation}
Differentiate this identity before restricting to the inverse image:
\begin{equation}
 du=-A(t)^{-1}(dA(t))A(t)^{-1}(v-\ell p_0)
        +A(t)^{-1}(dv-p_0\,d\ell).
 \label{wb:h24:ret:eq:full-derivative}
\end{equation}
On \(M\) the vector \(v-\ell p_0\) is exactly zero.  Therefore
\begin{equation}
 du|_M=A_0^{-1}(dv-p_0\,d\ell),\qquad A_0=A(t_0).
 \label{wb:h24:ret:eq:derivative-on-M}
\end{equation}
This explicitly accounts for all derivatives of the original bounded
period terms; none are omitted or estimated away.

Identify \(T_w\mathbb C^3\) with \(\mathbb C^3\) in the given coordinate chart,
and define real tangent vectors
\begin{equation}
 L_j=(0,\ldots,w_j,\ldots,0),\qquad
 Q_j=(0,\ldots,iw_j,\ldots,0)\qquad(0\leq j\leq2).
 \label{wb:h24:ret:eq:LQ}
\end{equation}
Thus \(d\log|w_i|(L_j)=\delta_{ij}\),
\(d\arg w_i(Q_j)=\delta_{ij}\), and the crossed evaluations vanish.
The ordered tangent vectors of the parametrization \eqref{wb:h24:ret:eq:M} are
\begin{equation}
 T_a=-Q_0+Q_1,\qquad T_b=-Q_0+Q_2.
 \label{wb:h24:ret:eq:tangents}
\end{equation}
They are independent because their \(w_1\) and \(w_2\) components,
respectively, are nonzero.  For the target order
\((\ell,\phi,u_1,u_2)\), define the following four lifts of its
coordinate vectors:
\begin{align}
 N_\ell&=\tfrac12L_0+\tfrac14L_1+\tfrac14L_2,
       &N_\phi&=Q_0,\label{wb:h24:ret:eq:normal-first}\\
 N_i&=-\bigl((A_0)_{1i}+(A_0)_{2i}\bigr)L_0
           +(A_0)_{1i}L_1+(A_0)_{2i}L_2
       &&(i=1,2).\label{wb:h24:ret:eq:normal-u}
\end{align}
Let \(P=(\ell,\phi,u_1,u_2)\).  Evaluation using
\eqref{wb:h24:ret:eq:derivative-on-M} gives the entire differential table
\begin{equation}
\begin{array}{c|rrrrrr}
 &N_\ell&N_\phi&N_1&N_2&T_a&T_b\\\hline
 d\ell&1&0&0&0&0&0\\
 d\phi&0&1&0&0&0&0\\
 du_1&0&0&1&0&0&0\\
 du_2&0&0&0&1&0&0
\end{array}.
 \label{wb:h24:ret:eq:differential-table}
\end{equation}
In particular \(P\), and therefore \(H\) in its open target cell,
is a submersion along \(M\).  The classes of
\((N_\ell,N_\phi,N_1,N_2)\) in \(TX|_M/TM\) give the specified
normal four-framing.  Orthogonal representatives are not required for
this quotient-bundle definition; projecting the displayed lifts along
\(TM\) to any chosen orthogonal complement preserves the quotient
frame.

If the normal frame is specified using the literal four real coordinates
in \eqref{wb:h24:ret:eq:four-cell}, their differential at the chosen value is
\[
 dH_{\rm cell}=K\,dP,\qquad
 K=\operatorname{diag}\!\left(
 \frac{\pi r_0}{\epsilon}\csc^2\frac{\pi r_0}{\epsilon},
 \frac12,2\pi,2\pi\right).
\]
The first factor includes \(dr=r_0\,d\ell\), the second uses
\(\phi_0=\pi\), and the final two use the actual representatives
\((3/4,1/4)\) of \(u_0\). Thus, writing
\(\mathcal N=(N_\ell,N_\phi,N_1,N_2)\), the literal dual frame is
\(\mathcal N K^{-1}\). All four entries of \(K\) are positive
constants on \(M\). The explicitly positive diagonal matrices
\((1-v)K^{-1}+vI_4\), \(0\leq v\leq1\), give a homotopy
from that literal frame to \(\mathcal N\). This retains the original
target coordinates and establishes that either coordinate description
gives the same framed bordism class.

For completeness, the original period matrix admits the following
explicit homotopy of this normal framing.  Put
\begin{equation}
 A_{0,v}=\ell_0B_0+vR(t_0),\qquad 0\leq v\leq1,
 \label{wb:h24:ret:eq:A-homotopy}
\end{equation}
and replace \(A_0\) by \(A_{0,v}\) only in
\eqref{wb:h24:ret:eq:normal-u}.  Equation \eqref{wb:h24:ret:eq:invertibility} proves
invertibility for every \(v\).  Their images under the actual
\((du_1,du_2)\) are the columns of \(A_0^{-1}A_{0,v}\), so these
are normal frames throughout the homotopy.  This is a homotopy of
the computed frame on the same embedded torus, not a replacement of
the original degeneration by a different one.

\begin{lemma}[The exact orientation sign]\label{wb:h24:ret:orientation-lemma}
\StatementRecord{wb:h24:ret:orientation-lemma}{proved}{Historical events: GPT-6 Astra (OpenAI), ChatGPT Pro mode, effort Pro, S6\_09\_PAPERS, 17 November 2026: targeted proof check. Also checked in this writing. Pre-continuation writing events: GPT-6 Astra (OpenAI), ChatGPT Pro mode, effort Pro, S6\_10\_PAPERS, 18 November 2026: targeted proof check. Continuation metadata events: GPT-6 Astra (OpenAI), ChatGPT Pro mode, effort Pro, S6\_10\_PAPERS continuation, 19 November 2026: used and rechecked in this run; status-vocabulary reconciliation; metadata reconciliation only; no new mathematical check. \emph{Pre-continuation recorded proof scope.} /root, effort Pro, 2026-11-17; also checked in this writing by GPT-6 Astra (OpenAI), ChatGPT Pro mode, effort Pro, 18 November 2026}{the hypotheses, notation and definitions stated here; \cref{wb:h24:ret:eq:normal-u}; \cref{wb:h24:ret:eq:A-homotopy}}

Use the complex real-coordinate orientation, whose ordered coordinates are
\[
 (\operatorname{Re}w_0,\operatorname{Im}w_0,
  \operatorname{Re}w_1,\operatorname{Im}w_1,
  \operatorname{Re}w_2,\operatorname{Im}w_2).
\]
Then one has
\begin{equation}
 \det(N_\ell,N_\phi,N_1,N_2,T_a,T_b)
       =-\det(A_0)|w_0w_1w_2|^2<0.
 \label{wb:h24:ret:eq:sign}
\end{equation}
Hence the tangent orientation induced by the ordered normal frame is
the orientation of \((-T_a,T_b)\).
\end{lemma}


\begin{proof}
\begin{enumerate}[label=\arabic*.,leftmargin=*]
\item The ordered frame \((L_0,Q_0,L_1,Q_1,L_2,Q_2)\) has determinant
\(|w_0|^2|w_1|^2|w_2|^2>0\) in real coordinates.  Regroup its rows
as \((L_0,L_1,L_2,Q_0,Q_1,Q_2)\); this permutation has three
inversions and sign \(-1\).  Regroup the columns of the frame in
\eqref{wb:h24:ret:eq:sign} as
\((N_\ell,N_1,N_2,N_\phi,T_a,T_b)\); this has two inversions and
sign \(+1\).  The angular block is
\(\left(\begin{smallmatrix}1&-1&-1\\0&1&0\\0&0&1\end{smallmatrix}\right)\),
of determinant one.  The radial block has first column
\((1/2,1/4,1/4)^{\mathsf T}\) and the other columns given by
\eqref{wb:h24:ret:eq:normal-u}.  Adding its last two rows to its first row gives
first row \((1,0,0)\), with remaining lower-right block \(A_0\).
Its determinant is \(\det A_0\).  This proves the equality.
Equation \eqref{wb:h24:ret:eq:A-homotopy} joins \(A_0\) to \(\ell_0B_0\)
through invertible matrices, so
\(\operatorname{sign}\det A_0=\operatorname{sign}\det(\ell_0B_0)=+1\).
The inequality and the tangent orientation follow.
\end{enumerate}
\end{proof}

\paragraph{An explicit contraction of the full frame}

Define a complex-linear matrix, also regarded as a real-linear matrix,
\begin{equation}
 D(a,b)=\operatorname{diag}(e^{-i(a+b)},e^{ia},e^{ib})\in\operatorname{SU}(3).
 \label{wb:h24:ret:eq:D}
\end{equation}
By \eqref{wb:h24:ret:eq:M}, every \(L_j\) and \(Q_j\) at \((a,b)\) is
\(D(a,b)\) times its value at \((0,0)\).  The coefficients in
\eqref{wb:h24:ret:eq:normal-first}, \eqref{wb:h24:ret:eq:normal-u}, and
\eqref{wb:h24:ret:eq:tangents} are constant on \(M\), since \(t=t_0\).
It follows that the positively oriented full frame is exactly
\begin{equation}
 (N_\ell,N_\phi,N_1,N_2,-T_a,T_b)_{(a,b)}
       =D(a,b)G_0,
 \label{wb:h24:ret:eq:full-frame}
\end{equation}
where \(G_0\in\operatorname{GL}^+(6,\mathbb R)\) is its single value at \((0,0)\).
Thus neither an unspecified phase nor a position-dependent
real-linear matrix remains.

Here is a periodic contraction of \(D\), with all parameters
specified.  For \(v\in[0,1]\) put
\[
 s_v=\sin(\pi v/2),\qquad c_v=\cos(\pi v/2),\qquad
 Q_v(\theta)=
 \begin{pmatrix}s_ve^{-i\theta}&-c_v\\c_v&s_ve^{i\theta}\end{pmatrix},
 \qquad K_v(\theta)=Q_v(\theta)Q_v(0)^*.
\]
The two columns of \(Q_v\) have norm one and zero Hermitian inner
product, and \(\det Q_v=s_v^2+c_v^2=1\).  Thus
\(Q_v,K_v\in\operatorname{SU}(2)\).  Direct substitution gives
\[
 K_0(\theta)=I,\quad
 K_1(\theta)=\operatorname{diag}(e^{-i\theta},e^{i\theta}),\quad
 K_v(0)=I.
\]
Let \(\iota_{01},\iota_{02}:\operatorname{SU}(2)\to\operatorname{SU}(3)\) act in the
ordered coordinate planes \((0,1)\) and \((0,2)\), respectively,
and as the identity in the unused coordinate.  Then
\begin{equation}
 D_v(a,b)=\iota_{01}(K_v(a))\,\iota_{02}(K_v(b))
 \label{wb:h24:ret:eq:D-homotopy}
\end{equation}
is a well-defined continuous map on
\([0,1]\times(\mathbb R/(2\pi\mathbb Z))^2\), with
\(D_0=I\) and \(D_1=D\).  This proves the needed null-homotopy
directly, including both phase seams.  In particular
\(D_v(a,b)G_0\) contracts the comparison matrix in
\eqref{wb:h24:ret:eq:full-frame} to one constant positive matrix.

\begin{proposition}\label{wb:h24:ret:framing-proposition}
\StatementRecord{wb:h24:ret:framing-proposition}{proved}{Historical events: GPT-6 Astra (OpenAI), ChatGPT Pro mode, effort Pro, S6\_09\_PAPERS, 17 November 2026: targeted proof check. Also checked in this writing. Pre-continuation writing events: GPT-6 Astra (OpenAI), ChatGPT Pro mode, effort Pro, S6\_10\_PAPERS, 18 November 2026: targeted proof check. Continuation metadata events: GPT-6 Astra (OpenAI), ChatGPT Pro mode, effort Pro, S6\_10\_PAPERS continuation, 19 November 2026: used and rechecked in this run; status-vocabulary reconciliation; metadata reconciliation only; no new mathematical check. \emph{Pre-continuation recorded proof scope.} /root, effort Pro, 2026-11-17; also checked in this writing by GPT-6 Astra (OpenAI), ChatGPT Pro mode, effort Pro, 18 November 2026}{the hypotheses, notation and definitions stated here; \cref{wb:h24:ret:eq:sign}; \cref{wb:h24:ret:eq:full-frame}; \cref{wb:h24:ret:eq:D-homotopy}}

The stable tangential framing induced by the actual normal four-frame
of \(H^{-1}(y_0)\) agrees, up to homotopy and the orientation sign
\eqref{wb:h24:ret:eq:sign}, with the Lie framing
\((\partial_a,\partial_b)\) on the native two-torus.
\end{proposition}


\begin{proof}
\begin{enumerate}[label=\arabic*.,leftmargin=*]
\item In the coordinate chart, the stable tangent trivialization is the
isomorphism
\[
 \mathbb R^4\oplus TM\longrightarrow T\mathbb C^3|_M,
 \qquad (v,V)\longmapsto
 v_1N_\ell+v_2N_\phi+v_3N_1+v_4N_2+V.
\]
Using the intrinsic frame \((-\partial_a,\partial_b)\) on \(TM\),
its matrix is precisely \eqref{wb:h24:ret:eq:full-frame}.  The homotopy
\eqref{wb:h24:ret:eq:D-homotopy} therefore makes this comparison constant.
A constant positive matrix changes no homotopy class of stable
framing: Gram--Schmidt deforms it to an element of \(SO(6)\), and
an orthogonal matrix of determinant one can be reduced to the identity
by finitely many oriented plane rotations, each joined to the
identity by varying its angle.  This proves the asserted comparison
within the chart.

\item It remains to compare the chart with the ambient sphere framing.
Pull back the usual stable tangent trivialization
\(TS^6\oplus\mathbb R\cong\mathbb R^7\) under the fixed based diffeomorphism
\(X\cong S^6\).  On the contractible \(w\)-ball containing \(M\),
it differs from the coordinate trivialization plus one real line by
a continuous map into \(\operatorname{GL}(7,\mathbb R)\).  Contracting the ball to its
center contracts this comparison map to a constant.  Thus its
restriction to \(M\) contributes no varying stable framing class.
If the chosen sphere orientation disagrees with the complex
orientation, the remaining constant has negative determinant and
reverses the resulting framed bordism class.  Both that sign and the
explicit sign \eqref{wb:h24:ret:eq:sign} disappear in the order-two class
computed below.  This argument compares actual ambient
trivializations on a contractible chart; it does not suppose that
\(\mathbb C^3\) globally trivializes \(TX\).
\end{enumerate}
\end{proof}

\paragraph{Nonzero framed bordism and the exact unstable class}

We give the nonvanishing test for the resulting torus explicitly.
A stable framing gives a spin structure by lifting its oriented frame
to the trivial \(\operatorname{Spin}\)-bundle.  On the Lie-framed
circle this is the disconnected double cover, whereas the spin
structure induced from the boundary of a framed disc is the
connected double cover: as the boundary tangent turns once, its
\(2\pi\)-rotation lifts from \(1\) to \(-1\) in
\(\operatorname{Spin}(2)\).  We call the former nonbounding and the
latter bounding.

For a closed oriented spin surface \(S\), the spin quadratic
refinement
\[
 q_S:H_1(S;\mathbb F_2)\longrightarrow\mathbb F_2
\]
takes a class represented by an embedded circle to zero or one
according to whether the induced spin structure, using its normal
line in the surface, is bounding or nonbounding.  It satisfies
\begin{equation}
 q_S(x+y)=q_S(x)+q_S(y)+(x\mathbin{\cdot}y).
 \label{wb:h24:ret:eq:quadratic}
\end{equation}
This is the usual spin quadratic form: resolving transverse
intersections adds one tangent turn per intersection, and the
\(2\pi\)-spin-lift calculation above gives the last term.  The
well-defined correspondence between spin structures and such
quadratic refinements is given in
\cite[Theorem 3.2 and Remark 3.3]{DebrayGunningham2018Root}; its original
surface proof is \cite{Johnson1980Root}.

For a symplectic basis \((e_i,f_i)_{i=1}^g\), put
\begin{equation}
 \operatorname{Arf}(S)=\sum_{i=1}^gq_S(e_i)q_S(f_i)\in\mathbb F_2.
 \label{wb:h24:ret:eq:Arf}
\end{equation}
The basis independence can be seen without selecting a preferred
basis: factor the finite sum over the symplectic planes to obtain
\[
 \sum_{x\in H_1(S;\mathbb F_2)}(-1)^{q_S(x)}
   =\prod_{i=1}^g
     \sum_{a,b\in\mathbb F_2}(-1)^{a q_S(e_i)+b q_S(f_i)+ab}
   =2^g(-1)^{\sum_i q_S(e_i)q_S(f_i)}.
\]
The left side has no choice of basis, proving the assertion.

Here is also the bordism-invariance argument needed for the present
obstruction.  A spin bordism has a handle decomposition.  An index
zero or three handle creates or removes a spin sphere, whose first
homology and Arf invariant are zero.  An index one handle either
connects two components, in which case the first-homology quadratic
forms add orthogonally, or increases the genus of one component.
In the latter case the new symplectic pair contains the meridian
of the handle.  That meridian bounds the cocore disc with its
induced spin structure, so its quadratic value is zero; the pair
contributes zero to \eqref{wb:h24:ret:eq:Arf}, regardless of the value on its
dual.  Existing cycles can be represented outside the attaching
discs, so their quadratic values are preserved.  An index two
handle is the reverse of an index one handle.  Therefore each
handle preserves the total Arf invariant of the boundary, and a
spin boundary has Arf invariant zero.  This also proves that a
stably framed boundary has Arf invariant zero, because forgetting
the framing gives a spin bordism.  The full spin bordism statement
is recorded in \cite[Theorem 3.6]{DebrayGunningham2018Root}.

The intrinsic Lie-torus calculation is the known computation in \cite[Example~3.8, p.~6]{DebrayGunningham2018Root}; we retain its direct derivation. For the intrinsic Lie framing on
\(M=(\mathbb R/(2\pi\mathbb Z))^2\), take
\[
 e=[a\mapsto(a,0)],\qquad f=[b\mapsto(0,b)].
\]
Both the tangent and normal vectors in the surface are constant in
the Lie frame along either circle.  Thus both circles have the
nonbounding spin structure, and
\begin{equation}
 e\mathbin{\cdot}f=1,\qquad q_M(e)=q_M(f)=1,
 \qquad q_M(xe+yf)=x+xy+y,
 \qquad \operatorname{Arf}(M)=1.
 \label{wb:h24:ret:eq:actual-Arf}
\end{equation}
Reversing the first Lie tangent vector changes the orientation but
does not change these mod-two values.  In particular the
specifically framed inverse image computed above is not a framed
boundary.

To identify the element, let \(\eta\in\pi_1^{\rm S}\) denote the
stable class of the complex Hopf map.  Its Pontryagin--Thom inverse
image is the circle with its Lie framing.  This can be verified
directly on the Hopf map \(S^3\to\mathbb{CP}^1\): near the fibre
\((z_1,z_2)=(e^{ia},0)\), the target coordinate is \(z_2/z_1\), so
the two horizontal normal vectors are \((0,e^{ia})\) and
\((0,ie^{ia})\).  Add the tangent vector \((ie^{ia},0)\) and the
outward radial vector \((e^{ia},0)\).  In \(\mathbb C^2\) the resulting
full frame makes one complete rotation in each of its two real
coordinate planes.  This loop is contractible explicitly: for the
fixed real-linear involution \(J(z_1,z_2)=(z_1,\overline z_2)\),
\[
 J\operatorname{diag}(e^{ia},e^{ia})J=\operatorname{diag}(e^{ia},e^{-ia})=K_1(-a),
\]
and \(JK_v(-a)J\) contracts the original loop to the identity using
the periodic homotopy already constructed.  Thus the induced stable
tangent frame on the circle is the Lie frame.  This is the basic
framed-circle identification in \cite[Theorems~20--21, pp.~96--101]{Pontryagin1955Root}.

The product of two normally framed inverse images represents the
product in stable homotopy: its tubular collapse is the smash of
the two tubular collapse maps, as follows immediately from the
product tubular neighborhood \(\nu_1\times\nu_2\) and its ordered
normal coordinates.  Consequently the product of the two Lie
circles represents \(\eta^2\).  Equation \eqref{wb:h24:ret:eq:actual-Arf}
proves \(\eta^2\ne0\) by an explicit bordism invariant.  The
stable group \(\pi_2^{\rm S}=\mathbb Z/2\) is Pontryagin's computation
\cite[pp.~957--959]{Pontryagin1950Root} and \cite[Theorem~24, pp.~109--112]{Pontryagin1955Root}; the stable table in
\cite[opening table, p.~339, and Corollary~4.24, pp.~360--361]{Hatcher2002} gives the same group.  Thus the sign from
\eqref{wb:h24:ret:eq:sign} satisfies \(-\eta^2=\eta^2\).

\begin{proof}[Expanded proof of \cref{maps:src:B.3.1}]
\begin{enumerate}[label=\arabic*.,leftmargin=*]
\item The globally continuous quotient \(H\) is smooth on a neighborhood
of its compact inverse image \(M\), by
\eqref{wb:h24:ret:eq:exact-u} and \eqref{wb:h24:ret:eq:four-cell}. First make it constant
near the source basepoint, before choosing any approximation bound.
Choose a target chart \(\kappa:O\to\mathbb R^4\) about \(*\) with
\(\kappa(*)=0\), convex ball image, and \(y_0\notin O\).
Choose a source neighborhood \(U\) of \(x_{\rm base}\), disjoint
from \(M\), with \(H(U)\subset O\). Let \(\beta\) be a smooth
function supported in \(U\), equal to one near \(x_{\rm base}\),
with \(0\leq\beta\leq1\). On \(U\) use the based homotopy
\[
 H_s(x)=\kappa^{-1}\bigl((1-s\beta(x))\kappa(H(x))\bigr),
 \qquad 0\leq s\leq1,
\]
and leave \(H\) unchanged off \(U\). Its endpoint \(H_b\) is
smooth and constant near the basepoint; it agrees with \(H\)
near \(M\) and has no new inverse image of \(y_0\).

\item Choose one closed tubular neighborhood \(U_1\) of \(M\), with
\(M\subset\operatorname{int}U_1\), contained in the unchanged
smooth region and disjoint from the basepoint neighborhood.
Compactness gives the positive bound
\[
 d_0=\min_{x\in X\setminus\operatorname{int}U_1}
             d\bigl(H_b(x),y_0\bigr)>0.
\]
Relative smooth approximation, fixing the same closed set \(U_1\)
and a neighborhood of the basepoint, produces a smooth based map
\(H_{\rm sm}\) whose error from \(H_b\) is less than \(d_0/2\)
on its complement. Choose the approximation closer also than a
fixed target injectivity radius; the short target geodesics give a
based homotopy. Thus the exact inverse image and its differential
are unchanged on \(U_1\), and the complement acquires no new
inverse image. The relative Pontryagin construction applies as in
\cite[\S7, Theorem B and Lemma 4, pp.~43, 48--50]{CuspNullMilnor1965}.

\item Pontryagin--Thom therefore identifies the stabilization of
\([H]\) with the bordism class of \(M\) carrying the four normal
vectors \eqref{wb:h24:ret:eq:normal-first}--\eqref{wb:h24:ret:eq:normal-u}.  The comparison
\eqref{wb:h24:ret:eq:full-frame} and its explicit contraction identify that
class with the Lie torus, with the harmless order-two orientation
sign already computed.  Hence
\[
 E^\infty[H]=\eta^2\ne0.
\]
Freudenthal suspension is an isomorphism
\[
 E:\pi_{n+2}(S^n)\longrightarrow\pi_{n+3}(S^{n+1})
 \qquad(n\geq4),
\]
because \(n+2\leq2n-2\).  Therefore stabilization from
\(\pi_6(S^4)\) is an isomorphism onto \(\pi_2^{\rm S}\).
The composite \(\eta_4\circ\eta_5:S^6\to S^4\) stabilizes to
\(\eta^2\), so injectivity gives the exact equality
\cref{maps:src:B.3.1}.  Finally, a based null-homotopy of \(q_2\)
would postcompose with \(p_3\) to a based null-homotopy of \(H\),
contradicting \cref{maps:src:B.3.1}.
\end{enumerate}
\end{proof}

\paragraph{Exact scope of the comparison with the angle-forgetting map.}
The map here retains the angular coordinate before collapsing the
coordinate two-skeleton.  Its regular inverse image is the specific
normally four-framed \(T^2\) above.  The angle-forgetting
postcomposition has a normally three-framed \(T^3\) inverse image.
The calculation above is independent of its class and supplies an
actual obstruction to null-homotopy of the original \(q_2\), even
when that other postcomposition is null-homotopic.  The exact
morphism carrying this obstruction is
\[
 (p_3)_*:[X,\mathcal A]_*\longrightarrow[X,S^4]_*,\qquad
 [q_2]\longmapsto[H]=\eta_4\eta_5.
\]
No conclusion about an arbitrary earlier quotient follows merely
from the diffeomorphism type of one inverse image; the specified
four-framing is what proves this value.

\begin{corollary}[The exact order-two quotient detected by the original cusp map]
\label{wb:h24:ret:q2-detected-quotient}
\StatementRecord{wb:h24:ret:q2-detected-quotient}{proved}{Historical events: GPT-6 Astra (OpenAI), ChatGPT Pro mode, effort Pro, S6\_09\_PAPERS, 17 November 2026: targeted proof check. Also checked in this writing. Pre-continuation writing events: GPT-6 Astra (OpenAI), ChatGPT Pro mode, effort Pro, S6\_10\_PAPERS, 18 November 2026: targeted proof check. Continuation metadata events: GPT-6 Astra (OpenAI), ChatGPT Pro mode, effort Pro, S6\_10\_PAPERS continuation, 19 November 2026: used and rechecked in this run; status-vocabulary reconciliation; metadata reconciliation only; no new mathematical check. \emph{Pre-continuation recorded proof scope.} /root, effort Pro, 2026-11-17; also checked in this writing by GPT-6 Astra (OpenAI), ChatGPT Pro mode, effort Pro, 18 November 2026}{the hypotheses, notation and definitions stated here; \cref{maps:src:B.3.1}}

Let \(d:(S^6,*)\to(X,x_{\mathrm{base}})\) be the fixed based
diffeomorphism, and put
\[
\alpha=[q_2\circ d]\in\pi_6(\mathcal A,v_0).
\]
The original top-cell postcomposition induces a surjective homomorphism
\[
\varepsilon=(p_3)_*:\pi_6(\mathcal A,v_0)
  \longrightarrow\pi_6(S^4,*)\cong\mathbb Z/2,
\qquad \varepsilon(\alpha)=\eta_4\circ\eta_5.
\]
It gives the exact quotient isomorphism
\[
\pi_6(\mathcal A,v_0)/\ker\varepsilon
 \xrightarrow{\ [\beta]\mapsto\varepsilon(\beta)\ }
\mathbb Z/2.
\]
Every odd multiple of \(\alpha\) is nonzero. Its order is either
infinite or an even positive integer; this calculation does not determine
that order or assert a splitting of the displayed quotient.
\end{corollary}


\begin{proof}
\begin{enumerate}[label=\arabic*.,leftmargin=*]
\item Postcomposition distributes over the pinch sum on the source sphere,
so \((p_3)_*\) is a group homomorphism. Theorem~\ref{maps:src:B.3.1}
gives \((p_3)_*(\alpha)=[p_3q_2d]=\eta_4\eta_5\ne0\).
This is the nonzero element of the cyclic group of order two, proving
surjectivity. The first isomorphism theorem gives the stated quotient
with exactly the displayed map. For every integer \(n\),
\[
\varepsilon(n\alpha)=(n\bmod2)\eta_4\eta_5.
\]
An odd multiple therefore has nonzero image and cannot vanish.
If \(m\alpha=0\), applying \(\varepsilon\) forces \(m\)
to be even. These arguments prove all the assertions without inferring
the exact order of \(\alpha\).
\end{enumerate}
\end{proof}





\begin{lemma}[The three-dimensional coned Seifert target]
\label{mapsIII:carrier-three-complex-wedge}
\StatementRecord{mapsIII:carrier-three-complex-wedge}{proved}{No earlier separately identified check is recorded. Checked only in this writing. Pre-continuation writing events: GPT-6 Astra (OpenAI), ChatGPT Pro mode, effort Pro, S6\_10\_PAPERS, 18 November 2026: targeted proof check. Continuation metadata events: GPT-6 Astra (OpenAI), ChatGPT Pro mode, effort Pro, S6\_10\_PAPERS continuation, 19 November 2026: used and rechecked in this run; status-vocabulary reconciliation; metadata reconciliation only; no new mathematical check. \emph{Pre-continuation recorded proof scope.} GPT-6 Astra (OpenAI, ChatGPT Pro mode), S6\_09\_PAPERS, effort Pro, 16-17 November 2026; also checked in this writing by GPT-6 Astra (OpenAI), ChatGPT Pro mode, effort Pro, 18 November 2026}{the hypotheses, notation and definitions stated here; the explicit constructions and calculations in the proof}

Let $C$ be a simply connected CW complex with
$\widetilde H_2(C;\mathbb Z)=\widetilde H_3(C;\mathbb Z)=\mathbb Z$
and no other reduced integral homology. Then
$C\simeq S^2\vee S^3$.
\end{lemma}



\begin{proof}
\begin{enumerate}[label=\arabic*.,leftmargin=*]
\item The Hurewicz isomorphism $\pi_2(C)\to H_2(C;\mathbb Z)$ supplies a
based map $a:S^2\to C$ inducing an isomorphism on $H_2$ and on
$\pi_2$. Replace it by its mapping-cylinder inclusion
$S^2\hookrightarrow M_a$, where $M_a\simeq C$.
The pair $(M_a,S^2)$ is $2$-connected: both spaces are simply
connected, and the map on $\pi_2$ is an isomorphism.
Relative Hurewicz therefore gives
\[
 \pi_3(M_a,S^2)\xrightarrow{\ \cong\ }
 H_3(M_a,S^2;\mathbb Z).
\]
The homology exact sequence and $a_*:H_2(S^2)\cong H_2(M_a)$
identify this relative group with $H_3(M_a;\mathbb Z)$.
The homotopy exact sequence makes
$\pi_3(M_a)\to\pi_3(M_a,S^2)$ surjective, because its next map
lands in the kernel of the isomorphism
$\pi_2(S^2)\to\pi_2(M_a)$.
Naturality of the Hurewicz homomorphism then supplies a map
$b:S^3\to C$ representing an integral generator of $H_3(C)$.
The wedge $a\vee b:S^2\vee S^3\to C$ is an integral homology
equivalence between simply connected CW complexes. The homological
Whitehead theorem, proved in the classical topology appendix of Paper I,
makes it a homotopy equivalence.
\end{enumerate}
\end{proof}

For the coned Seifert target $\mathcal C$, the previously displayed
presentation kills the remaining boundary generators and gives
$\pi_1(\mathcal C)=1$, while Poincaré--Lefschetz duality gives exactly
the homology in this lemma. This provides the needed wedge equivalence
without an additional appeal to Whitehead's exact sequence.



\begin{proposition}[The actual degree-three finite-band map]
\label{mapsaudit:finite-degree-three}
\StatementRecord{mapsaudit:finite-degree-three}{proved}{No earlier separately identified check is recorded. Checked only in this writing. Pre-continuation writing events: GPT-6 Astra (OpenAI), ChatGPT Pro mode, effort Pro, S6\_10\_PAPERS, 18 November 2026: targeted proof check. Continuation metadata events: GPT-6 Astra (OpenAI), ChatGPT Pro mode, effort Pro, S6\_10\_PAPERS continuation, 19 November 2026: used and rechecked in this run; status-vocabulary reconciliation; metadata reconciliation only; no new mathematical check. \emph{Pre-continuation recorded proof scope.} GPT-6 Astra (OpenAI, ChatGPT Pro mode), S6\_09\_PAPERS, effort Pro, 16-17 November 2026; also checked in this writing by GPT-6 Astra (OpenAI), ChatGPT Pro mode, effort Pro, 18 November 2026}{the hypotheses, notation and definitions stated here; the explicit constructions and calculations in the proof}

Let $V=\mathbb Z\langle\gamma,u,w,\delta\rangle=H^1(F;\mathbb Z)$,
with $\operatorname{vol}=\gamma u w\delta$, and write
\[
 b=\gamma uw,\qquad
 c_1=uw\delta-4\gamma u\delta-2\gamma w\delta,\qquad
 c_2=uw\delta-3\gamma u\delta-3\gamma w\delta.
\]
For the degree-three and degree-four finite coverings $\pi_j:F\to S_j$,
\[
 \pi_1^*H^3(S_1;\mathbb Z)=\mathbb Zb\oplus\mathbb Zc_1,
 \qquad
 \pi_2^*H^3(S_2;\mathbb Z)=\mathbb Z(2b)\oplus\mathbb Zc_2.
\]
Choose the unique cohomology bases having those ordered pullbacks and the
Kronecker-dual bases of $H_3(S_j;\mathbb Z)$. Relative to the ordered
homology basis dual to $(\gamma uw,\gamma u\delta,\gamma w\delta,uw\delta)$,
the finite-band Mayer--Vietoris map is
\[
 \alpha_3=(\pi_{1*},-\pi_{2*})=
 \begin{pmatrix}
 1&0&0&0\\
 0&-4&-2&1\\
 -2&0&0&0\\
 0&3&3&-1
 \end{pmatrix}.
\]
Its Smith form has nonzero entries $(1,1,1)$ and its kernel is generated
by $(0,1,1,6)^t$. The order-three pushforward on $H_3$ is surjective;
the order-four pushforward has index two. In particular neither is
identified with its degree-one pushforward by integral changes of bases.
\end{proposition}


\begin{proof}
\begin{enumerate}[label=\arabic*.,leftmargin=*]
\item The integral fixed subgroups of $\bigwedge^3 V$ are
$\langle b,c_1\rangle$ and $\langle b,c_2\rangle$, respectively, by
solving the two integral invariance equations. The degree-one pullback
lattices are $\langle3\gamma,\eta_1\rangle$ and
$\langle4\gamma,\eta_2\rangle$, where
$\eta_1=2u+w+3\delta$ and $\eta_2=u+w+2\delta$.
The pairings of these bases with $(b,c_j)$ have matrices
\[
 \begin{pmatrix}0&3\\-3&0\end{pmatrix},\qquad
 \begin{pmatrix}0&4\\-2&0\end{pmatrix}.
\]
The downstairs $H^1$--$H^3$ pairing is unimodular, because the finite
surfaces have torsion-free integral cohomology and are oriented closed
four-manifolds. Pullback multiplies every pairing by the covering degree.
Thus its degree-three image has index one in the first invariant lattice
and index two in the second. If $rb+sc_2$ descends, pairing with the
descended class $\eta_2$ is divisible by four, so $-2r$ is divisible by
four. Hence its coefficient $r$ is even. This parity sublattice already
has index two, so it is the whole image. This proves the two displayed
pullback formulas, including sufficiency.

\item For $z\in H_3(F)$ and $a\in H^3(S_j)$, naturality of the actual
Kronecker pairing gives $\langle a,\pi_{j*}z\rangle
=\langle\pi_j^*a,z\rangle$. Evaluating the four listed classes therefore
gives the stated matrix, with the minus sign on the second summand.
Its $3\times3$ minors include $1$, and its rank is three. Solving its
kernel equations gives $z_1=0$, $z_3=z_2$, $z_4=6z_2$, proving the
remaining assertions. The first two rows have image all of $\mathbb Z^2$;
the last two have image $2\mathbb Z\oplus\mathbb Z$. In degree one the
corresponding image indices are three and four. These indices are
unchanged by unimodular changes of basis and rule out the asserted
degree-one/degree-three identification in the working proof.
\end{enumerate}
\end{proof}

\subsection{An extension homotopy for the angle-forgotten carriers}
\label{mapsIII:carrier-extension-section}

The literal flat quotient on the punctured cusp need not extend
continuously across the central fibre. Its homotopy class as a map to a
torus does extend. This distinction gives a second proof of the
angle-forgotten rank-two nullhomotopy and settles the one-copy circle
question left open in the working-edition homotopy-status remark.

\begin{lemma}[Suspending a torus-valued map that extends up to homotopy]
\label{mapsIII:carrier-extension-lemma}
\StatementRecord{mapsIII:carrier-extension-lemma}{proved}{Historical events: GPT-6 Astra (OpenAI), ChatGPT Pro mode, effort Pro, S6\_09\_PAPERS, 17 November 2026: targeted proof check. Also checked in this writing. Pre-continuation writing events: GPT-6 Astra (OpenAI), ChatGPT Pro mode, effort Pro, S6\_10\_PAPERS, 18 November 2026: targeted proof check. Continuation metadata events: GPT-6 Astra (OpenAI), ChatGPT Pro mode, effort Pro, S6\_10\_PAPERS continuation, 19 November 2026: used and rechecked in this run; status-vocabulary reconciliation; metadata reconciliation only; no new mathematical check. \emph{Pre-continuation recorded proof scope.} GPT-6 Astra (OpenAI, ChatGPT Pro mode), effort Pro, 2026-11-17; also checked in this writing by GPT-6 Astra (OpenAI), ChatGPT Pro mode, effort Pro, 18 November 2026}{the hypotheses, notation and definitions stated here; the explicit constructions and calculations in the proof}

Let $X$ be a compact Hausdorff space, and let $N\subset X$ be a compact
manifold with boundary which is a closed neighbourhood and satisfies
$\operatorname{int}_X(N)=N\setminus\partial N$.
Thus its frontier in $X$ is exactly its manifold boundary.
Let $r:N\to[0,\epsilon]$ be continuous, with
$W=r^{-1}(0)$, $\partial N=r^{-1}(\epsilon)$, and put
$U=N\setminus W$. Suppose $U$ is connected and locally path connected.
Let $\alpha:U\to T^d=\mathbb R^d/\mathbb Z^d$ be continuous.
Suppose there is a continuous $h:N\to T^d$ for which
$\alpha$ and $h|_U$ induce the same homomorphism to
$\pi_1(T^d)=\mathbb Z^d$. Fix $x_0\in W$, and choose $h(x_0)=0$.
For a based continuous map $c:(T^d,0)\to(Z,z_0)$ to a compact
Hausdorff based space, define
\[
 G(x)=
 \begin{cases}
 [r(x),c(\alpha(x))],&x\in U,\\
 *,&x\in W\cup(X\setminus\operatorname{int}N),
 \end{cases}
 \qquad
 \Sigma_r Z=
 \frac{[0,\epsilon]\times Z}
 {(\{0,\epsilon\}\times Z)\cup([0,\epsilon]\times\{z_0\})}.
\]
Then $G:X\to\Sigma_rZ$ is based nullhomotopic, by a homotopy
stationary on $(X\setminus\operatorname{int}N)\cup\{x_0\}$.
\end{lemma}



\begin{proof}
\begin{enumerate}[label=\arabic*.,leftmargin=*]
\item The difference $\alpha-h|_U:U\to T^d$ induces the zero homomorphism
on fundamental groups. The lifting criterion for the covering
$\mathbb R^d\to T^d$ therefore supplies a continuous function
$L:U\to\mathbb R^d$ with
\[
 \alpha(x)=h(x)+[L(x)].
\]
For $s\in[0,1]$, put
\[
 A_s(x)=h(x)+[(1-s)L(x)],\qquad
 G_s(x)=[r(x),c(A_s(x))]\quad(x\in U),
\]
and set $G_s=*$ on $W$ and outside $\operatorname{int}N$.
These formulas agree on the outer boundary. They are jointly continuous
at $W$ and at $\partial N$, even if $L$ is unbounded toward $W$:
every neighbourhood of the suspension basepoint contains the images of
$[0,\delta)\times Z$ and $(\epsilon-\delta,\epsilon]\times Z$
for some $\delta>0$, by compactness of $Z$. This estimate is uniform
in $s$ and in $A_s(x)$. Thus $G_s$ is a based homotopy from $G$
to the map $x\mapsto[r(x),c(h(x))]$ on $N$.

\item Now use that $h$ is defined on all of $N$. Set
\[
 H_s(x)=[(1-s)r(x)+s\epsilon,c(h(x))]
 \quad(x\in N),
 \qquad H_s(x)=*\quad(x\notin\operatorname{int}N).
\]
This is continuous by the closed pasting lemma. The height on
$\partial N$ is always $\epsilon$, so the formulas agree there.
At $x_0$ the second coordinate is $c(0)=z_0$, so $H_s(x_0)=*$
for every $s$. Finally $H_0=G_1$ and $H_1=*$.
Concatenating the two homotopies proves the assertion.
\end{enumerate}
\end{proof}

\begin{proposition}[Both angle-forgotten one-copy carriers are nullhomotopic]
\label{mapsIII:angle-forgotten-one-copy-null}
\StatementRecord{mapsIII:angle-forgotten-one-copy-null}{proved}{Historical events: GPT-6 Astra (OpenAI), ChatGPT Pro mode, effort Pro, S6\_09\_PAPERS, 17 November 2026: targeted proof check. Also checked in this writing. Pre-continuation writing events: GPT-6 Astra (OpenAI), ChatGPT Pro mode, effort Pro, S6\_10\_PAPERS, 18 November 2026: targeted proof check. Continuation metadata events: GPT-6 Astra (OpenAI), ChatGPT Pro mode, effort Pro, S6\_10\_PAPERS continuation, 19 November 2026: used and rechecked in this run; status-vocabulary reconciliation; metadata reconciliation only; no new mathematical check. \emph{Pre-continuation recorded proof scope.} GPT-6 Astra (OpenAI, ChatGPT Pro mode), effort Pro, 2026-11-17; also checked in this writing by GPT-6 Astra (OpenAI), ChatGPT Pro mode, effort Pro, 18 November 2026}{the hypotheses, notation and definitions stated here; \cref{mapsIII:carrier-extension-lemma}}

For the actual cusp neighbourhood $N_0=f^{-1}(|t|\leq\epsilon)$ and
the original flat marking, the maps
\[
 \widehat q_3:X\longrightarrow S^3,
 \qquad
 \widehat q_\gamma:X\longrightarrow S^2
\]
defined by forgetting the base angle and taking the reduced suspension
are based nullhomotopic. The assertion concerns their original formulas
and basepoint $x_0\in W$. It also gives based nullhomotopies of every
positive smash power of either map.
\end{proposition}



\begin{proof}
\begin{enumerate}[label=\arabic*.,leftmargin=*]
\item Write
\[
 \Lambda=\mathbb Z\langle\widehat\gamma,\widehat u,
                    \widehat w,\widehat\delta\rangle,
 \quad
 K=\mathbb Z\langle\widehat w,\widehat\delta\rangle,
 \quad
 \bar\Lambda=\Lambda/K.
\]
The geometric fundamental-group calculation in
\cref{mapsI:top:prop:geometric-pi1} proves the following
local information before taking any global quotient:
\[
 \pi_1(N_0)=\bar\Lambda,\qquad
 \pi_1(N_0\setminus W)=\Lambda\rtimes_{M_0}\mathbb Z,
 \qquad
 (\lambda,n)\longmapsto\bar\lambda.
\]
Take the radius $\epsilon$ strictly inside the punctured-overlap chart.
Because the fibration is a submersion over a small annulus about
$|t|=\epsilon$, a radial product collar identifies the fundamental
groups of this closed neighbourhood with those of the slightly larger
open neighbourhood used in that calculation, and identifies the
punctured inclusions as well.  This same collar makes $N_0$ a closed
codimension-zero neighbourhood in $X$, with
$\operatorname{int}_X(N_0)=N_0\setminus\partial N_0$ and frontier
$\partial N_0=f^{-1}(|t|=\epsilon)$, as required by the lemma.
Here the generator of the second factor is the zero-section meridian.
It bounds the actual holomorphic disc
$t\mapsto(t,1,1)$ in the toric chart, so there is no additional
translation in the displayed inclusion. Also
$(M_0-I)\Lambda=K$, as follows directly from
$(M_0-I)(a,b,c,d)=(0,0,b,-a)$.

\item The native map
$\alpha:N_0\setminus W\to\bar T^2=
(\bar\Lambda\otimes\mathbb R)/\bar\Lambda$
therefore induces precisely $(\lambda,n)\mapsto\bar\lambda$.
There is a map $h:N_0\to\bar T^2$ inducing the identity on
$\pi_1(N_0)=\bar\Lambda$. To see existence without any extra
classification assumption, take a CW model for the compact smooth
manifold $N_0$, map its one-cells with the specified integral winding
numbers, and extend over each two-cell because its attaching word
has trivial image in $\bar\Lambda$. All higher extensions exist
because the universal cover $\mathbb R^2$ of $\bar T^2$ is
contractible, so $\pi_k(\bar T^2)=0$ for $k\geq2$.
Compose with the CW comparison map to obtain $h$ on $N_0$.
Translation in the target arranges $h(x_0)=0$ and changes none of
the induced fundamental-group homomorphisms.

\item For $\widehat q_3$, apply Lemma~\ref{mapsIII:carrier-extension-lemma}
with $r=|t|$ and with
\[
 c:\bar T^2\longrightarrow
 \bar T^2/\bigl((S^1\times\{0\})\cup(\{0\}\times S^1)\bigr)
 \cong S^2.
\]
Its reduced suspension is exactly the original
$\widehat q_3(x)=[|t|,c(\alpha(x))]$ on the punctured cusp.
For $\widehat q_\gamma$, compose $\alpha$ and $h$ with the first
coordinate character $\bar T^2\to S^1$ and take $c$ to be the
identity of the based circle. The same lemma gives the nullhomotopy.
Both homotopies fix $x_0$, so their smash products with any based maps
are well defined and nullhomotopic.
\end{enumerate}
\end{proof}

\paragraph{Scope and preservation of the other proof.}
This argument uses the proved local fundamental-group inclusion and does
not produce a pointwise extension of the original flat character across
$W$. In particular, it does not alter the nonextension result for that
literal character. The explicit three-framed torus and its embedded
nullbordism provide a distinct proof of $\widehat q_3\simeq*$ and
are retained. The retained-angle map $q_2$ and its nonzero
$S^4$-valued postcomposition have different targets; the present
homotopy does not apply to them because their base-angle coordinate has
nonzero winding on the meridian that is killed in $N_0$.




